% ============================================================
%  DERECHO CIVIL --- PARTE GENERAL
%  Apuntes integrados de estudio
%  Compilar con: pdflatex (dos o tres pasadas)
% ============================================================
\documentclass[12pt,oneside]{book}

\usepackage[utf8]{inputenc}
\usepackage[T1]{fontenc}
\usepackage[spanish,es-nodecimaldot]{babel}
\usepackage[
top=2.5cm,
bottom=2.5cm,
left=3cm,
right=2.5cm,
headheight=15pt
]{geometry}
\usepackage{amsmath}
\usepackage{amssymb}
\usepackage{mathtools}
\usepackage{graphicx}
\usepackage{float}
\usepackage{caption}
\usepackage{subcaption}
\usepackage{booktabs}
\usepackage{array}
\usepackage{longtable}
\usepackage{tabularx}
\usepackage{enumitem}
\usepackage{xcolor}
\usepackage[most]{tcolorbox}
\usepackage{fancyhdr}
\usepackage{ragged2e}
\usepackage{titlesec}
\usepackage[
colorlinks=true,
linkcolor=blue!50!black,
citecolor=green!40!black,
urlcolor=blue!60!black,
bookmarks=true
]{hyperref}

\hypersetup{
pdftitle={Derecho Civil --- Parte General: apuntes integrados},
pdfauthor={Isaac Edgar Camacho Ocampo},
pdfsubject={Derecho Civil --- Parte General}
}

% ── Colores ───────────────────────────────────────────────────
\definecolor{acento}{RGB}{30,60,110}
\definecolor{fondodef}{RGB}{232,240,252}
\definecolor{fondoimp}{RGB}{255,246,226}
\definecolor{fondoeje}{RGB}{238,248,238}
\definecolor{fondoreg}{RGB}{244,244,244}
\definecolor{fondonota}{RGB}{246,246,250}
\definecolor{fondoadv}{RGB}{253,236,234}
\definecolor{fondocon}{RGB}{236,243,236}

% ── Control tipográfico ──────────────────────────────────────
\setlength{\parindent}{0pt}
\setlength{\parskip}{0.5em}
\clubpenalty=10000
\widowpenalty=10000
\emergencystretch=1.5em
\setlist{nosep,leftmargin=1.6em,itemsep=2pt,topsep=3pt}
\setcounter{secnumdepth}{3}
\setcounter{tocdepth}{2}
\captionsetup{font=small,labelfont=bf}

% ── Títulos ──────────────────────────────────────────────────
\titleformat{\chapter}[hang]{\normalfont\LARGE\bfseries\color{acento}\raggedright}{\thechapter.}{0.6em}{}
\titlespacing*{\chapter}{0pt}{-10pt}{22pt}
\titleformat{\section}{\normalfont\Large\bfseries\color{acento}}{\thesection}{0.6em}{}
\titlespacing*{\section}{0pt}{16pt}{6pt}
\titleformat{\subsection}{\normalfont\large\bfseries}{\thesubsection}{0.6em}{}
\titlespacing*{\subsection}{0pt}{12pt}{4pt}
\titleformat{\subsubsection}{\normalfont\normalsize\bfseries\itshape}{\thesubsubsection}{0.6em}{}
\titlespacing*{\subsubsection}{0pt}{10pt}{3pt}

% ── Encabezados y pies ───────────────────────────────────────
\pagestyle{fancy}
\fancyhf{}
\renewcommand{\chaptermark}[1]{\markboth{#1}{}}
\fancyhead[L]{\small\itshape\leftmark}
\fancyhead[R]{\small\itshape Derecho Civil --- Parte General}
\fancyfoot[C]{\thepage}
\renewcommand{\headrulewidth}{0.4pt}
\fancypagestyle{plain}{%
\fancyhf{}%
\fancyfoot[C]{\thepage}%
\renewcommand{\headrulewidth}{0pt}}

% ── Cajas académicas ─────────────────────────────────────────
\tcbset{aca/.style={enhanced,breakable,boxrule=0.4pt,arc=1.5pt,
left=6pt,right=6pt,top=4pt,bottom=4pt,fonttitle=\bfseries\small,
before skip=8pt,after skip=8pt}}
\newtcolorbox{definicion}[1][]{aca,colback=fondodef,colframe=acento,
title={Definición\if\relax\detokenize{#1}\relax\else: #1\fi}}
\newtcolorbox{concepto}[1][]{aca,colback=fondocon,colframe=green!40!black,
title={Concepto clave\if\relax\detokenize{#1}\relax\else: #1\fi}}
\newtcolorbox{importante}[1][]{aca,colback=fondoimp,colframe=orange!70!black,
title={Importante\if\relax\detokenize{#1}\relax\else: #1\fi}}
\newtcolorbox{ejemplo}[1][]{aca,colback=fondoeje,colframe=green!50!black,
title={Ejemplo\if\relax\detokenize{#1}\relax\else: #1\fi}}
\newtcolorbox{regla}[1][]{aca,colback=fondoreg,colframe=black!60,
title={Regla\if\relax\detokenize{#1}\relax\else: #1\fi}}
\newtcolorbox{ejercicio}[1][]{aca,colback=fondonota,colframe=violet!60!black,
title={Ejercicio\if\relax\detokenize{#1}\relax\else: #1\fi}}
\newtcolorbox{nota}[1][]{aca,colback=fondonota,colframe=black!45,
title={Nota\if\relax\detokenize{#1}\relax\else: #1\fi}}
\newtcolorbox{advertencia}[1][]{aca,colback=fondoadv,colframe=red!60!black,
title={Advertencia\if\relax\detokenize{#1}\relax\else: #1\fi}}

% ── Columnas de tabla ────────────────────────────────────────
\newcolumntype{Y}{>{\raggedright\arraybackslash}X}
\newcolumntype{L}[1]{>{\raggedright\arraybackslash}p{#1}}
\newcolumntype{B}[1]{>{\bfseries\raggedright\arraybackslash}p{#1}}

% ── Cuadro Cornell ───────────────────────────────────────────
\newcommand{\cq}[2]{\textbf{#1} & #2 \\ \hline}
\newcommand{\cornell}[3]{%
\par\bigskip\noindent
\begin{tabularx}{\textwidth}{|>{\raggedright\arraybackslash}p{0.30\textwidth}|Y|}
\hline
\multicolumn{2}{|l|}{\textbf{\small Cuadro Cornell --- #1}} \\ \hline
\textbf{\small Preguntas / palabras clave} & \textbf{\small Notas / respuestas} \\ \hline
#2
\multicolumn{2}{|p{\dimexpr\textwidth-2\tabcolsep-2\arrayrulewidth\relax}|}{\textbf{Resumen.} #3} \\ \hline
\end{tabularx}
\par\medskip}

% NOTA EDITORIAL: los archivos fuente 04 y 05 no fueron provistos; el documento integra únicamente el contenido de los archivos recibidos.
% ============================================================
\begin{document}
\frontmatter
\pagestyle{empty}

% ── PORTADA ──────────────────────────────────────────────────
\begin{titlepage}
\centering
\vspace*{1cm}
% TODO: el logotipo institucional de las fuentes (../img/descarga.jpeg) no se incluye con los archivos; se inserta solo si existe.
\IfFileExists{../img/descarga.jpeg}{\includegraphics[scale=.4]{../img/descarga.jpeg}\par\vspace{1cm}}{}
{\large\textsc{Universidad de Buenos Aires}}\par\vspace{0.3cm}
{\large\textsc{Facultad de Derecho}}\par\vspace{2cm}
{\color{acento}\rule{\linewidth}{1pt}}\par\vspace{0.6cm}
{\Huge\bfseries Derecho Civil}\par\vspace{0.3cm}
{\LARGE Parte General}\par\vspace{0.5cm}
{\large\itshape Derecho, codificación, ley y persona humana: nociones, atributos y capacidad}\par\vspace{0.6cm}
{\color{acento}\rule{\linewidth}{1pt}}\par\vspace{1.2cm}
{\large Apuntes integrados de estudio}\par
\vfill
{\large Isaac Edgar Camacho Ocampo}\par\vspace{0.3cm}
{\large Buenos Aires, Argentina}\par\vspace{0.3cm}
{\large 2026}\par\vspace{0.8cm}
\begin{flushleft}\small\itshape
Este es un modesto aporte a los estudiantes de «Derecho Civil» y no pretende sustituir el material oficial de la materia.
\end{flushleft}
\end{titlepage}

\pagestyle{plain}
\setcounter{tocdepth}{2}
\tableofcontents

\mainmatter
\pagestyle{fancy}

\chapter*{Introducción general}
\addcontentsline{toc}{chapter}{Introducción general}
\markboth{Introducción general}{}

La persona humana es el centro del derecho civil. Una máxima del derecho romano recogida en el Digesto lo expresa así: todo el derecho está constituido en razón del hombre. Este volumen recorre el camino que va desde la idea de derecho y de justicia hasta la regulación concreta de la persona humana, sus atributos y su capacidad en el Código Civil y Comercial de la Nación (CCyC).

La exposición avanza de lo más general a lo más concreto:

\begin{itemize}
\item La \textbf{Parte I} establece qué se entiende por derecho y por justicia, cuáles son los sentidos de la palabra «derecho» y cómo se ubica el derecho civil dentro del derecho privado.
\item La \textbf{Parte II} explica el fenómeno de la codificación y el recorrido que conduce desde el Código de Vélez Sarsfield hasta el Código Civil y Comercial de 2014.
\item La \textbf{Parte III} estudia el marco general que rige todo el código: la constitucionalización del derecho privado, el Título Preliminar, los principios, los efectos de la ley en el tiempo y el orden público.
\item La \textbf{Parte IV} se ocupa de la persona humana: su fundamento, su reconocimiento y el comienzo de su existencia.
\item La \textbf{Parte V} estudia los atributos de la persona: nombre, domicilio y estado civil.
\item La \textbf{Parte VI} trata la capacidad: el régimen general, el de las personas menores de edad y las restricciones a la capacidad por discapacidad.
\end{itemize}

Cada capítulo comienza con una breve presentación de su problema central y, cuando la densidad conceptual lo justifica, termina con un cuadro Cornell para el estudio activo. Las citas de artículos reproducen los textos tal como figuran en las fuentes de estas notas; cuando las propias fuentes presentan inconsistencias, se las señala en comentarios internos del archivo \LaTeX{} con la marca \texttt{TODO}.

% ============================================================
\part{El derecho, la justicia y el derecho civil}
% ============================================================

\chapter{La noción de derecho y de justicia}
\label{cap:justicia}

Todo el estudio del derecho civil descansa sobre una pregunta previa: qué es el derecho. Este capítulo parte del significado original de la palabra, define la justicia como virtud que ordena las relaciones entre personas, explica la opción metodológica adoptada por estas notas y presenta los dos principios que estructuran el derecho privado: la dignidad de la persona y el bien común.

\section{Origen etimológico de la palabra derecho}

La palabra \textbf{derecho} proviene del latín \textit{directum}, que significa «lo que está alineado, enderezado»: aquello contrario a lo torcido, lo que ha sido sujeto a una regla y dirigido conforme a ella.

\begin{definicion}
Históricamente, el derecho se asoció también a la voz latina \textit{ius} y a la voz griega \textit{tó díkaion}, ambas expresivas de la idea de lo \textbf{justo}, entendido como el objeto de la justicia. Por una transición conceptual, la palabra \textit{derecho} llegó a identificarse con la noción de \textit{cosa justa}.
\end{definicion}

\section{La justicia como dar a cada uno lo suyo}

\begin{definicion}[Derecho y justicia]
El \textbf{derecho} es lo justo. \textbf{Lo justo} es el objeto de la \textbf{justicia}, y la justicia consiste en \textbf{dar a cada uno lo suyo}: una relación real de igualdad que vincula personas, cosas y comunidades según un criterio de \textbf{bien común}.
\end{definicion}

Los elementos ---personas, cosas y comunidad--- permiten determinar qué es lo justo en cada situación. A lo largo de toda la vida jurídica y de las distintas ramas del derecho está siempre en juego esa determinación. Lo justo puede estar fijado de dos modos:

\begin{itemize}
\item por la \textbf{ley positiva}, cuando la norma establece directamente qué debe ocurrir;
\item por la \textbf{sentencia judicial}, cuando es el juez quien determina qué es lo justo en el caso concreto.
\end{itemize}

\begin{ejemplo}[Nulidad y restitución de frutos]
Tras una nulidad, la ley determina qué es lo justo: el poseedor de buena fe devuelve los frutos desde cierto momento, mientras que el poseedor de mala fe los devuelve desde que los recibió. La restitución queda regulada por la ley, que fija el criterio de justicia para cada caso.
\end{ejemplo}

\begin{ejemplo}[Accidente de tránsito]
Cuando se discute el alcance del deber de reparar los daños de un accidente de tránsito, es la sentencia judicial la que determina qué es lo justo: lo que se debe entre las personas, en torno a los daños, en el marco de la comunidad y con criterio de bien común.
\end{ejemplo}

\subsection{Tipos de justicia}

El derecho civil tiene como tema de fondo la justicia, entendida como la virtud que ordena la relación con los demás. Dentro de ella se distinguen tres tipos, según la dirección de la relación que regulan:

\begin{table}[H]
\centering
\begin{tabularx}{\textwidth}{B{0.22\textwidth} Y B{0.24\textwidth}}
\toprule
\textbf{Tipo de justicia} & \textbf{Descripción} & \textbf{Campo de aplicación} \\
\midrule
Justicia legal & Del individuo hacia el todo: obliga a los miembros de la comunidad hacia el bien común. & Derecho público / constitucional \\
Justicia distributiva & Del todo hacia la parte: distribución de cargas y beneficios desde la comunidad hacia los individuos. & Derecho administrativo / tributario \\
Justicia conmutativa & Entre las partes: busca la igualdad en los intercambios y la reciprocidad. & Derecho civil / contratos / daños \\
\bottomrule
\end{tabularx}
\caption{Tipos de justicia}
\end{table}

\begin{definicion}[Justicia conmutativa]
Es la justicia propia de las relaciones entre particulares, que busca la igualdad en los intercambios y la reciprocidad. Es el campo en el que se mueve el derecho civil, en lo que se ha llamado la reciprocidad de los cambios.
\end{definicion}

\section{Opción metodológica: el iusnaturalismo}

La concepción del derecho vinculada a la justicia implica una \textbf{opción metodológica} que se aparta de dos posiciones alternativas:

\begin{itemize}
\item El \textbf{positivismo jurídico}, que reduce el fenómeno jurídico al texto de la ley escrita o a una pura forma procedimental, sin contenido sustantivo.
\item El \textbf{sociologismo jurídico}, que identifica al derecho con la percepción de lo que la mayoría desea, sin asegurar su vinculación con la idea de justicia.
\end{itemize}

La opción adoptada reconoce el valor de la ley positiva, pero sostiene que esta debe tener un \textbf{contenido sustantivo vinculado a principios fundamentales} accesibles a la \textbf{recta razón} humana.

\begin{definicion}[Concepción clásica del derecho]
El \textbf{derecho} es un \textbf{orden social}, fundado en la \textbf{virtud de la justicia}, que ordena las relaciones de \textbf{alteridad} ---relaciones de convivencia entre personas--- y cuyo contenido es dar a cada uno lo suyo.
\end{definicion}

\section{Dignidad de la persona y bien común}
\label{sec:dignidad-bien-comun}

De la concepción del derecho vinculada a la justicia se desprenden dos \textbf{principios estructurantes} del derecho privado.

\begin{importante}
\begin{enumerate}
\item \textbf{Dignidad de la persona humana.} El sujeto del derecho es la persona humana; según el Digesto, todo el derecho está constituido en orden a ella. La \textbf{dignidad} expresa la excelencia y el valor inherente a la persona: un valor inviolable por el que la persona nunca debe ser tratada como medio, sino siempre como un \textbf{fin en sí misma}.
\item \textbf{Bien común.} El respeto a la dignidad se complementa con la noción de \textbf{bien común político}: el conjunto de condiciones que permiten que las personas, las familias y las asociaciones logren sus propios fines y participen del bien dado a toda la comunidad.
\end{enumerate}
\end{importante}

Ambos principios constituyen el telón de fondo del derecho y de la justicia: la persona necesita bienes para realizar sus fines, los obtiene en el marco de relaciones con los demás y, en esas relaciones, debe respetarse la dimensión de justicia.

En el centro del derecho civil está, entonces, la persona humana, con su dignidad resaltada en el artículo 51 del CCyC, y protegida desde la concepción hasta la muerte natural (art.~19 CCyC). El derecho civil busca además el bien común político: el conjunto de condiciones que permite que las personas, las familias y las sociedades alcancen su perfección. Ambas normas se estudian en detalle en el capítulo~\ref{cap:persona} y en el capítulo~\ref{cap:por-nacer}.

\cornell{La noción de derecho y de justicia}{
\cq{¿De dónde proviene la palabra «derecho»?}{Del latín \textit{directum}, «lo alineado, enderezado». Se asoció también a \textit{ius} y a \textit{tó díkaion}, voces vinculadas a lo justo.}
\cq{¿Qué es lo justo y cómo puede fijarse?}{Es el objeto de la justicia, que consiste en dar a cada uno lo suyo con criterio de bien común. Lo fija la ley positiva (restitución de frutos tras una nulidad) o la sentencia judicial (deber de reparar en un accidente de tránsito).}
\cq{¿Cuáles son los tres tipos de justicia y cuál es propia del derecho civil?}{Legal (individuo hacia el todo), distributiva (todo hacia la parte) y conmutativa (entre las partes). El derecho civil se mueve en la justicia conmutativa.}
\cq{¿Por qué la concepción adoptada es iusnaturalista?}{Porque exige que la ley positiva tenga contenido sustantivo vinculado a principios accesibles a la recta razón, y rechaza reducir el derecho a la norma escrita (positivismo) o a la opinión mayoritaria (sociologismo).}
\cq{¿Cuáles son los dos principios estructurantes del derecho privado?}{La dignidad de la persona humana (valor inviolable; la persona es un fin en sí misma) y el bien común político (condiciones para que personas, familias y asociaciones alcancen sus fines).}
}{El derecho es lo justo, y la justicia consiste en dar a cada uno lo suyo en relaciones de alteridad ordenadas al bien común. La ley y la sentencia determinan lo justo; el derecho civil opera en la justicia conmutativa. Dignidad de la persona y bien común estructuran el derecho privado.}

% ------------------------------------------------------------
\chapter{Derecho objetivo, subjetivo y normativo}
\label{cap:objetivo}

La palabra «derecho» no tiene un único significado. Este capítulo distingue sus sentidos, expone las principales teorías sobre el derecho subjetivo ---la noción que más discusión ha generado--- y presenta las críticas que desde posiciones opuestas se le han dirigido.

\section{Las tres dimensiones de la palabra derecho}

La palabra \textbf{derecho} es una \textbf{palabra análoga}: posee varios significados que guardan relación entre sí. El primer analogado es el concepto de \textit{lo justo}; a partir de él, por analogía, se habla de derecho en otros sentidos.

\begin{definicion}[Derecho objetivo]
La misma \textbf{cosa justa}, el objeto de la justicia: aquello que se debe dar a la otra persona en una relación de justicia.
\end{definicion}

\begin{definicion}[Derecho subjetivo]
La \textbf{prerrogativa o facultad} reconocida a la persona para exigir de los demás un determinado comportamiento. Mira al sujeto y a su potestad de reclamar.
\end{definicion}

\begin{definicion}[Derecho normativo]
También llamado derecho efectivo o vigente: el conjunto de \textbf{normas} dispuestas por el legislador en un determinado país, es decir, las disposiciones y reglas dadas para determinar qué es lo justo.
\end{definicion}

La cosa justa (derecho objetivo) a veces está determinada por el derecho normativo; otras veces la ley no la determina directamente y debe establecerla la sentencia judicial en el caso concreto. Además de estos tres sentidos principales, también se denomina \textit{derecho} a la \textbf{sentencia judicial}, en tanto precisa lo justo en el caso concreto, y a la \textbf{ciencia del derecho}, la disciplina que estudia todo lo relacionado con el derecho y la justicia.

\section{Teorías clásicas del derecho subjetivo}

La noción de derecho subjetivo ha dado lugar a numerosas teorías. Las principales posiciones clásicas son las siguientes:

\begin{table}[H]
\centering
\begin{tabularx}{\textwidth}{B{0.27\textwidth} Y}
\toprule
\textbf{Autor} & \textbf{Concepción del derecho subjetivo} \\
\midrule
Francisco Suárez (1548--1617) & Facultad moral sobre lo propio: facultad de exigir aquello que es propio del sujeto (escolástica española). \\
F.~C.~von Savigny (1779--1861) & Poder atribuido a la voluntad. Visión marcadamente individualista que enfatiza la voluntad como elemento esencial. \\
Rudolf von Jhering & Interés jurídicamente protegido con seguridad en el goce. El interés es decisivo, pero falta el elemento de la voluntad y la potestad de reclamar. \\
\bottomrule
\end{tabularx}
\caption{Teorías clásicas del derecho subjetivo}
\end{table}

\begin{importante}
La teoría de \textbf{Savigny} presenta el riesgo propio del individualismo: al enfatizar la voluntad sin establecer un límite claro, puede dar lugar a ejercicios abusivos o arbitrarios del derecho. La teoría de \textbf{Jhering} señala correctamente que el interés es decisivo en materia de derecho, pero el goce en abstracto no puede concretarse en ausencia de la voluntad de ejercerlo: falta el elemento de la voluntad y la potestad de reclamar.
\end{importante}

\section{Críticas colectivista y positivista al derecho subjetivo}

Dos corrientes niegan, desde perspectivas opuestas, la existencia o la relevancia del derecho subjetivo.

\begin{definicion}[Posición de León Duguit (escuela colectivista)]
El hombre no tiene derechos, sino únicamente \textbf{deberes} derivados de su pertenencia a la comunidad. Todo individuo tiene una función que cumplir y tareas que ejecutar. El fundamento de la regla de derecho no es un derecho subjetivo que el ser humano pueda reclamar para sí, sino el deber social de obrar y desenvolver su individualidad conforme a esa función, tanto para gobernantes como para gobernados.
\end{definicion}

\begin{definicion}[Posición de Hans Kelsen (escuela positivista)]
La noción de derecho subjetivo es \textbf{superflua}, porque lo único que importa es el \textbf{derecho normativo}, que establece las obligaciones jurídicas. El derecho subjetivo sería la contracara de la obligación: un concepto reflejo de utilidad didáctica, pero prescindible para la descripción científica. Lo que se denomina derecho subjetivo no es más que la conducta encuadrada en una norma que, dado cierto presupuesto, adjudica determinadas consecuencias jurídicas.
\end{definicion}

Frente a estas posiciones se reconoce la \textbf{importancia del derecho subjetivo} como expresión de la \textbf{dignidad de la persona humana}, especialmente en lo relativo a los derechos humanos y los derechos fundamentales. A partir de las teorías clásicas, el derecho subjetivo puede entenderse como una \textbf{prerrogativa reconocida a la persona por el ordenamiento jurídico} para realizar un comportamiento tendiente a la satisfacción de \textbf{intereses humanos}, debiendo ejercerse en conformidad con esa finalidad.

\section{La teoría de John Finnis}

\textbf{John Finnis} (profesor de Oxford) concibe el derecho subjetivo como \textbf{exigencia derivada del principio de razonabilidad práctica}. Según esta concepción, el derecho subjetivo incluye:

\begin{itemize}
\item la \textbf{exigencia positiva o negativa} impuesta sobre otro: exigir un comportamiento determinado o exigir que el otro no interfiera con el propio desarrollo o disfrute de un bien;
\item la \textbf{capacidad de reclamar} el cumplimiento de esa exigencia, o la posibilidad de quedar inmune frente a un intento de sometimiento.
\end{itemize}

\begin{definicion}[Razonabilidad práctica (Finnis)]
Principio que remite a una concepción de \textbf{ley natural}: principios fundamentales superiores al ordenamiento positivo, que el ser humano puede conocer a través de la recta razón y a la luz de los cuales es posible valorar el contenido del derecho vigente.
\end{definicion}

\cornell{Sentidos del derecho y derecho subjetivo}{
\cq{¿Qué significa que «derecho» es una palabra análoga?}{Que sus sentidos (objetivo, subjetivo, normativo) guardan relación con un primer analogado: lo justo.}
\cq{¿Cómo se distinguen los tres sentidos principales?}{Objetivo: la cosa justa. Subjetivo: la facultad de exigir un comportamiento. Normativo: el conjunto de normas del legislador.}
\cq{¿Cómo definieron el derecho subjetivo Suárez, Savigny y Jhering?}{Suárez: facultad moral sobre lo propio. Savigny: poder de la voluntad, con riesgo de ejercicio abusivo. Jhering: interés jurídicamente protegido, sin voluntad ni potestad de reclamar.}
\cq{¿Por qué Duguit y Kelsen niegan el derecho subjetivo?}{Duguit: solo hay deberes sociales según la función de cada individuo. Kelsen: es superfluo; solo importa el derecho normativo y la obligación jurídica.}
\cq{¿Cómo concibe Finnis el derecho subjetivo?}{Como exigencia derivada de la razonabilidad práctica, con capacidad de reclamar su cumplimiento, remitida a principios de ley natural.}
}{El derecho es una palabra análoga cuyo primer sentido es lo justo. El derecho subjetivo, expresión de la dignidad humana, ha sido explicado por la voluntad, el interés o la exigencia razonable, y negado por el colectivismo y el positivismo.}

% ------------------------------------------------------------
\chapter{Las grandes divisiones del derecho}
\label{cap:divisiones}

El derecho puede clasificarse según el fundamento de sus principios (natural y positivo) y según los sujetos que intervienen en la relación regulada (público y privado). Estas dos divisiones permiten situar luego al derecho civil.

\section{Derecho natural y derecho positivo}

\begin{definicion}[Derecho natural]
Está constituido por \textbf{principios fundamentales} que corresponden a la naturaleza humana, accesibles al ser humano mediante la \textbf{recta razón} y derivados de las inclinaciones fundamentales del ser hacia sus fines propios.
\end{definicion}

Se distinguen tres grandes \textbf{inclinaciones fundamentales} del ser humano, de las que se desprenden principios del derecho natural:

\begin{enumerate}
\item la inclinación a \textbf{conservar la propia vida};
\item la inclinación a \textbf{transmitir la vida} y todo lo vinculado a la procreación;
\item la inclinación a \textbf{buscar el bien según la naturaleza racional} y a organizar la vida conforme a esa búsqueda.
\end{enumerate}

Estos principios son con frecuencia genéricos y requieren \textbf{concreciones}. Allí interviene el derecho positivo.

\begin{definicion}[Derecho positivo]
Es el \textbf{dispuesto por el legislador} en orden al bien común. Comprende todas las normas escritas: la Constitución, los tratados, las leyes y los decretos.
\end{definicion}

\begin{table}[H]
\centering
\begin{tabularx}{\textwidth}{B{0.27\textwidth} Y}
\toprule
\textbf{Posición} & \textbf{Tesis fundamental} \\
\midrule
Positivismo jurídico & No existe un principio superior al derecho positivo. El derecho es lo que establece la autoridad competente; su justicia se debate en el plano político legislativo. \\
Iusnaturalismo & Existen principios superiores al derecho positivo, accesibles a la recta razón, a la luz de los cuales es posible valorar el contenido y la justicia de la ley vigente. \\
Corrientes trans-positivistas & Reconocen principios o valores superiores a la ley positiva, pero sin identificarlos con la ley natural clásica. Incluyen el neoconstitucionalismo y posiciones consensualistas o procedimentalistas. \\
\bottomrule
\end{tabularx}
\caption{Posiciones sobre la relación entre derecho natural y derecho positivo}
\end{table}

\section{Derecho público y derecho privado}

Dentro del derecho positivo, la división fundamental distingue el derecho público del derecho privado.

\begin{definicion}[Derecho público y derecho privado]
El \textbf{derecho público} regula las relaciones en que \textbf{interviene el Estado} como comunidad frente a los individuos. El \textbf{derecho privado} regula las relaciones \textbf{entre particulares}, sin intervención del Estado en posición de poder.
\end{definicion}

\begin{table}[H]
\centering
\begin{tabularx}{\textwidth}{B{0.38\textwidth} Y}
\toprule
\textbf{Rama} & \textbf{Ámbito} \\
\midrule
Derecho constitucional, administrativo y penal & Público \\
Derecho civil y comercial & Privado \\
Derecho del trabajo & Privado (con elementos de derecho público) \\
Derecho procesal & Intersección entre público y privado \\
\bottomrule
\end{tabularx}
\caption{Principales ramas del derecho según su ámbito}
\end{table}

% ------------------------------------------------------------
\chapter{El derecho civil, la relación jurídica y su Parte General}
\label{cap:civil}

Situado el derecho privado, este capítulo define el derecho civil como su rama madre, explica el concepto de relación jurídica con sus tres elementos y describe cómo el Código Civil y Comercial organiza su contenido, con especial atención a la Parte General.

\section{El derecho civil como rama madre del derecho privado}

\begin{definicion}[Derecho civil]
Es la \textbf{rama madre del derecho privado}. Se ocupa del \textbf{hombre como sujeto de derecho}, sin distinción de cualidades accidentales, y de las \textbf{relaciones jurídicas} ---tanto patrimoniales como extrapatrimoniales--- que lo constituyen como sujeto, regulando las instituciones básicas y constituyendo el punto de partida de las demás ramas del derecho privado.
\end{definicion}

La palabra \textit{civil} proviene del derecho romano: el \textit{ius civile} era el derecho de los ciudadanos (\textit{civis}) de la ciudad (\textit{civitas}) romana. Tras la caída del Imperio Romano, este derecho subsistió en las distintas civilizaciones como derecho de los particulares, organizando las relaciones más fundamentales de la persona: obligaciones, contratos, derechos reales, familia y sucesiones.

\begin{importante}
El rasgo distintivo del derecho civil frente a otras ramas del derecho privado es que \textbf{mira a la persona en cuanto tal}, desprovista de cualidades accidentales. El derecho comercial la mira en cuanto comerciante; el derecho laboral, en cuanto trabajador; el derecho societario, en cuanto integrante de una sociedad. El derecho civil estudia a la persona humana y sus atributos (existencia, nombre, domicilio, estado, capacidad, patrimonio, derechos personalísimos), a la persona jurídica, a los bienes y las cosas, y a los hechos y actos jurídicos.
\end{importante}

\section{La relación jurídica: sujeto, objeto y causa}

Un concepto central para el estudio del derecho civil es el de relación jurídica.

\begin{definicion}[Relación jurídica]
Es el \textbf{vínculo jurídico} que se establece entre uno o varios \textbf{sujetos}, en torno a un \textbf{objeto determinado} y fundamentado en el \textbf{derecho}. Da lugar a consecuencias jurídicas entre las partes.
\end{definicion}

\begin{ejemplo}[Contrato de mutuo]
Cuando una persona presta dinero a otra se genera una relación jurídica entre dos sujetos (prestamista y prestatario), en torno a un objeto (el dinero), con una causa que es el contrato de mutuo, el cual fija las pautas de la relación y obliga a la restitución.
\end{ejemplo}

\subsection{Sujetos (elemento personal)}

Los \textbf{sujetos} son las personas que entablan relaciones jurídicas. El CCyC distingue dos clases:

\begin{itemize}
\item \textbf{Personas humanas:} todo ser humano, por el solo hecho de serlo, es persona. El término \textit{persona} denota el valor y la excelencia en el ser de quien es capaz de entablar relaciones jurídicas; no es sinónimo de ciudadano, habitante o individuo.
\item \textbf{Personas jurídicas:} surgen de la asociación de personas humanas para fines útiles, dando lugar a formas de actuación jurídicamente relevantes. Pueden ser \textit{públicas} (el Estado nacional, provincial y municipal; los Estados extranjeros; los organismos públicos internacionales; la Iglesia Católica, entre otras) o \textit{privadas} (las sociedades, las asociaciones civiles, las fundaciones, las iglesias y entidades religiosas, las cooperativas, el consorcio de propiedad horizontal, entre otras).
\end{itemize}

\subsection{Objeto (elemento real)}

El \textbf{objeto} es la \textbf{cosa o bien} sobre el que recae o se proyecta la relación jurídica. Los \textbf{bienes} son todos los objetos materiales e inmateriales susceptibles de tener valor, incluyendo ideas, derechos intelectuales, invenciones artísticas y científicas, y derechos creditorios. Las \textbf{cosas} son los objetos materiales susceptibles de tener valor y se clasifican, entre otras categorías, en muebles e inmuebles.

\subsection{Causa (elemento normativo)}

La \textbf{causa} es aquello que da origen a la relación jurídica. Puede provenir de:

\begin{itemize}
\item \textbf{hechos de la naturaleza}, como el granizo que arruina una cosecha y genera consecuencias jurídicas;
\item \textbf{hechos ilícitos}, como un accidente de tránsito;
\item \textbf{actos jurídicos}, como un contrato o un testamento.
\end{itemize}

\section{La Parte General en el Código Civil y Comercial}
\label{sec:parte-general}

El contenido del derecho civil se organiza en el Código Civil y Comercial de la Nación a través de seis libros:

\begin{table}[H]
\centering
\begin{tabularx}{\textwidth}{B{0.17\textwidth} Y}
\toprule
\textbf{Libro} & \textbf{Contenido} \\
\midrule
Primero & \textbf{Parte General.} La persona humana y sus atributos (existencia, nombre, domicilio, estado civil, capacidad, patrimonio, derechos personalísimos y situaciones especiales de capacidad); la persona jurídica y su clasificación; los bienes y las cosas; los hechos y los actos jurídicos como causa de las relaciones jurídicas. \\
Segundo & Relaciones de familia. \\
Tercero & Obligaciones y contratos, incluyendo todo lo relativo a la responsabilidad civil. \\
Cuarto & Derechos reales: derechos de las personas sobre las cosas. \\
Quinto & Sucesiones. \\
Sexto & Disposiciones comunes a los derechos reales y personales: derecho internacional privado, prescripción y otras disposiciones comunes. \\
\bottomrule
\end{tabularx}
\caption{Los seis libros del Código Civil y Comercial}
\end{table}

\begin{importante}
La \textbf{Parte General} (Libro Primero) no constituye un sistema lógico deductivo cerrado. Su función es filosófica e introductoria: fijar los contenidos y márgenes fundamentales del derecho privado ---persona, atributos, bienes y causa de las relaciones jurídicas--- para comprender cómo surgen las relaciones jurídicas y qué las rige. Es la puerta de entrada conceptual a todas las demás ramas del derecho privado.
\end{importante}

\cornell{El derecho civil y su Parte General}{
\cq{¿Por qué el derecho civil es la «rama madre» del derecho privado?}{Porque proviene del \textit{ius civile} romano, organizó las relaciones más fundamentales de la persona y la mira en cuanto tal, a diferencia de otras ramas que la contemplan como comerciante, trabajador o socio.}
\cq{¿Cómo se distinguen derecho natural y derecho positivo?}{El natural son principios accesibles a la recta razón, derivados de las inclinaciones humanas fundamentales; el positivo es el dispuesto por el legislador para concretarlos en orden al bien común.}
\cq{¿Cómo se distinguen derecho público y privado?}{El público regula las relaciones en que interviene el Estado; el privado, las relaciones entre particulares. El derecho procesal se sitúa en la intersección.}
\cq{¿Qué elementos componen la relación jurídica?}{Sujetos (personas humanas y jurídicas), objeto (bienes y cosas) y causa (hechos de la naturaleza, hechos ilícitos o actos jurídicos).}
\cq{¿Cuál es la función de la Parte General?}{No es un sistema deductivo cerrado: es una introducción filosófica y conceptual que fija persona, atributos, bienes y causa como puerta de entrada al derecho privado.}
}{El derecho civil, heredero del \textit{ius civile}, estudia a la persona en cuanto tal y sus relaciones jurídicas. Toda relación jurídica se compone de sujeto, objeto y causa, y la Parte General del CCyC reúne esos elementos como introducción a todo el derecho privado.}

% ============================================================
\part{La codificación del derecho civil}
% ============================================================

\chapter{La codificación: concepto, caracteres y orígenes}
\label{cap:codificacion}

El derecho civil y la codificación son fenómenos inseparables: no puede comprenderse cabalmente el primero sin conocer el segundo. Este capítulo define qué es codificar, qué caracteres tiene un código, qué ventajas e inconvenientes presenta y cómo nació el movimiento codificador moderno.

\section{Concepto y caracteres de la codificación}

\begin{definicion}[Codificación]
Codificar es reunir orgánicamente en un cuerpo único todas las normas vigentes sobre una materia determinada del derecho. Un \textbf{código} es un cuerpo sistemático, orgánico y metódico de todas las normas y principios jurídicos relativos a una rama específica del ordenamiento jurídico.
\end{definicion}

Un código presenta tres caracteres fundamentales:

\begin{itemize}
\item \textbf{Unitario}: unifica en un solo cuerpo las normas de la materia.
\item \textbf{Exclusivo}: no existen normas de la materia por fuera del código; todo lo referido a ella queda comprendido en él.
\item \textbf{Sistemático}: posee un orden, una lógica y un método propios, lo que facilita la comprensión del fenómeno jurídico regulado.
\end{itemize}

\section{Ventajas e inconvenientes de la codificación}

La codificación da unidad y coherencia al sistema legislativo y brinda \textbf{seguridad jurídica}: ante una duda, el código es el lugar único donde se espera encontrar la respuesta, lo que facilita el conocimiento del ordenamiento. Presenta, no obstante, inconvenientes. Pretender que todo el ordenamiento quepa en un único instrumento responde a un cierto racionalismo, pero en la práctica la materia tiende a desbordar el código. Además, un código puede \textbf{cristalizar} un momento histórico; cuando después resulta necesario reformarlo, existe el riesgo de que las reformas alteren su carácter sistemático y generen incoherencias internas, afectando la unidad, la exclusividad y el orden que lo caracterizan.

\begin{table}[H]
\centering
\begin{tabularx}{\textwidth}{YY}
\toprule
\textbf{Ventajas} & \textbf{Inconvenientes} \\
\midrule
Unidad y coherencia del sistema legislativo & Pretensión de exhaustividad nunca totalmente alcanzada \\
Seguridad jurídica: un lugar donde encontrar todo & Cristalización: el derecho puede quedar desactualizado \\
Facilita el conocimiento del ordenamiento & Las reformas parciales generan incoherencias internas \\
\bottomrule
\end{tabularx}
\caption{Ventajas e inconvenientes de la codificación}
\end{table}

\section{Orígenes históricos de la codificación moderna}

El movimiento codificador se vincula con el pensamiento de la Ilustración y el racionalismo europeo del siglo XVIII. Comienza en Prusia, donde se elabora el \textit{Código General de los Estados Prusianos}, obligatorio desde 1794 como derecho territorial de los estados y no limitado a la materia civil.

El código más importante de este movimiento es el \textbf{Código Civil francés de 1804} (Código Napoleón), que instaura el modelo de las codificaciones nacionales y ejerce una influencia decisiva, a inicios del siglo XIX, sobre numerosos países, incluidos los de América Latina. Los grandes codificadores latinoamericanos son \textbf{Freitas} (Brasil), \textbf{Andrés Bello} (Chile) y \textbf{Vélez Sarsfield} (Argentina); todos reciben una fuerte influencia del Código Napoleón, aunque se apartan de él en distintos aspectos.

\begin{nota}
La codificación no implica una ruptura con las instituciones preexistentes: el derecho civil abreva en una tradición jurídica que se remonta al derecho romano. La codificación pretende dotar a esas instituciones de unidad, sistematicidad y exclusividad, en consonancia con el espíritu racionalista e iluminista que confiaba en que la sola razón permitía abarcar y compendiar todo el conocimiento de una rama del derecho.
\end{nota}

\cornell{La codificación}{
\cq{¿Qué es codificar y qué es un código?}{Reunir orgánicamente en un cuerpo único las normas vigentes de una materia. El código es un cuerpo sistemático, orgánico y metódico de normas y principios de una rama del derecho.}
\cq{¿Cuáles son los tres caracteres de un código?}{Unitario, exclusivo y sistemático.}
\cq{¿Qué ventajas e inconvenientes tiene?}{Ventajas: unidad, coherencia y seguridad jurídica. Inconvenientes: la materia desborda el código, cristaliza un momento histórico y las reformas parciales generan incoherencias.}
\cq{¿Dónde y cuándo nace la codificación moderna?}{En la Ilustración y el racionalismo del siglo XVIII; Prusia (obligatorio en 1794) y, sobre todo, el Código Civil francés de 1804.}
}{Codificar es reunir en un cuerpo unitario, exclusivo y sistemático las normas de una materia, con ventaja de seguridad jurídica y riesgo de cristalización. El movimiento nace con la Ilustración y se consolida con el Código Napoleón.}

% ------------------------------------------------------------
\chapter{El Código de Vélez Sarsfield y la reforma de la Ley 17.711}
\label{cap:velez}

La codificación argentina comienza con el Código Civil de Vélez Sarsfield y atraviesa más de un siglo de proyectos de reforma antes de su sustitución completa. Este capítulo describe el proceso de elaboración de ese código, sus fuentes y su método, y la primera gran reforma que lo modificó: la Ley 17.711.

\section{Contexto constitucional y elaboración del Código de Vélez}

La Constitución Nacional contemplaba que fuese el \textbf{Congreso de la Nación} quien dictara el Código Civil para todo el país. Al organizarse el Estado argentino, las provincias delegaron en el Congreso esa potestad, definición de gran importancia porque otorgaba unidad a una nación surgida de guerras fratricidas y fuertes disputas internas, que necesitaba consolidar su identidad y encontrar elementos de cohesión. Esa cohesión debía expresarse, entre otros ámbitos, en las instituciones básicas: la persona, la familia, los contratos, la propiedad y las sucesiones. Las provincias conservaron la potestad de dictar sus propios códigos procesales y organizar sus sistemas judiciales.

Por ley de 1863 se autorizó al Poder Ejecutivo a nombrar una comisión redactora. En 1864, el presidente Mitre encomendó la tarea a \textbf{Vélez Sarsfield}, jurista cordobés, quien se dedicó a ella durante cinco años y terminó el proyecto en 1869. El presidente Sarmiento lo envió al Congreso, que lo aprobó «libro cerrado» mediante la Ley 340.

\begin{importante}[Ley 340 (1869)]
Aprobó el Código Civil de Vélez Sarsfield. Entró en vigencia el 1.º de enero de 1871.
\end{importante}

\section{Fuentes del Código de Vélez}

El Código de Vélez se nutrió de las siguientes fuentes:

\begin{enumerate}[label=\alph*)]
\item \textbf{Derecho romano}: llega a través de los estudios universitarios del propio Vélez Sarsfield y de la influencia de Savigny en su pensamiento.
\item \textbf{Legislación española y patria}: Vélez no pretendía marcar una ruptura con las formas y leyes del pueblo argentino, por lo que se advierte continuidad de instituciones provenientes de la legislación española.
\item \textbf{Derecho canónico}: influye especialmente en materia de matrimonio, cuyos efectos civiles quedan a cargo de Vélez Sarsfield (situación que cambiaría con la ley de matrimonio civil).
\item \textbf{Código Napoleón}: influye de forma directa e indirecta, a través de los comentaristas conocidos por Vélez, señalados en las notas al pie de su código.
\item \textbf{Código chileno de Andrés Bello, Código de Luisiana y Código uruguayo}.
\item \textbf{Proyecto de García Goyena} (España), obra de distintos juristas.
\end{enumerate}

\begin{importante}
La influencia más importante fue la de \textbf{Augusto Teixeira de Freitas}, jurista brasileño, autor de una obra monumental ---la \textit{Consolidación de las Leyes} y el \textit{Esbozo del Código Civil}--- de la que Vélez tomó instituciones y contenidos. Es considerada la mayor influencia recibida por el Código Civil argentino, especialmente en cuanto a su método.
\end{importante}

\section{Método del Código de Vélez}

El Código sancionado por la Ley 340 se estructuraba en dos títulos preliminares, cuatro libros y un título complementario:

\begin{table}[H]
\centering
\begin{tabularx}{\textwidth}{B{0.27\textwidth} Y}
\toprule
\textbf{Parte} & \textbf{Contenido} \\
\midrule
Título Preliminar I & De las leyes: teoría general de la ley \\
Título Preliminar II & Del modo de contar los intervalos del derecho \\
Libro I & De las personas (personas en general y derechos procedentes de las relaciones de familia) \\
Libro II & De los derechos personales en las relaciones civiles (obligaciones, hechos y actos jurídicos, contratos) \\
Libro III & De los derechos reales (vínculos directos entre persona y cosas) \\
Libro IV & Disposiciones comunes a los derechos reales y personales (sucesiones, privilegios, prescripción) \\
Título Complementario & Disposiciones de aplicación \\
\bottomrule
\end{tabularx}
\caption{Estructura del Código Civil de Vélez Sarsfield (Ley 340)}
\end{table}

El método fue siempre muy valorado por tratar en un único título los hechos y actos jurídicos y por unificar los derechos reales. El matrimonio no era tratado entre los contratos sino dentro del derecho de familia propiamente dicho, y las sucesiones estaban separadas de las donaciones. Se le criticaba, sin embargo, la ausencia de una \textbf{parte general}.

\section{Proyectos de reforma anteriores a 1968}

Hubo distintos proyectos de reforma al Código de Vélez antes de su primera gran transformación:

\begin{table}[H]
\centering
\begin{tabularx}{\textwidth}{B{0.10\textwidth} B{0.30\textwidth} Y}
\toprule
\textbf{Año} & \textbf{Proyecto} & \textbf{Características} \\
\midrule
1926 & Anteproyecto Bibiloni & Jurista a cargo de la redacción inicial \\
1936 & Proyecto del 36 & Comisión que tomó como base el trabajo de Bibiloni \\
1954 & Anteproyecto Llambías & Comisión presidida por Jorge J. Llambías; proyecto completo de reformas \\
\bottomrule
\end{tabularx}
\caption{Proyectos de reforma anteriores a la Ley 17.711}
\end{table}

\section{La Ley 17.711 (1968): primera reforma integral}

La reforma efectiva llegó cien años después del Código de Vélez.

\begin{importante}[Ley 17.711 (1968)]
Primera reforma integral del Código Civil argentino. Sancionada en abril de 1968, en vigencia desde el 1.º de julio de 1968.
\end{importante}

% TODO: verificar nombres de la comisión redactora presentes en las fuentes originales; la transcripción de los apellidos parece alterada (por ejemplo «Gorda», «Aprobar», «Tomar», «Híbrido»).
La comisión redactora estuvo inicialmente integrada, según las fuentes, por Fleitas, Semon, Gorda, José Media López Olaciregui, Piensa, Aprobar, Tomar, Tiene Ruiz y Alberto. Tras distintas renuncias, el proyecto fue firmado por Martínez Ruiz, Fleitas e Híbrido. \textbf{Borda}, ministro del interior, es considerado el principal autor de la iniciativa.

La ley constituyó una transformación estructural muy importante, aunque no una derogación completa: conservó la estructura general del código e introdujo reformas en aspectos sustanciales.

\subsection{Institutos incorporados}

\begin{definicion}[Lesión subjetiva (art.~954 del Código Civil)]
Vicio del acto jurídico que permite solicitar la nulidad o modificación del acto cuando una parte, explotando la necesidad, ligereza o inexperiencia de la otra, obtiene una ventaja patrimonial evidentemente desproporcionada. Vélez la había rechazado expresamente en la nota al artículo 943; la Ley 17.711 la incorporó recogiendo las propuestas del Tercer Congreso de Derecho Civil realizado en Córdoba.
\end{definicion}

\begin{definicion}[Teoría de la imprevisión]
Permite resolver un contrato cuando ocurren circunstancias sobrevinientes que alteran sustancialmente la igualdad entre las prestaciones, tornándolas desproporcionadas para una de las partes. Su finalidad es restablecer la justicia contractual.
\end{definicion}

\begin{definicion}[Abuso del derecho]
Reconoce el carácter funcional de los derechos subjetivos: estos no pueden ejercerse de manera abusiva o contraria a los fines para los cuales fueron reconocidos.
\end{definicion}

Además, la reforma incorporó el factor moral en el derecho positivo, la protección de la buena fe, la corrección de un individualismo excesivo presente en algunos institutos del Código de Vélez y la equidad ---también presente en la lesión subjetiva--- como herramienta para corregir situaciones en el caso concreto. Su reforma del régimen de los efectos de la ley en el tiempo se estudia en el capítulo~\ref{cap:tiempo}.

\subsection{Reformas posteriores y proyectos de unificación}

Con posterioridad se sucedieron reformas en materia de familia (leyes 23.264 y 23.515), en materia cambiaria y en materia de obligaciones, con la cuestión de la convertibilidad. Luego comenzaron los proyectos de reforma integral orientados a la \textbf{unificación} civil y comercial: en 1987 se conformó una comisión; en 1993 se presentaron el proyecto de la Cámara de Diputados ---elaborado por una comisión federal--- y un proyecto del Poder Ejecutivo; y en 1998 se presentó un nuevo proyecto, elaborado por una comisión designada por el Poder Ejecutivo. Este último ejerció una influencia decisiva sobre el Código Civil y Comercial de 2014, pues fue la base tomada por la comisión de 2011-2012.

\cornell{El Código de Vélez y la Ley 17.711}{
\cq{¿Cómo se elaboró y cuándo rigió el Código de Vélez?}{El Congreso encargó la comisión (ley de 1863); Vélez lo redactó entre 1864 y 1869; fue aprobado «libro cerrado» por la Ley 340 y rigió desde el 1.º de enero de 1871.}
\cq{¿Cuál fue su influencia principal?}{Teixeira de Freitas, en especial en cuanto al método, junto con el derecho romano, la legislación española, el derecho canónico y el Código Napoleón.}
\cq{¿Qué crítica se le hacía al método de Vélez?}{La ausencia de una parte general.}
\cq{¿Qué institutos incorporó la Ley 17.711?}{Lesión subjetiva, teoría de la imprevisión y abuso del derecho, junto con el factor moral, la buena fe y la equidad.}
}{El Código de Vélez, sancionado en 1869 y vigente desde 1871, tuvo a Freitas como principal influencia y careció de parte general. Cien años después, la Ley 17.711 lo reformó sin derogarlo; los proyectos de unificación posteriores desembocaron en el Código Civil y Comercial.}

% ------------------------------------------------------------
\chapter{El Código Civil y Comercial de la Nación}
\label{cap:ccyc}

El Código Civil y Comercial de 2014 no es una reforma como la Ley 17.711: deroga el Código de Vélez y lo sustituye por un cuerpo nuevo. Este capítulo reconstruye su elaboración y trámite, compara su método con el del código anterior y expone sus grandes líneas.

\section{Comisión redactora y elaboración del anteproyecto}

En febrero de 2011, el \textbf{Decreto 191/2011} designó una comisión integrada por dos jueces de la Corte Suprema ---el juez \textbf{Lorenzetti} y la doctora \textbf{Highton de Nolasco}--- y una ex jueza de la Suprema Corte mendocina, la doctora \textbf{Aída Kemelmajer de Carlucci}. Estos tres juristas convocaron a 33 subcomisiones y a la colaboración de 90 profesores, y tomaron como base el proyecto de 1998. Las subcomisiones tuvieron poco diálogo entre sí, lo que se refleja en algunas inconsistencias internas entre tramos del proyecto; la comisión redactora unificó los aportes con la tarea de un coordinador, Federico De Lorenzo, aunque subsisten algunas inconsistencias.

Un año después la comisión presentó al Poder Ejecutivo el anteproyecto, que incluye fundamentos que reemplazan a las notas al pie y deroga completamente el Código de Vélez. El Poder Ejecutivo tomó tres meses adicionales para realizar sus propios ajustes y elevó el proyecto al Congreso el 8 de junio de 2012.

\section{Trámite parlamentario}

El proyecto ingresó por el Senado (expediente S-1157/2012). Se conformó una \textbf{comisión bicameral} de diputados y senadores, que estableció un mecanismo de trabajo con audiencias públicas en todo el país, con participación de particulares, organizaciones no gubernamentales, organizaciones públicas, jueces y abogados, y una expresión federal. Los temas más debatidos fueron la existencia de la persona, la maternidad subrogada, la propiedad indígena, la regulación de los consumidores y diversos temas de familia, propiedad, conjuntos inmobiliarios y otros derechos reales.

La comisión concluyó 2012 sin resultados y el proceso se mantuvo con perfil bajo durante 2013. Hacia octubre-noviembre de ese año se reencauzó: la comisión emitió su dictamen y el Senado dio media sanción el 27 de noviembre de 2013. En Diputados el tratamiento se retomó en septiembre y la sanción definitiva se produjo en la sesión del 1.º de octubre de 2014. El Poder Ejecutivo promulgó y publicó la norma en el Boletín Oficial el 8 de octubre de 2014.

\begin{importante}[Ley 26.994]
Sancionada el 1.º de octubre de 2014 y promulgada el 8 de octubre de 2014, aprueba el Código Civil y Comercial de la Nación. Entró en vigencia el 1.º de agosto de 2015 (la Ley 27.077 adelantó la fecha originalmente prevista).
\end{importante}

\section{Método del Código Civil y Comercial frente al Código de Vélez}

El nuevo código cuenta con un \textbf{título preliminar} y \textbf{seis libros}, mientras que el de Vélez tenía dos títulos preliminares, cuatro libros y un título complementario. Su título preliminar se estudia en el capítulo~\ref{cap:titulo-prel}.

\begin{longtable}{@{}B{0.20\textwidth}L{0.34\textwidth}L{0.34\textwidth}@{}}
\toprule
\textbf{Aspecto} & \textbf{Código de Vélez (1871)} & \textbf{Código Civil y Comercial (2014)} \\
\midrule
\endfirsthead
\toprule
\textbf{Aspecto} & \textbf{Código de Vélez (1871)} & \textbf{Código Civil y Comercial (2014)} \\
\midrule
\endhead
\bottomrule
\endfoot
Títulos preliminares & Dos (leyes y modo de contar intervalos) & Uno, dividido en cuatro capítulos \\
Libros & Cuatro libros más título complementario & Seis libros \\
Parte general & No tenía & Libro Primero: Parte General (novedad metodológica) \\
Personas y familia & Libro I: personas y familia juntos & Separados: Libro I (personas / parte general) y Libro II (familia) \\
Personas humanas y jurídicas & Comenzaba por personas jurídicas & Comienza por personas humanas (Título I) \\
Sucesiones & En Libro IV (disposiciones comunes) & Libro Quinto (libro propio) \\
Notas al pie & Sí (fuentes y fundamentos de Vélez) & Reemplazadas por los Fundamentos del Anteproyecto \\
Unificación civil-comercial & No (Código de Comercio separado) & Sí (unifica materia civil y comercial) \\
\caption{Comparación metodológica entre el Código de Vélez y el Código Civil y Comercial}\\
\end{longtable}

\section{Grandes líneas del Código Civil y Comercial}

\subsection{En materia de persona y familia}

La \textbf{persona humana} está regulada como principio fundamental al inicio del Libro Primero. Se deroga la definición de persona del Código de Vélez, por considerarse una noción previa al derecho: toda persona es un ser humano y no podría existir un ser humano que no sea persona. Se otorga centralidad al concepto de \textbf{dignidad}. Se ratifica que la existencia de la persona comienza desde la concepción, pero se elimina la referencia al seno materno que contenía el código anterior, en atención al problema de los embriones.

Otras novedades en materia de personas son el capítulo de \textbf{derechos personalísimos} y una fuerte reforma en materia de capacidad, a la luz de las Convenciones de los Derechos del Niño y de las Personas con Discapacidad, con tendencia hacia la \textbf{autonomía progresiva} de los niños y de las personas con discapacidad, acompañada de medidas de protección.

En materia de familia se advierte una marcada primacía de la autonomía individual: se debilita la institución matrimonial y aparece el \textbf{divorcio exprés}; en el régimen de bienes del matrimonio se habilita una opción de separación que antes no existía; se regula la \textbf{unión convivencial}, registrada o no; y en materia de filiación se incorporan las técnicas de reproducción asistida bajo la idea de voluntad procreacional, con problemas vinculados al derecho a la identidad. La maternidad subrogada, que había estado en el proyecto inicial, se excluyó del texto final, lo que dejó abierto un debate judicial en el que distintos jueces han dictado sentencias admitiéndola; conforme a las notas de origen, no está admitida en el ordenamiento argentino. % TODO: verificar consistencia entre fuentes originales: la fuente señala a la vez que algunas sentencias la admiten y que no está admitida.
La \textit{patria potestad} pasó a llamarse \textbf{responsabilidad parental}. En materia sucesoria, el cambio más relevante es la reducción de la porción legítima ---la de los descendientes pasó de cuatro quintos a dos tercios, y la de los ascendientes, de dos tercios a un medio--- y se incorporó la posibilidad de mejorar la porción legítima a favor del heredero con discapacidad.

\begin{longtable}{@{}B{0.20\textwidth}L{0.34\textwidth}L{0.34\textwidth}@{}}
\toprule
\textbf{Institución} & \textbf{Código de Vélez} & \textbf{Código Civil y Comercial} \\
\midrule
\endfirsthead
\toprule
\textbf{Institución} & \textbf{Código de Vélez} & \textbf{Código Civil y Comercial} \\
\midrule
\endhead
\bottomrule
\endfoot
Persona humana & Definida en el código (art.~30) & Concepto previo al derecho; se deroga la definición \\
Capacidad & Sistema rígido de incapacidades & Autonomía progresiva y apoyos para personas con discapacidad \\
Derechos personalísimos & Dispersos en el código & Capítulo específico en el Libro I \\
Matrimonio & Divorcio vincular (Ley 23.515) & Divorcio exprés, sin causales ni plazos \\
Régimen de bienes & Solo régimen de comunidad & Opción: comunidad o separación de bienes \\
Unión convivencial & No regulada & Regulada (registrada o no) \\
Filiación & Natural y matrimonial & Se agrega la filiación por técnicas de reproducción asistida \\
Maternidad subrogada & No regulada & Excluida del código (debate judicial abierto) \\
Patria potestad & Patria potestad (art.~264) & Responsabilidad parental \\
Legítima de descendientes & 4/5 & 2/3 \\
Legítima de ascendientes & 2/3 & 1/2 \\
\caption{Comparación de institutos de persona y familia entre ambos códigos}\\
\end{longtable}

\subsection{En materia de derecho privado patrimonial}

Se incorpora parcialmente el estatuto del consumidor y se regulan los contratos de consumo, aunque continúa vigente la ley de defensa del consumidor. En responsabilidad civil se incorpora la \textbf{función preventiva}, conservando la resarcitoria; no se regula la función punitiva, lo que dio lugar a un debate doctrinario. Se amplía la regulación de contratos, incorporando numerosos contratos antes no tipificados como nominados, y se incorporan nuevos derechos reales, como los conjuntos inmobiliarios. Se incorpora también, en el Libro Sexto, un título específico sobre derecho internacional privado.

\begin{importante}[Función preventiva de la responsabilidad civil]
Antes del Código Civil y Comercial solo existía la función resarcitoria. El nuevo código reconoce expresamente la posibilidad de prevenir el daño antes de que ocurra (arts.~1710 a 1713 CCyC).
\end{importante}

\cornell{El Código Civil y Comercial}{
\cq{¿Cómo se elaboró?}{Comisión designada por el Decreto 191/2011 (Lorenzetti, Highton de Nolasco y Kemelmajer de Carlucci), con 33 subcomisiones y 90 profesores, sobre la base del proyecto de 1998; trámite en comisión bicameral con audiencias públicas.}
\cq{¿Cuándo se sancionó y cuándo rige?}{Ley 26.994: sancionada el 1.º de octubre de 2014, promulgada el 8 de octubre de 2014; vigente desde el 1.º de agosto de 2015.}
\cq{¿Qué cambia metodológicamente respecto de Vélez?}{Un título preliminar y seis libros; parte general; persona humana antes que persona jurídica; familia y sucesiones en libros propios; unificación civil-comercial; fundamentos en lugar de notas.}
\cq{¿Qué novedades trae en persona y familia?}{Persona humana sin definición, dignidad central, derechos personalísimos, autonomía progresiva, divorcio exprés, uniones convivenciales, reproducción asistida y responsabilidad parental.}
\cq{¿Qué novedad aporta en responsabilidad civil?}{La función preventiva del daño (arts.~1710 a 1713), además de la resarcitoria.}
}{El Código Civil y Comercial deroga y sustituye al de Vélez, unifica la materia civil y comercial, abre con la persona humana y reconoce la autonomía progresiva, la autonomía individual en familia y la prevención del daño.}

% ============================================================
\part{El marco general del Código Civil y Comercial}
% ============================================================

\chapter{La constitucionalización del derecho privado}
\label{cap:constit}

El Código Civil y Comercial no puede leerse aislado de la Constitución y de los tratados de derechos humanos. Este capítulo explica qué significa constitucionalizar el derecho privado, por qué vías opera y qué ventajas y riesgos se han señalado, siguiendo a Corral Talciani.

\section{Qué significa constitucionalizar el derecho privado}

\begin{definicion}[Constitucionalización del derecho privado]
Influencia que ejercen la Constitución Nacional y los tratados de derechos humanos ---con jerarquía constitucional desde la reforma de 1994 (art.~75 inc.~22 CN)--- sobre el derecho privado, proyectando efectos sobre todas las ramas jurídicas, incluida la del derecho civil.
\end{definicion}

El tema adquiere particular relevancia a partir de la reforma constitucional de 1994, que incorporó con jerarquía constitucional los tratados internacionales de derechos humanos.

\section{Las tres vías de constitucionalización}

Corral Talciani explica que la constitucionalización opera a través de tres vías: la reformadora, la hermenéutica y la de aplicación directa.

\subsection{La vía reformadora}

A partir de los tratados de derechos humanos y de sus disposiciones se generan cambios en el ámbito legislativo.

\begin{ejemplo}[Vía reformadora]
En 2008 Argentina ratificó la Convención sobre los Derechos de las Personas con Discapacidad, cuyo artículo 12 establece exigencias sobre la regulación de la capacidad de ejercicio y de la capacidad jurídica en general. Su recepción llevó a que en 2010 se sancionara una ley que modificó parcialmente el Código de Vélez en lo referido a las sentencias en materia de capacidad de ejercicio. De modo similar, la Convención sobre los Derechos del Niño ---con principios como el interés superior del niño y el derecho del niño a ser oído--- fue generando cambios legislativos sucesivos hasta la sanción del Código Civil y Comercial, que recepta las indicaciones de esa y de otras convenciones.
\end{ejemplo}

\subsection{La vía hermenéutica}

Refiere al modo de interpretar la ley: los principios y los tratados de derechos humanos operan como criterio para determinar su alcance y como pauta al momento de juzgar.

\begin{regla}[Art.~2 CCyC: interpretación]
La ley debe ser interpretada teniendo en cuenta sus palabras, sus finalidades, las leyes análogas, las disposiciones que surgen de los tratados sobre derechos humanos, los principios y los valores jurídicos, de modo coherente con todo el ordenamiento.
\end{regla}

Este artículo recoge en el texto del código la importancia de los tratados de derechos humanos como guía interpretativa.

\subsection{La vía de aplicación directa}

En algunos casos se verifica una aplicación directa de los tratados y de los principios constitucionales. La Constitución de 1994 incorporó, por ejemplo, normas sobre el consumidor y sobre hábeas data, que se aplican generalmente a través de legislación específica ---la Ley de Defensa del Consumidor y la Ley de Protección de Datos Personales---, pero en ciertos casos puede aplicarse directamente la norma constitucional. Esta situación es excepcional, porque el texto constitucional suele aparecer mediado por la legislación y por la tarea del Poder Legislativo.

\begin{ejemplo}[Caso emblemático (2019)]
La Corte Suprema recurrió al principio de la dignidad y al derecho a la vida para modificar, en favor de una persona con discapacidad, el orden en que se cobran los créditos en el régimen de privilegios de la quiebra (Ley de Concursos y Quiebras). Fue un caso de aplicación directa de un principio constitucional para modificar una ley, reconocido como excepcional.
\end{ejemplo}

\section{Ventajas y riesgos de la constitucionalización}

\subsection{Ventajas}

\begin{longtable}{@{}B{0.20\textwidth}L{0.34\textwidth}L{0.34\textwidth}@{}}
\toprule
\textbf{Ventaja} & \textbf{Descripción} & \textbf{Reflejo en el CCyC} \\
\midrule
\endfirsthead
\toprule
\textbf{Ventaja} & \textbf{Descripción} & \textbf{Reflejo en el CCyC} \\
\midrule
\endhead
\bottomrule
\endfoot
Centralidad de la persona y su dignidad & Se coloca a la persona humana en el centro del ordenamiento, resaltando su dignidad. & Art.~51 (dignidad), art.~19 (inicio de la existencia), art.~57 (protección del embrión) \\
Re-materialización del derecho & Incorporación de contenidos sustanciales frente a la idea de un derecho meramente procedimental o formal. & Todo el Título Preliminar y los artículos de la parte general sobre persona y familia \\
Rehabilitación de la dimensión práctica o valorativa & Se recupera la importancia de la argumentación y de los principios en la tarea de juzgar. & Art.~2 (interpretación), art.~3 (deber de decidir de modo razonablemente fundado) \\
La Constitución como fuente del derecho & La Constitución, sus principios y sus valores se transforman en fuente del derecho aplicable. & Art.~1 (fuentes: Constitución y tratados junto a la ley) \\
\end{longtable}

\subsection{Riesgos}

Desde la filosofía jurídica, quienes sostienen una posición iusnaturalista entienden que existe el riesgo de pasar de un positivismo legalista a un positivismo constitucionalista, perdiendo la dimensión de los principios superiores al orden jurídico que iluminan e inciden en la valoración del ordenamiento. Se trata de un debate profundo sobre el rol del juez y sobre la jurisprudencia como fuente: en qué medida el juez crea el derecho o lo aplica cuando existe una ley que regula la situación.

\begin{longtable}{@{}B{0.26\textwidth}L{0.64\textwidth}@{}}
\toprule
\textbf{Riesgo} & \textbf{Descripción} \\
\midrule
\endfirsthead
\toprule
\textbf{Riesgo} & \textbf{Descripción} \\
\midrule
\endhead
\bottomrule
\endfoot
Pérdida de especificidad del derecho civil & La expansión del derecho constitucional puede afectar las metodologías de otras ramas. El derecho civil tiene principios sectoriales propios ---como las reglas para calcular una indemnización--- que pueden perderse cuando todo se constitucionaliza. \\
Positivismo judicial & El juez que, al aplicar la Constitución, se aparta excesivamente de la finalidad de la ley estaría determinando qué es lo obligatorio, alterando la división de poderes y erosionando la construcción republicana. \\
Inseguridad jurídica & Si en las relaciones civiles no se supiera a qué atenerse, la inseguridad resultante paralizaría la vida económica y social, generaría mayores costos y crearía problemas nuevos. \\
Desborde de fuentes del derecho & Por la pérdida de relevancia de la ley aparecen como fuentes fallos de tribunales internacionales, recomendaciones y convenciones de otros sistemas, lo que dificulta determinar el derecho aplicable. \\
Relativismo filosófico & A veces se sostiene que estas posiciones solo son admisibles en una construcción relativista. Desde el iusnaturalismo se sostiene que es posible conocer, a partir de la racionalidad humana, principios fundamentales: que la vida es un bien y que quitar la vida a una persona es un mal. \\
\end{longtable}

\cornell{La constitucionalización del derecho privado}{
\cq{¿Qué es y desde cuándo es relevante?}{La influencia de la Constitución y de los tratados de derechos humanos sobre el derecho privado; adquiere relevancia desde la reforma de 1994 (art.~75 inc.~22 CN).}
\cq{¿Cuáles son las tres vías?}{Reformadora (cambios legislativos), hermenéutica (criterio de interpretación, art.~2 CCyC) y de aplicación directa (excepcional, como el caso de 2019).}
\cq{¿Qué ventajas se señalan?}{Centralidad de la persona, re-materialización del derecho, rehabilitación de la dimensión valorativa y la Constitución como fuente.}
\cq{¿Qué riesgos se señalan?}{Pérdida de especificidad del derecho civil, positivismo judicial, inseguridad jurídica, desborde de fuentes y relativismo filosófico.}
}{La constitucionalización opera por vías reformadora, hermenéutica y de aplicación directa. Aporta centralidad de la persona y de los principios, pero plantea riesgos para la especificidad del derecho civil, la seguridad jurídica y la división de poderes.}

% ------------------------------------------------------------
\chapter{El Título Preliminar y la unificación civil y comercial}
\label{cap:titulo-prel}

El Título Preliminar es el núcleo de significaciones generales del código. Este capítulo describe su estructura y el alcance de la unificación de la materia civil y comercial que el código concreta.

\section{Importancia y estructura del Título Preliminar}

El artículo 1 del CCyC señala la relevancia de la Constitución y de los tratados internacionales de derechos humanos como pautas para resolver los casos regidos por el código.

\begin{nota}
La Constitución puede compararse con la estructura de un edificio: define dónde pueden ir las paredes y cuánto peso pueden soportar. El Código Civil y Comercial es el interior del edificio ---habitaciones, puertas, contratos, familia, herencias---, cuyo diseño debe respetar esa estructura madre. El Título Preliminar es el plano que explica cómo leer todo el edificio: qué fuentes son válidas, cómo interpretar las normas y qué principios prevalecen.
\end{nota}

El Título Preliminar se divide en cuatro capítulos y consta de 18 artículos. Según la elevación del proyecto, estos artículos proyectan sus efectos sobre los seis libros del código.

\begin{table}[H]
\centering
\begin{tabularx}{\textwidth}{B{0.08\textwidth} B{0.22\textwidth} Y B{0.16\textwidth}}
\toprule
\textbf{Cap.} & \textbf{Título} & \textbf{Contenido principal} & \textbf{Artículos} \\
\midrule
1 & El derecho & Fuentes y aplicación. Distinción entre derecho y ley. & 1, 2 y 3 \\
2 & La ley & Ámbito de vigencia temporal y espacial. Efectos de la ley en relación al tiempo. & 4 a 8 \\
3 & Ejercicio de los derechos & Principio del abuso del derecho y principio de buena fe. & 9 a 14 \\
4 & Derechos y bienes & Clasificación de bienes jurídicos, cosas, patrimonio, cuerpo humano. & 15 a 18 \\
\bottomrule
\end{tabularx}
\caption{Estructura del Título Preliminar}
\end{table}

\begin{importante}
El Título Preliminar ofrece un núcleo de significaciones generales para todo el código: en torno a él se encuentran los principios fundamentales que influyen sobre los demás libros.
\end{importante}

\section{La distinción entre derecho y ley}

En los fundamentos del proyecto se distinguió entre «derecho» y «ley», señalando que el derecho es más que la ley. Los tres primeros artículos lo reflejan: el artículo 1 trata las fuentes del derecho (Constitución, tratados, ley, costumbre); el artículo 2, la interpretación de la ley (véase el capítulo~\ref{cap:constit}); y el artículo 3, el deber de emitir una decisión razonablemente fundada.

\begin{regla}[Art.~1 CCyC: fuentes y aplicación]
Los casos que este Código rige deben ser resueltos según las leyes que resulten aplicables, conforme con la Constitución Nacional y los tratados de derechos humanos en los que la República sea parte. A tal efecto, se tendrá en cuenta la finalidad de la norma. Los usos, prácticas y costumbres son vinculantes cuando las leyes o los interesados se refieren a ellos o en situaciones no regladas legalmente, siempre que no sean contrarios a derecho.
\end{regla}

La importancia de «el derecho», como noción que supera a la ley, se enfatiza en el Título Preliminar, aunque en la práctica la ley sigue siendo la principal de las fuentes, como se aprecia en el capítulo 2 del mismo título.

\section{Unificación de la materia civil y comercial}

En la Argentina existieron distintas iniciativas de unificación de la materia civil y comercial (véase el capítulo~\ref{cap:velez}). El Código Civil y Comercial la concreta, aunque, según los autores que han estudiado el tema, se trata de una unificación \textbf{parcial}. Se incorporan la teoría general de las condiciones y los contratos, con sus instituciones comunes, y ciertos temas en materia de actos jurídicos y de contabilidad y registros contables.

\subsection{Materias que permanecen fuera del código}

Continúan reguladas por sus propias leyes especiales: las sociedades comerciales (Ley General de Sociedades), las quiebras y concursos (Ley de Concursos y Quiebras) y la defensa de la competencia. La empresa como institución y el empresario no están incluidos ni regulados específicamente en el código. Tampoco existe en él una definición de la materia comercial: quedó derogado el Código de Comercio, que contenía un artículo específico sobre el acto de comercio y las definiciones de acto de comercio y de comerciante.

\subsection{Contenido comercial incorporado}

\begin{itemize}
\item La contabilidad y la rendición de cuentas.
\item Toda la teoría de la representación (arts.~358 y siguientes).
\item La parte especial de contratos, con numerosos contratos comerciales típicos.
\item Las relaciones de consumo.
\item Las reglas de interpretación de contratos que estaban en el Código de Comercio.
\item El valor de los usos y las costumbres, con mayor presencia como fuente, especialmente en materia de contratos.
\item Los títulos de crédito.
\end{itemize}

% ------------------------------------------------------------
\chapter{Principios del derecho civil: buena fe y equidad}
\label{cap:principios}

Además de la justicia, la dignidad y el bien común (véase la sección~\ref{sec:dignidad-bien-comun}), el derecho civil se apoya en principios que el Título Preliminar proyecta sobre todo el código. Este capítulo estudia los de buena fe y equidad.

\section{El principio de buena fe}

\begin{regla}[Art.~9 CCyC: principio de buena fe]
Los derechos deben ser ejercidos de buena fe.
\end{regla}

Este principio se incorpora en el Título Preliminar y proyecta sus efectos sobre todo el código. Presenta dos grandes dimensiones:

\begin{table}[H]
\centering
\begin{tabularx}{\textwidth}{B{0.17\textwidth} Y Y}
\toprule
\textbf{Dimensión} & \textbf{Descripción} & \textbf{Ejemplo} \\
\midrule
Buena fe objetiva & Conformidad con ciertas reglas; obrar conforme a estándares objetivos de comportamiento. & El poseedor es de buena fe cuando no conoce ---ni puede conocer--- que carece de derecho; por ejemplo, quien compra algo con un vicio en la cadena de transmisión sin saberlo. \\
Buena fe subjetiva & Deber de lealtad; se vincula con el modo en que las partes ejecutan y desenvuelven los contratos: con honestidad, sin ocultar ni retener información relevante. & Deber de honestidad entre contratantes durante las negociaciones y la ejecución; una parte no puede retener información que la otra necesita para decidir. \\
\bottomrule
\end{tabularx}
\caption{Dimensiones de la buena fe}
\end{table}
% TODO: verificar consistencia entre fuentes originales: la fuente asigna el ejemplo del poseedor (desconocimiento de la carencia de derecho) a la «buena fe objetiva» y el deber de lealtad a la «subjetiva»; la denominación podría estar invertida respecto de la doctrina habitual.

\section{El principio de equidad}

Otro principio central es la equidad, vinculada con la justicia en el caso concreto: según la formulación aristotélica, la «rectificación de lo justo legal en el caso concreto». Ya se encontraba en la Ley 17.711 y en el Código Civil anterior; en el CCyC se la encuentra especialmente en los artículos 332 y 1742.

\begin{regla}[Art.~332 CCyC: lesión]
Puede demandarse la nulidad o la modificación de los actos jurídicos cuando una de las partes, explotando la necesidad, debilidad síquica o inexperiencia de la otra, obtuviera por medio de ellos una ventaja patrimonial evidentemente desproporcionada y sin justificación. El juez debe apreciar subjetivamente las circunstancias y puede hacer reajustes equitativos cuando se presenta una situación desproporcionada.
\end{regla}

\begin{regla}[Art.~1742 CCyC: atenuación de la responsabilidad]
El juez, al fijar la indemnización, puede atenuarla si es equitativo en función del patrimonio del deudor, la situación personal de la víctima y las circunstancias del hecho. Esta facultad no es aplicable en caso de dolo del responsable.
\end{regla}

En ambos casos los jueces pueden efectuar rectificaciones o reajustes equitativos, procurando corregir la situación injusta que se presenta en el caso concreto.

\cornell{Título Preliminar, unificación y principios}{
\cq{¿Cómo se estructura el Título Preliminar?}{Cuatro capítulos y 18 artículos: el derecho (1 a 3), la ley (4 a 8), ejercicio de los derechos (9 a 14) y derechos y bienes (15 a 18).}
\cq{¿Qué distinción introduce entre derecho y ley?}{Que el derecho es más que la ley: art.~1 (fuentes), art.~2 (interpretación) y art.~3 (decisión razonablemente fundada).}
\cq{¿Qué alcance tiene la unificación civil y comercial?}{Es parcial: se incorporan contratos, representación, consumo, contabilidad y títulos de crédito; quedan fuera sociedades, concursos y quiebras y defensa de la competencia.}
\cq{¿Qué dimensiones tiene la buena fe?}{Una vinculada a reglas objetivas de comportamiento y otra al deber de lealtad y honestidad en la ejecución de los contratos (art.~9).}
\cq{¿Qué es la equidad y dónde aparece?}{La rectificación de lo justo legal en el caso concreto; arts.~332 (lesión) y 1742 (atenuación de la indemnización).}
}{El Título Preliminar contiene el núcleo de principios del código. Distingue derecho de ley, concreta una unificación parcial de la materia civil y comercial y consagra la buena fe y la equidad como principios que permiten corregir el caso concreto.}

% ------------------------------------------------------------
\chapter{Los efectos de la ley con relación al tiempo}
\label{cap:tiempo}

Este capítulo aborda una cuestión central del derecho civil: qué sucede cuando una ley cambia mientras una relación o situación jurídica se extiende en el tiempo. El problema recorre tres momentos: el sistema del Código de Vélez (1871), la reforma de la Ley 17.711 (1968) y el régimen vigente del artículo 7 del CCyC. La exposición sigue ese orden: primero se plantea el conflicto; luego se examina la solución del código originario, basada en la distinción entre derecho \textbf{adquirido} y \textbf{expectativa}; después se explica por qué esa solución resultó insuficiente y fue reemplazada por el criterio de los hechos \textbf{cumplidos}; y finalmente se muestra que el código vigente conserva esa lógica, con un agregado de protección al consumidor.

\section{Planteamiento del problema}

La pregunta que estructura el capítulo es: ¿qué sucede si hay cambios en la legislación aplicables a situaciones y relaciones jurídicas que se extienden en el tiempo? ¿Puede el legislador dictar normas retroactivas? ¿Con qué límites?

\begin{ejemplo}[El préstamo en dólares]
Una persona presta a otra diez mil dólares a una tasa del cinco por ciento anual, mediante un contrato de mutuo regido por el Código Civil. Mientras la relación está en curso, el Código cambia y la nueva ley dispone que quien prestó dólares debe recibir la devolución en pesos. Se genera un conflicto entre las leyes a lo largo del tiempo: la relación jurídica no nació y concluyó bajo una misma ley, sino que se desarrolló en el tiempo y durante ese desarrollo cambió la norma.
\end{ejemplo}

\subsection{Tres modos de aplicación de la ley en el tiempo}

\begin{enumerate}
\item \textbf{Aplicación inmediata}: la nueva ley rige las consecuencias no cumplidas de relaciones jurídicas en curso.
\item \textbf{Aplicación retroactiva}: la nueva ley pretende regir situaciones jurídicas ya concluidas bajo la ley anterior.
\item \textbf{Aplicación para el futuro}: la nueva ley rige únicamente las relaciones jurídicas que nazcan a partir de su vigencia.
\end{enumerate}

\begin{ejemplo}[Aplicación retroactiva sobre un préstamo ya cumplido]
Si el dinero se prestó en enero de 2015 y se devolvió en julio de 2015, y el nuevo Código ---vigente desde agosto de 2015--- pretendiera aplicarse a lo ya cumplido, habría aplicación retroactiva. La aplicación para el futuro rige solo a partir de la entrada en vigencia del nuevo código, hacia adelante.
\end{ejemplo}

\begin{ejemplo}[El divorcio y los alimentos]
Bajo el Código Civil anterior, el cónyuge declarado culpable en un divorcio podía ser condenado a pagar alimentos al cónyuge inocente. El Código Civil y Comercial eliminó la figura de la culpa en el divorcio y dispuso que ya no correspondía el pago de alimentos por esa causa. Frente a los reclamos de quienes debían seguir pagando alimentos bajo la ley anterior, los jueces entendieron que la nueva ley se aplicaba de inmediato a la situación jurídica en curso, de modo que el obligado dejaba de pagar desde la entrada en vigencia de la nueva norma.
\end{ejemplo}

\section{El sistema del Código de Vélez (1871)}

El código originario (Ley 340) regulaba la materia en su \textbf{artículo 3}, inspirado en el Código Civil francés.

\begin{regla}[Art.~3, Código de Vélez (texto original)]
«La ley ha de disponer para lo futuro; no tiene efecto retroactivo, ni puede alterar los derechos adquiridos.»
\end{regla}

Este artículo consagra el \textbf{principio de irretroactividad}: en principio, una ley rige el presente y el futuro y no alcanza a lo que ya ha ocurrido. Todo el sistema estaba organizado en torno a la noción de \textbf{derecho adquirido}.

\begin{definicion}[Derecho adquirido]
Un derecho está adquirido cuando se cumplieron todos los recaudos que la ley de fondo dispone para que ingrese al patrimonio de la persona, de modo que pueda considerarse su titularidad. Se distingue del derecho en expectativa (art.~4044) y de la mera facultad (art.~4045).
\end{definicion}

\subsection{El derecho en expectativa (art.~4044)}

\begin{regla}[Art.~4044, Código de Vélez]
La nueva ley que entra en vigencia no se considera retroactiva si priva a los particulares de meros derechos en expectativa.
\end{regla}

\begin{definicion}[Derecho en expectativa]
Aquel respecto del cual todavía no se han configurado los presupuestos necesarios para que la adquisición ocurra; por eso no está protegido frente a una nueva ley.
\end{definicion}

\begin{ejemplo}[La expectativa de heredar]
Si una persona tiene la expectativa de heredar a un familiar y una nueva ley cambia el régimen de modo que ya no puede heredar, no puede sostener que la ley sea retroactiva, porque el familiar todavía no ha muerto y no se produjo la adquisición del patrimonio. Distinto sería que el familiar ya hubiera muerto y la nueva ley dispusiera que, de las personas fallecidas en los últimos dos años, solo heredan determinados familiares, excluyendo a quien reclama: allí habría retroactividad sobre un derecho ya adquirido, porque ya se habían cumplido todos los recaudos para que el derecho sucesorio ingresara al patrimonio del heredero.
\end{ejemplo}

\subsection{La mera facultad (art.~4045)}

\begin{regla}[Art.~4045, Código de Vélez]
Si la ley anterior otorgaba una facultad de reclamar algo y esa facultad no fue ejercida, la nueva ley que priva de esa facultad no es retroactiva.
\end{regla}

Quien tenía la facultad y no la ejerció es privado de algo que, en rigor, no llegó a tener.

\subsection{Las excepciones del sistema}

El sistema se completaba con dos artículos tradicionalmente considerados excepciones al principio de irretroactividad: los artículos 4 y 5.

\subsubsection{Las leyes interpretativas (art.~4): una falsa excepción}

\begin{regla}[Art.~4, Código de Vélez]
Las leyes que tengan por objeto aclarar o interpretar otras leyes no tienen efecto respecto de los casos ya juzgados.
\end{regla}

\begin{ejemplo}[Ley interpretativa sobre la tasa de interés]
Si una ley dispone que la tasa de interés no puede superar un monto razonable y al año siguiente otra ley aclara que ese monto no puede superar el cinco por ciento, la segunda interpreta a la primera. La ley interpretativa se aplica a los casos en curso, pero no a los ya juzgados.
\end{ejemplo}

La doctrina la considera una \textbf{falsa excepción}: al ser interpretativa, en realidad se sigue aplicando la ley originaria y solo cambia la interpretación que se le da.

\subsubsection{Las leyes de orden público (art.~5): la verdadera excepción}

\begin{regla}[Art.~5, Código de Vélez]
Ninguna persona puede tener derechos irrevocablemente adquiridos contra una ley de orden público.
\end{regla}

\begin{importante}
El artículo 5 es una excepción genuina: una nueva ley de orden público ---es decir, que expresa principios fundamentales relativos a la subsistencia de la organización social--- podía aplicarse retroactivamente. Esto generaba una considerable inseguridad jurídica, por lo que la Corte Suprema estableció que la ley de orden público no podía afectar derechos como la propiedad, garantizados por la Constitución. La noción de orden público se estudia en el capítulo~\ref{cap:op}.
\end{importante}

\begin{table}[H]
\centering
\begin{tabularx}{\textwidth}{B{0.17\textwidth} Y}
\toprule
\textbf{Artículo} & \textbf{Contenido} \\
\midrule
Art.~3 & Principio de irretroactividad de la ley y protección de los derechos adquiridos. \\
Art.~4044 & El derecho en expectativa no está protegido frente a la nueva ley. \\
Art.~4045 & La mera facultad no ejercida no está protegida frente a la nueva ley. \\
Art.~4 & Las leyes interpretativas se aplican a casos en curso, no a casos juzgados (falsa excepción). \\
Art.~5 & Las leyes de orden público pueden afectar derechos adquiridos (excepción real). \\
\bottomrule
\end{tabularx}
\caption{Síntesis del sistema del Código de Vélez}
\end{table}

\subsection{Crítica al sistema de Vélez}

Una crítica que se volvió clásica, siguiendo al autor francés Roubier, sostiene que la teoría de los derechos adquiridos no servía para explicar cuándo una ley era retroactiva, especialmente en situaciones o relaciones que se prolongan en el tiempo.

\begin{ejemplo}[El préstamo en cuotas]
Si un préstamo se pactó en varias cuotas, la primera cumplida y la segunda pendiente, y cambia la ley disponiendo que en lugar de dólares se pueden entregar pesos, el deudor pretende pagar en pesos y el acreedor cobrar en dólares conforme a la ley anterior. La idea de derecho adquirido no resuelve el conflicto: el acreedor sostendrá que el derecho está adquirido desde el primer momento; el deudor, que es una mera expectativa.
\end{ejemplo}

\section{La reforma de la Ley 17.711 (1968)}

La Ley 17.711 es el punto de inflexión decisivo: en lo sustancial, su sistema continúa rigiendo en la Argentina, pues el CCyC lo mantuvo en su parte central. Modificó el \textbf{artículo 3} y derogó los artículos 4, 5, 4044 y 4045.

\begin{importante}
La Ley 17.711 abandona la noción de derechos adquiridos y adopta la \textbf{teoría de los hechos cumplidos, o del consumo jurídico}, como criterio para determinar la retroactividad de la ley.
\end{importante}

\begin{definicion}[Teoría de los hechos cumplidos o del consumo jurídico]
Un hecho jurídico es una contingencia susceptible de producir consecuencias jurídicas. Cuando un hecho está totalmente cumplido ---se ha consumido, agotó la virtualidad de producir consecuencias jurídicas--- se dice que se ha producido el consumo jurídico. Este criterio permite comprender que existe un elemento dinámico en las relaciones y situaciones jurídicas: los hechos pueden agotar todas sus consecuencias en un instante o cumplirse a lo largo del tiempo.
\end{definicion}

\begin{ejemplo}[Bicicleta y venta en cuotas]
La compra de una bicicleta, en la que se paga y se entrega la cosa en un mismo acto, agota sus consecuencias jurídicas de manera instantánea. Una venta pactada en cuotas se cumple a lo largo del tiempo: la relación se consume recién cuando se paga la última cuota.
\end{ejemplo}

\begin{regla}[Art.~3, Código Civil, según la Ley 17.711]
«A partir de su entrada en vigencia, las leyes se aplican a las consecuencias de las relaciones y situaciones jurídicas existentes. Las leyes no tienen efecto retroactivo, sean o no de orden público, salvo disposición en contrario. La retroactividad establecida por la ley en ningún caso podrá afectar derechos amparados por garantías constitucionales. A los contratos en curso de ejecución no son aplicables las nuevas leyes supletorias.»
\end{regla}

Este artículo incorpora tres principios, que reemplazan a los cinco artículos del sistema originario.

\subsection{Primer principio: efecto inmediato de la nueva ley imperativa}

\begin{definicion}[Efecto inmediato]
A partir de su entrada en vigencia, la nueva ley se aplica de inmediato a las \textbf{consecuencias} de las relaciones y situaciones jurídicas existentes, que se extienden en el tiempo. Se refiere a las \textbf{leyes imperativas}, es decir, las que se imponen a los particulares y que estos no pueden dejar de lado.
\end{definicion}

\begin{ejemplo}[Cuota pendiente de un préstamo]
Si se prestó dinero y la primera cuota fue pagada pero la segunda no, esta última es una consecuencia no cumplida de la relación jurídica; la nueva ley se aplica a esa consecuencia pendiente.
\end{ejemplo}

\subsection{Segundo principio: irretroactividad de la ley}

Este principio se desagrega en cuatro componentes:

\begin{enumerate}
\item Las leyes no tienen efecto retroactivo: se aplican para el futuro, no hacia el pasado.
\item Deja de ser relevante si la ley es o no de orden público, en coincidencia con la derogación del artículo 5.
\item La propia ley puede disponer su retroactividad («excepto disposición en contrario»).
\item La retroactividad que establezca la ley no puede afectar derechos amparados por garantías constitucionales.
\end{enumerate}

Respecto del segundo componente, la irretroactividad y el orden público dejan de confundirse: se trata siempre de leyes imperativas, pero la imperatividad puede o no coincidir con el orden público, tanto en el Código Civil como en el CCyC. En consecuencia, el carácter de orden público ya no resulta decisivo para la retroactividad: la regla general sigue siendo la irretroactividad, salvo disposición expresa en contrario, sea o no la ley de orden público.

\begin{ejemplo}[Retroactividad expresa: filiación por técnicas de reproducción asistida]
Al sancionarse, el CCyC aplicó retroactivamente las normas sobre filiación derivada de la fecundación artificial a personas nacidas con anterioridad al código, para solucionar el problema filiatorio de esas personas. Es un caso de retroactividad establecida expresamente por la ley.
\end{ejemplo}

\begin{importante}
El legislador tiene la potestad de ordenar hacia el pasado, pero con un límite: la retroactividad no puede afectar derechos amparados por garantías constitucionales. Una ley que dispusiera que todos los préstamos en dólares de los últimos dos años deben transformarse a pesos afectaría un derecho de propiedad ya adquirido.
\end{importante}

\subsection{Tercer principio: efecto diferido de la nueva ley supletoria}

\begin{definicion}[Ley imperativa y ley supletoria]
La \textbf{ley supletoria} es aquella que las partes de un contrato pueden dejar de lado mediante pacto en contrario; si nada dicen, se aplica lo que la ley dispone. La \textbf{ley imperativa} se impone a las partes, que carecen de libertad para apartarse de ella.
\end{definicion}

\begin{ejemplo}[La sublocación]
En el CCyC, si las partes de una locación no dicen nada sobre la posibilidad de sublocar, se entiende que sí pueden hacerlo, pero pueden pactar algo distinto. Si el Código cambia esa regla supletoria, el cambio no rige la cláusula que las partes ya integraron a su contrato.
\end{ejemplo}

\begin{ejemplo}[El plazo mínimo de la locación de vivienda]
La locación destinada a vivienda tiene un plazo mínimo de dos años que las partes no pueden dejar de lado. Esa imperatividad es distinta del carácter de la ley supletoria.
\end{ejemplo}

Una nueva ley supletoria tiene, entonces, \textbf{efecto diferido}: no se aplica de inmediato, sino recién a las relaciones jurídicas que se celebren desde su vigencia; no se aplica a un contrato en curso de ejecución.

\section{El régimen vigente: artículo 7 del CCyC}

El CCyC, vigente desde 2015, mantuvo en lo sustancial la redacción de 1968, ahora en su \textbf{artículo 7}.

\begin{regla}[Art.~7 CCyC]
«A partir de su entrada en vigencia, las leyes se aplican a las consecuencias de las relaciones y situaciones jurídicas existentes. Las leyes no tienen efecto retroactivo, sean o no de orden público, excepto disposición en contrario. La retroactividad establecida por la ley no puede afectar derechos amparados por garantías constitucionales. Las nuevas leyes supletorias no son aplicables a los contratos en curso de ejecución, con excepción de las normas más favorables al consumidor en las relaciones de consumo.»
\end{regla}

Los dos primeros principios no presentan cambios. La novedad está en el tercero.

\begin{importante}
Las nuevas leyes supletorias no se aplican a los contratos en curso de ejecución, \textbf{con excepción de las normas más favorables al consumidor en las relaciones de consumo}: si el contrato en curso es de consumo, la nueva norma más favorable al consumidor se aplica de inmediato.
\end{importante}

\begin{nota}
Se ha observado críticamente que este agregado transforma, en los hechos, la norma del consumidor en una norma imperativa, dado que se aplica de inmediato y las normas de protección al consumidor suelen ser imperativas. Resultaría algo inconsistente técnicamente, porque se formula como excepción referida a normas supletorias cuando la norma favorable al consumidor suele ser imperativa. El cambio se explica por la incorporación del derecho del consumidor al nuevo código, ausente en el anterior.
\end{nota}

\begin{table}[H]
\centering
\begin{tabularx}{\textwidth}{B{0.20\textwidth} Y Y}
\toprule
\textbf{Aspecto} & \textbf{Vélez (1871): arts.~3, 4, 5, 4044, 4045} & \textbf{Ley 17.711 (1968) / CCyC (2015): arts.~3 y 7} \\
\midrule
Criterio rector & Derechos adquiridos & Hechos cumplidos / consumo jurídico \\
Irretroactividad & Principio general (art.~3) & Principio general; admite retroactividad expresa con límite constitucional \\
Orden público & Excepción real a la irretroactividad (art.~5) & Irrelevante a estos fines: no se confunde con la imperatividad \\
Leyes interpretativas & Falsa excepción (art.~4) & No reguladas como categoría aparte \\
Leyes supletorias y contratos en curso & No reguladas expresamente & Efecto diferido; excepción a favor del consumidor (solo CCyC, art.~7) \\
\bottomrule
\end{tabularx}
\caption{Comparación de los tres sistemas normativos}
\end{table}

\cornell{Efectos de la ley en el tiempo}{
\cq{¿Cuál es el problema y qué modos de aplicación existen?}{Qué ocurre con la ley cuando una relación jurídica se extiende en el tiempo y cambia la norma. Modos: aplicación inmediata, retroactiva y para el futuro.}
\cq{¿Cómo resolvía Vélez el problema?}{Con la noción de derecho adquirido (art.~3), distinguida de la expectativa (art.~4044) y de la mera facultad (art.~4045), con dos excepciones: leyes interpretativas (falsa excepción) y leyes de orden público (excepción real).}
\cq{¿Por qué fue criticado ese sistema?}{Porque, según Roubier, la teoría de los derechos adquiridos no explica la retroactividad en relaciones que se prolongan en el tiempo (ejemplo del préstamo en cuotas).}
\cq{¿Qué adoptó la Ley 17.711?}{La teoría de los hechos cumplidos o del consumo jurídico, con tres principios: efecto inmediato de la ley imperativa, irretroactividad e efecto diferido de la ley supletoria.}
\cq{¿Qué agrega el art.~7 CCyC?}{Mantiene la redacción de 1968 y agrega la aplicación inmediata de las normas más favorables al consumidor en contratos de consumo en curso.}
}{El sistema de derechos adquiridos de Vélez fue reemplazado en 1968 por el criterio de los hechos cumplidos, articulado en efecto inmediato, irretroactividad (sea o no de orden público, salvo disposición expresa y con límite constitucional) y efecto diferido de la ley supletoria. El art.~7 del CCyC lo conserva y agrega la excepción favorable al consumidor.}

% ------------------------------------------------------------
\chapter{El orden público}
\label{cap:op}

El orden público es uno de los conceptos más transversales del derecho civil: aparece cada vez que una persona firma un contrato, alquila una vivienda, hereda un bien, se casa, se divorcia o negocia en otro país. Comprenderlo permite saber cuándo un acuerdo es válido y cuándo puede ser declarado nulo por afectar principios que la sociedad considera indisponibles, evitar redactar contratos inválidos, asesorar sobre cláusulas abusivas o leoninas y anticipar conflictos en alquileres, sucesiones internacionales y matrimonios celebrados en el extranjero.

Todos los temas del capítulo giran alrededor de una misma idea: la libertad de contratar tiene un límite. Es como las reglas de tránsito: cada conductor es libre de elegir su destino y su camino, pero no puede saltarse un semáforo en rojo porque pone en riesgo a toda la comunidad. Las personas son libres de pactar lo que quieran (autonomía de la voluntad), pero no pueden pasar por encima de ciertos principios básicos de la sociedad, ya sea dentro del país, frente a una ley extranjera o a lo largo del tiempo.

\section{Punto de partida: la libertad de contratar}

Dentro del derecho civil, las personas son libres para hacer contratos, convenciones y acuerdos. Este principio es la capacidad de autonomía, que el derecho civil clásicamente denomina \textbf{autonomía de la voluntad} y que se identifica con la libertad personal consagrada en la Constitución Nacional.

\begin{regla}[Art.~19 Constitución Nacional]
«Nadie puede ser privado de lo que la ley no prohíbe ni obligado a lo que la ley no manda.»
\end{regla}

A partir de este principio, y siguiendo al derecho francés, existe una noción que limita los contratos: el \textbf{orden público}. No es posible celebrar una convención que tenga por finalidad afectar el orden público o que, aunque no se lo proponga, lo afecte de hecho.

El orden público no debe confundirse con el orden en la vía pública. Es un conjunto de principios fundamentales vinculados a la manera en que una sociedad se organiza y que limitan lo que las partes pueden pactar. Existen otros límites, como las leyes imperativas, que no necesariamente coinciden con él.

\section{Teorías sobre el orden público}

A lo largo del tiempo se intentó explicar qué es el orden público mediante distintas teorías:

\begin{itemize}
\item \textbf{Identificación con el derecho público.} Es inexacta, porque también existen normas de orden público dentro del derecho privado.
\item \textbf{Identificación con el interés general.} Aproximación más precisa, pero imperfecta, ya que existen muchas normas de interés general que no son de orden público.
\item \textbf{Identificación con la intuición del intérprete.} Identifica el orden público con lo que el juez debe decidir; no explica qué es, sino solamente quién decide sobre él.
\item \textbf{Identificación con la voluntad del legislador.} El legislador determina cuándo algo es de orden público; sin embargo, la Corte Suprema ha dicho que una ley no es necesariamente de orden público solo porque el legislador lo diga, sino que depende de los principios fundamentales en juego.
\item \textbf{Teoría de Bordas (leyes imperativas).} Para este autor, el orden público se identifica con las leyes imperativas. % TODO: verificar consistencia entre fuentes originales: la fuente escribe «Bordas» para la teoría y «Borda» al citar a los autores actuales.
\end{itemize}

\subsection{Leyes imperativas y leyes supletorias}

Para comprender esa última teoría hay que distinguir dos tipos de leyes. La \textbf{ley imperativa}, que el CCyC llama también indisponible, es la que las partes no pueden dejar de lado sin quebrantarla; de lo contrario, el acto sería nulo. Las \textbf{leyes supletorias} son las que las partes pueden dejar de lado mediante acuerdo (véase el capítulo~\ref{cap:tiempo}).

\begin{table}[H]
\centering
\begin{tabularx}{\textwidth}{YY}
\toprule
\textbf{Leyes imperativas (indisponibles)} & \textbf{Leyes supletorias} \\
\midrule
Las partes no pueden dejar de lado lo dispuesto por la ley; si lo hacen, el acto puede ser nulo. & Las partes pueden dejar de lado lo dispuesto por la ley y acordar algo distinto. \\
Ejemplo: el contrato de alquiler de vivienda no puede celebrarse por un plazo menor a dos años. & Ejemplo: el Código dice que el locatario puede subalquilar salvo disposición en contrario; las partes pueden acordar otra cosa. \\
\bottomrule
\end{tabularx}
\caption{Leyes imperativas y supletorias}
\end{table}

\begin{ejemplo}[Alquiler de vivienda por seis meses]
Quien quiere alquilar su casa es libre de contratar o no con una persona y de fijar el valor, pero el plazo mínimo del alquiler para vivienda es de dos años por una ley imperativa. Si las partes pactan seis meses, el contrato igualmente se entiende celebrado por dos años (sin perjuicio de que, a los seis meses, el inquilino pueda ejercer la facultad de rescindir, que es una cuestión distinta).
\end{ejemplo}

\subsection{Género y especie: la crítica a la teoría de Bordas}

Bordas sostenía que toda ley imperativa es de orden público y que todo orden público es imperativo. Los autores actuales (se citan Borda, Llambías, Rivera, Tobías y Alterini, entre otros) coinciden en general en que la ley imperativa es el género y las leyes de orden público una especie: no todas las leyes imperativas son de orden público, pero todas las de orden público son imperativas.

\begin{ejemplo}[Los tres prenombres]
El Código Civil dice que no se pueden imponer a un hijo más de tres prenombres. Quien quisiera anotar a su hijo con cinco nombres de pila encontraría que el registro le permite solo tres. Es una ley imperativa, pero no necesariamente de orden público: hace a la forma de organización social ---que alguien tenga tres nombres, cuatro o uno solo--- sin comprometer principios esenciales de la sociedad.
\end{ejemplo}

\subsection{Definición mayoritaria y caracteres del orden público}

\begin{definicion}[Orden público]
«Un conjunto de principios esenciales de orden político, jurídico, económico, moral y religioso, sobre los cuales se sustenta la organización social y la persistencia de la misma en el tiempo.»
\end{definicion}

Es un concepto amplio e indeterminado, que no puede encontrarse listado y que cambia con el tiempo. El propio legislador, en ocasiones, aclara que una ley es de orden público ---como la Ley de Defensa del Consumidor--- para enfatizar que no puede ser dejada de lado. Siguiendo a Borda, son dos sus características: es \textbf{eminentemente particular de cada país} (lo que es de orden público en uno puede no serlo en otro) y es \textbf{mutable}, lo que se vincula con el derecho internacional privado.

\begin{nota}[Tres conceptos que no deben confundirse]
La \textit{opinión pública} es un concepto sociológico, de la comunicación, sin implicancias jurídicas. El \textit{derecho público} regula al Estado y a los particulares en su relación con él. El \textit{orden público} es un conjunto de principios que inspiran a una sociedad y limitan lo que las partes pueden disponer; es un concepto propio del derecho privado, vinculado a la limitación de la autonomía de la voluntad.
\end{nota}

\section{Tres grandes supuestos de aplicación}

Pueden identificarse tres supuestos: el orden público frente a la autonomía de la voluntad, frente a la aplicabilidad de la ley extranjera y frente a los efectos de la ley en el tiempo.

\subsection{Orden público y autonomía de la voluntad}

\begin{regla}[Art.~958 CCyC: libertad de contratación]
Las partes son libres para celebrar un contrato y determinar su contenido, dentro de los límites impuestos por la ley, el orden público, la moral y las buenas costumbres.
\end{regla}

Celebrar un contrato significa hacerlo cuando se quiera y con quien se quiera; determinar su contenido también es libre, pero dentro de esos límites. El orden público es un límite distinto de la ley misma: señala principios que van más allá de lo que la ley dice expresamente, aunque a veces ambos se conjugan, cuando la ley declara que cierta materia es de orden público.

\begin{ejemplo}[La compraventa de órganos]
Quien ofrece vender su riñón encuentra un límite impuesto por la ley, que prohíbe la compraventa de órganos, y también por la moral y las buenas costumbres, que consideran que vender órganos afecta un sentido moral.
\end{ejemplo}

\paragraph{El artículo 12 del CCyC.} El artículo 958 se vincula con el artículo 12 del título preliminar, de origen en el Código francés y presente ya en el Código de Vélez.

\begin{regla}[Art.~12, primera parte, CCyC]
Las convenciones particulares no pueden dejar sin efecto las leyes en cuya observancia está interesado el orden público.
\end{regla}

La redacción anterior aludía también a la moral y las buenas costumbres; la mención se quitó, lo que se trataría de una omisión involuntaria, dado que las buenas costumbres subsisten en el artículo 958. El artículo 12 se refiere a algo más intenso que las leyes imperativas comunes: las partes pueden apartarse de las leyes supletorias, pero tampoco de las imperativas en general, aunque esto último no esté dicho expresamente y sea sobreentendido.

\begin{regla}[Art.~279 CCyC: objeto del acto jurídico]
El objeto del acto jurídico no debe ser un hecho imposible o prohibido por la ley, contrario a la moral, a las buenas costumbres, al orden público o lesivo de los derechos ajenos o de la dignidad humana.
\end{regla}

\paragraph{Nulidad de los actos contrarios al orden público.} Otras normas reiteran la lógica:

\begin{regla}[Art.~1004 CCyC: objeto de los contratos]
No pueden ser objeto de los contratos los hechos que son imposibles o están prohibidos por las leyes, son contrarios a la moral, al orden público, a la dignidad de la persona humana, o lesivos de los derechos ajenos.
\end{regla}

\begin{regla}[Art.~1014 CCyC: causa]
Es nulo el contrato cuando su causa es contraria a la moral, al orden público o a las buenas costumbres.
\end{regla}

\begin{regla}[Art.~386 CCyC: nulidad absoluta]
Son de nulidad absoluta los actos que contravienen el orden público, la moral o las buenas costumbres.
\end{regla}

El artículo 279 corresponde al Libro Primero (actos jurídicos como categoría general) y el 1004 al Libro Tercero (contratos): los actos jurídicos son el género y el contrato es una especie. El \textbf{acto jurídico} es el acto voluntario y lícito que tiene por fin inmediato producir efectos jurídicos (art.~259) y puede hacer nacer, modificar o extinguir relaciones o situaciones jurídicas; abarca contratos, testamentos y matrimonios.

\begin{definicion}[Nulidad absoluta y nulidad relativa]
La \textbf{nulidad} es una sanción que consiste en privar a un acto de sus efectos propios por un vicio existente al momento de su celebración. La \textbf{nulidad absoluta} se da cuando el acto contraría el orden público: no puede confirmarse, es imprescriptible y puede ser declarada de oficio si es manifiesta. La \textbf{nulidad relativa} protege un interés particular, por lo que el acto puede eventualmente ser confirmado.
\end{definicion}

\paragraph{Facultades del juez para modificar el contrato.}

\begin{regla}[Art.~960 CCyC]
Los jueces no tienen facultades para modificar las estipulaciones de los contratos, excepto que sea a pedido de una de las partes cuando lo autoriza la ley, o de oficio cuando se afecta, de modo manifiesto, el orden público.
\end{regla}

Es una facultad extraordinaria, muy importante, que no estaba en el código anterior.

\subsection{El fraude a la ley}

\begin{regla}[Art.~12, segunda parte, CCyC]
El acto respecto del cual se invoque el amparo de un texto legal, que persiga un resultado sustancialmente análogo al prohibido por una norma imperativa, se considera otorgado en fraude a la ley. En ese caso, el acto debe someterse a la norma imperativa que se trata de eludir.
\end{regla}

\begin{ejemplo}[El alquiler «turístico»]
Hay una norma que prohíbe alquilar una vivienda por menos de dos años. Quien invoca la norma que permite contratos temporarios por turismo y hace pasar por turismo lo que en realidad es un alquiler para vivienda por seis meses, con el fin de evitar el efecto de la inflación en un contrato largo, hace trampa a la ley. Su alquiler debe sujetarse a la norma imperativa que se trató de eludir: se considera que dura dos años aunque se lo haya denominado alquiler turístico.
\end{ejemplo}

El fraude a la ley no se analiza en el derecho civil como un delito que deba perseguirse penalmente. Puede ser cometido por ambas partes en connivencia, en cuyo caso el juez impone igualmente el orden público, o por una sola de ellas.

\subsection{La renuncia general a las leyes}

\begin{regla}[Art.~13 CCyC]
Está prohibida la renuncia general de las leyes. Los efectos de la ley pueden ser renunciados en el caso particular, excepto que el ordenamiento jurídico lo prohíba.
\end{regla}

Este texto guarda similitud con el artículo 19 de la Constitución Nacional y subyace en él la idea de que una renuncia general a la ley compromete, en general, el orden público.

\subsection{Orden público y ley extranjera}

El segundo supuesto parte del principio de territorialidad:

\begin{regla}[Art.~4 CCyC]
Las leyes son obligatorias para todos los que habitan el territorio de la República Argentina.
\end{regla}

Por las reglas del derecho internacional privado puede, sin embargo, invocarse en ciertas circunstancias la aplicación de una ley extranjera.

\begin{ejemplo}[Matrimonio con bienes en España]
Un matrimonio casado en la Argentina tiene bienes regidos por la ley española. Al divorciarse en la Argentina, el juez argentino aplica la ley española respecto de ese bien. Es un tema típico del derecho internacional privado.
\end{ejemplo}

\begin{regla}[Art.~2600 CCyC]
Las disposiciones de derecho extranjero aplicables deben ser excluidas cuando conducen a soluciones incompatibles con los principios fundamentales de orden público que inspiran el ordenamiento jurídico argentino.
\end{regla}

Su antecedente es el artículo 14 del código anterior, que hablaba de límites a la aplicación de la ley extranjera sin mencionar expresamente el orden público. El actual es más claro.

\begin{ejemplo}[Hijos matrimoniales y extramatrimoniales]
En un país «X» subsisten diferencias entre hijos matrimoniales y extramatrimoniales. Una familia de ese país se instala en la Argentina y se invoca, para la sucesión, una ley en la que el hijo extramatrimonial no hereda. Los jueces no aplican esa ley extranjera porque viola el orden público nacional, en el que los hijos son todos iguales (la herencia es igualitaria desde 1985).
\end{ejemplo}

\begin{nota}[Matrimonios y divorcios en el extranjero]
Un divorcio realizado en el exterior por una pareja casada en la Argentina puede tener validez, pero debe analizarse conforme al derecho internacional privado: normalmente se examina primero la validez del matrimonio (según la ley del lugar de celebración) y luego la disolución del vínculo, que puede regirse por la ley local del divorcio o por la ley de celebración, según el caso. Como ejemplo extremo, el matrimonio de una niña de doce años celebrado en su país de origen no puede inscribirse en la Argentina, porque viola el principio de orden público según el cual el matrimonio de niños está prohibido (en principio, solo con dispensa judicial para menores de 16 años o con autorización de los padres entre los 16 y los 18). El \textit{punto de conexión} es el criterio que determina cuál es la ley aplicable; el juez competente puede ser de un país y el derecho aplicable el de otro.
\end{nota}

\subsection{Orden público y efectos de la ley en el tiempo}

El tercer supuesto fue estudiado en el capítulo~\ref{cap:tiempo}. Bajo el artículo 5 del código de Vélez, las leyes de orden público podían aplicarse retroactivamente incluso contra derechos adquiridos. La reforma de 1968 modificó esa solución y el artículo 7 del CCyC dispone que las leyes no son retroactivas, sean o no de orden público, salvo disposición en contrario: el orden público dejó de ser decisivo respecto de la retroactividad. La Ley 17.711 hizo cambios importantes casi exclusivamente en este supuesto; en autonomía de la voluntad y ley extranjera no hubo mayores cambios.

\section{Cuestiones frecuentes sobre el orden público}

\begin{description}[leftmargin=0pt,style=unboxed]
\item[¿Cómo se sabe si una ley es de orden público?] A veces la ley lo dice expresamente (como la Ley de Defensa del Consumidor); otras veces hay que analizar si la norma protege los principios fundamentales del funcionamiento de la sociedad. Es un tema difícil y, en definitiva, el juez determina en cada caso si algo es de orden público. Se mencionó como ejemplo de actualidad la discusión sobre las normas de pago en dólares: no hay respuesta única y las resuelven los jueces.
\item[¿Todo queda librado a la subjetividad del juez?] No: conforme al artículo 3 del CCyC, todo juez debe dar una decisión razonable y fundada; cuando invoca la moral, las buenas costumbres o el orden público debe explicar por qué, en la situación concreta, se trata de un supuesto de orden público.
\item[¿Varía el orden público con el tiempo?] Sí; las leyes pueden seguir calificándose de orden público mientras la valoración social cambia. En principio se toma como de orden público lo que la ley así califica, y el juez determina si esa calificación se mantiene. Los jueces no pueden dejar de aplicar leyes imperativas ni las partes apartarse de ellas.
\item[¿Son los contratos leoninos contrarios al orden público?] No necesariamente. El contrato leonino abusa de la posición de una parte y se vincula más propiamente con el vicio de lesión (art.~332 CCyC); pueden cruzarse lesión, abuso de derecho y orden público, pero son figuras distintas.
\item[¿Siempre prevalecen las normas de derecho internacional?] No: depende del punto de conexión, y un juez argentino puede aplicar derecho extranjero o un juez extranjero, derecho argentino.
\end{description}

\cornell{El orden público}{
\cq{¿Qué es el orden público y qué límites impone?}{Conjunto de principios esenciales de orden político, jurídico, económico, moral y religioso que sustentan la organización social; limita la autonomía de la voluntad (arts.~958, 12 y 279).}
\cq{¿Cómo se distinguen leyes imperativas y supletorias, y qué relación hay con el orden público?}{Las imperativas no pueden dejarse de lado; las supletorias sí. Toda ley de orden público es imperativa, pero no toda imperativa es de orden público (ejemplo de los tres prenombres).}
\cq{¿Qué consecuencias tiene contravenirlo?}{Nulidad absoluta (art.~386): no confirmable, imprescriptible y declarable de oficio si es manifiesta; el juez puede modificar el contrato de oficio (art.~960).}
\cq{¿Qué es el fraude a la ley?}{Acto que invoca una norma para lograr un resultado análogo al prohibido por una norma imperativa (art.~12); se somete a la norma eludida.}
\cq{¿Cómo opera frente a la ley extranjera?}{Art.~2600: se excluye el derecho extranjero que conduce a soluciones incompatibles con los principios de orden público argentino.}
\cq{¿Qué rol tiene respecto de la retroactividad?}{Ninguno desde 1968: las leyes no son retroactivas, sean o no de orden público, salvo disposición expresa (art.~7).}
}{El orden público es el límite de la autonomía de la voluntad: su violación acarrea nulidad absoluta, el fraude a la ley se corrige sometiendo el acto a la norma eludida, excluye la ley extranjera incompatible y ya no determina la retroactividad de las leyes. Es mutable y propio de cada país, y lo determina el juez con decisión fundada.}

% ============================================================
\part{La persona humana}
% ============================================================

\chapter{La persona humana: fundamento y reconocimiento jurídico}
\label{cap:persona}

La persona humana constituye el centro del derecho civil y comprender su concepto es el punto de partida de toda práctica profesional: cada vez que un abogado redacta un contrato, representa a un cliente, interpreta una norma o litiga, opera sobre la categoría de persona. El concepto de dignidad humana y el reconocimiento de la personalidad no son abstracciones: son las razones por las que el derecho protege la vida, el nombre, el honor, la libertad y el patrimonio. Este capítulo avanza de lo más abstracto a lo más concreto: qué es una persona según la filosofía y la etimología, por qué el derecho la reconoce a partir de los tratados de derechos humanos y de la dignidad, y cómo la regula el CCyC.

\section{Fundamentos filosóficos y etimológicos}

\subsection{El origen etimológico de \emph{persona}}

La palabra \emph{persona} designaba en la antigüedad la máscara del teatro que representaba al personaje. La máscara, en latín, cumplía una función técnica: permitía que la voz del actor sonara más alto (\emph{per-sonare}). El término designaba tanto el objeto ---la máscara--- como el rol o estatus que el personaje representaba. Con el tiempo dejó de designar la máscara y pasó a identificarse con el mismo ser humano que la usaba; ese sentido de rol y estatus se trasladó luego a la filosofía.

\subsection{La definición filosófica de Boecio}

\begin{definicion}[Persona, según Boecio]
«Persona es la sustancia individual de naturaleza racional.»
\begin{itemize}
\item \textbf{Sustancia}: no es un accidente; es un ser que existe por sí mismo.
\item \textbf{Individual}: es indivisible, único y uno; cada persona es irrepetible y no tiene copia.
\item \textbf{De naturaleza racional}: tiene una esencia que implica la capacidad de raciocinio, lo que no significa que ejerza efectivamente esa capacidad en todo momento, sino que esa es su naturaleza.
\end{itemize}
\end{definicion}

\subsection{El ser humano como ser en relación}

Los seres humanos no están solos: vienen a la existencia en una comunidad ---la familia, el barrio, la ciudad, el país, el mundo--- y por eso tienen vínculos de \textbf{alteridad}. De esa dimensión relacional se derivan tres consecuencias jurídicas fundamentales:

\begin{table}[H]
\centering
\begin{tabularx}{\textwidth}{B{0.22\textwidth} Y}
\toprule
\textbf{Consecuencia} & \textbf{Explicación} \\
\midrule
La justicia & Siempre que hay relaciones de alteridad, hay relaciones de justicia: dar a cada uno lo suyo. \\
Los contratos & El ser humano tiene dominio de sí y es capaz de hacer compromisos o pactos vinculantes, lo que funda el orden contractual. \\
La norma jurídica & Se le pueden imponer normas porque tiene racionalidad y reconoce que hay autoridades que deben velar por el bien común. La ley es una ordenación de la razón en orden al bien común, dictada por quien tiene el cuidado de la comunidad. \\
\bottomrule
\end{tabularx}
\caption{Consecuencias jurídicas de la dimensión relacional de la persona}
\end{table}

La palabra \emph{persona} ---a diferencia de \emph{ciudadano}, \emph{individuo} o \emph{habitante}--- expresa un \textbf{valor y una dignidad} y, a la vez, una dimensión jurídica: son seres que se proyectan hacia los demás, individuales y relacionales.

\subsection{Todo ser humano es persona}

\begin{importante}
En el derecho no es posible ---ni legítimo--- distinguir entre ser humano y persona humana. \textbf{Todo ser humano es persona.} La personalidad no es una creación artificial del legislador, sino un dato de la naturaleza que la ley positiva reconoce.
\end{importante}

Se cita a Borda: la persona no nace porque el derecho objetivo le atribuya la capacidad para adquirir derechos y contraer obligaciones, sino que le reconoce esa capacidad. Nadie recibe la personalidad como una concesión al nacer; se es persona por el solo hecho de existir. La elección del verbo \textbf{reconocer} ---en contraposición a \emph{adjudicar} o \emph{conceder}--- admite que la personalidad es anterior al ordenamiento jurídico.

\begin{ejemplo}[El absurdo de los criterios arbitrarios]
Si el derecho dijera que son personas los seres humanos que usan anteojos y que los que no los usan no lo son, ¿cómo se justificaría ese criterio? Sería completamente arbitrario: todo ser humano es persona.
\end{ejemplo}

\section{La persona en los tratados internacionales de derechos humanos}

\subsection{La Declaración Universal de Derechos Humanos (1948)}

La Declaración Universal nació como reacción global ante el Holocausto y la Segunda Guerra Mundial: ante la constatación de que un gobierno había impuesto la degradación de ciertos seres humanos a los que no consideraba personas y los exterminaba, la comunidad internacional afirmó que ningún Estado puede adjudicar la personalidad.

\begin{importante}
Artículo 6: «Todo ser humano tiene derecho, en todas partes, al reconocimiento de su personalidad jurídica.»
Artículo 1: «Todos los seres humanos nacen libres e iguales en dignidad y en derechos.»
\end{importante}

A diferencia de una licencia de conducir ---que el Estado adjudica cuando se reúnen requisitos---, la personalidad jurídica no requiere cumplir ningún requisito: se reconoce por el solo hecho de ser humano.

\subsection{El Pacto de San José de Costa Rica}

\begin{regla}[Convención Americana de Derechos Humanos, art.~1.2]
«Para los efectos de esta Convención, persona es todo ser humano.»
\end{regla}

Lo mismo se encuentra en la Declaración Americana de Derechos y Deberes del Hombre y en el Pacto Internacional de Derechos Civiles y Políticos. Esto influye en la codificación argentina: el CCyC (art.~1) debe respetar los tratados de derechos humanos y, por lo tanto, no puede decir que son personas solo ciertos seres humanos.

\subsection{La dignidad humana: concepto y dimensiones}

La palabra \textbf{dignidad} es central en el sistema de derechos humanos y aparece en el artículo 51 del CCyC. Pueden identificarse tres formas de comprenderla:

\begin{table}[H]
\centering
\begin{tabularx}{\textwidth}{B{0.19\textwidth} Y Y}
\toprule
\textbf{Dimensión} & \textbf{Definición} & \textbf{Ejemplo} \\
\midrule
Funcional / del cargo & Respeto por la función que una persona desempeña en la sociedad. & El presidente de la Nación merece respeto por su cargo, no por ser «más» como ser humano. \\
Moral / de la conducta & Respeto por las acciones virtuosas o rechazo de las conductas indignas. & Los médicos que se expusieron durante la pandemia actuaron con dignidad; quienes lucraron con la crisis, de manera indigna. \\
Ontológica (la más importante para el derecho) & Merecimiento de respeto por el solo hecho de ser humano, independientemente de cualquier cualidad o conducta. & La tienen por igual la persona condenada por el delito más grave y la persona más honesta. \\
\bottomrule
\end{tabularx}
\caption{Dimensiones de la dignidad humana}
\end{table}

\begin{importante}
El concepto de dignidad que se toma es, sobre todo, el de la \textbf{dignidad ontológica}: un merecimiento de respeto por el solo hecho de ser humano, una excelencia en el ser. Es el fundamento de todos los derechos humanos.
\end{importante}

\section{La persona humana en el Código Civil y Comercial}

El CCyC (Ley 26.994) tiene un Título Preliminar de cuatro capítulos (véase el capítulo~\ref{cap:titulo-prel}) y un Libro Primero que comienza con el Título I, dedicado a la persona humana, y el Título II, a la persona jurídica; luego siguen los bienes, los hechos y actos jurídicos y la transmisión de los derechos.

\begin{table}[H]
\centering
\begin{tabularx}{\textwidth}{B{0.20\textwidth} Y Y}
\toprule
\textbf{Aspecto} & \textbf{Código de Vélez (Ley 340)} & \textbf{CCyC (Ley 26.994)} \\
\midrule
Terminología & Persona física / de existencia visible & Persona humana \\
Definición de persona & Art.~30: «entes susceptibles de adquirir derechos y contraer obligaciones» & No incluye definición (se reconoce, no se define) \\
Orden en el libro & Primero personas jurídicas, luego físicas & Primero persona humana (Título I), luego persona jurídica (Título II) \\
Definición de persona física & Art.~51: «entes con signos de humanidad sin distinción de cualidades o accidentes» & Eliminada; se da por supuesta \\
Marco de referencia & Derecho positivo nacional & Tratados de derechos humanos integrados al bloque constitucional \\
\bottomrule
\end{tabularx}
\caption{Terminología y método: Vélez y CCyC en materia de persona}
\end{table}

Esto es toda una definición: el nuevo código coloca en el centro del derecho civil y comercial a la persona humana.

\subsection{Los textos derogados de Vélez}

\begin{nota}[Art.~51 del Código de Vélez (derogado)]
«Personas físicas o de existencia visible son todos los entes que presenten signos característicos de humanidad sin distinción de cualidades o accidentes.»
\end{nota}

Vélez quería decir que todo ser humano es persona. Incluyó «sin distinción de cualidades o accidentes» porque en el derecho romano se discutía si era persona quien nacía con una malformación muy severa ---a quienes se llamaba «monstruosos» o «prodigiosos»--- o si la viabilidad era un requisito; para Vélez, ninguna cualidad o accidente quita la condición de persona.

\begin{nota}[Art.~30 del Código de Vélez (derogado)]
«Personas son todos los entes susceptibles de adquirir derechos y contraer obligaciones.»
\end{nota}

La definición englobaba a la persona humana y a las personas jurídicas, sujetos colectivos surgidos de la asociatividad. El nuevo código la eliminó por considerar que el concepto de persona es anterior al derecho positivo: el derecho no tiene por qué definir quiénes son personas, pues lo son todos los seres humanos. Los fundamentos siguen aquí al Proyecto de 1998, elaborado por una comisión encabezada, entre otros, por Alterini, Alegría, Highton de Nolasco y Méndez Costa.

\subsection{El artículo 51 vigente}

\begin{regla}[Art.~51 CCyC]
La persona humana es inviolable y en cualquier circunstancia tiene derecho al reconocimiento y respeto de su dignidad.
\end{regla}

\section{Cuestiones frecuentes sobre el concepto de persona}

\begin{description}[leftmargin=0pt,style=unboxed]
\item[Persona física y persona no física.] En el Código de Vélez se llamaba persona física o de existencia visible a la persona humana, por oposición a las personas jurídicas o de existencia ideal, que son los entes colectivos: una asociación civil, una sociedad, el Estado, una iglesia, una fundación.
\item[¿Puede un robot ser persona?] El Código de Vélez asociaba la personalidad con la capacidad de adquirir derechos y contraer obligaciones, y no existían robots en su época. El profesor Carlos Muñiz sostiene que el término persona se vincula centralmente con el ser humano, concepto históricamente asociado a él. Los animales merecen respeto, pero no pueden entablar relaciones jurídicas como los humanos; un proyecto de reforma del CCyC habla de nominarlos «seres sintientes» para darles un estatuto especial y regular su protección. El término persona, por involucrar una dignidad vinculada al dominio de sí, no sería fácilmente trasladable a un robot: hoy lo controla alguien, un programador humano que se hace cargo. Si un robot celebrara contratos por sí mismo habría que analizar la analogía con las personas jurídicas, pero siempre habría una persona humana o jurídica detrás.
\item[Persona jurídica y persona moral.] A la persona jurídica se la llama también persona moral o colectiva. El CCyC simplifica el esquema de Vélez en personas humanas y personas jurídicas. Una persona jurídica es una organización con personalidad jurídica, para lo cual hay que tramitar una autorización o que la ley se la conceda; no cualquier asociación lo es, debe asumir uno de los tipos reconocidos (arts.~141 y siguientes, especialmente las formas privadas de los arts.~147 y 148): sociedad anónima o de responsabilidad limitada si tiene forma societaria; fundación o asociación civil si no tiene fines de lucro; o entidad religiosa. La Iglesia Católica tiene personería jurídica pública porque la Constitución le da un estatuto especial (art.~2). El Estado, los Estados extranjeros y organizaciones internacionales como la ONU también son personas jurídicas.
\item[¿Son lo mismo honor y dignidad?] No exactamente. La dignidad es más amplia y es el fundamento de todos los derechos humanos; el honor se refiere a la reputación y tiene una dimensión personal (la honra) y otra social (la fama). Del reconocimiento de la dignidad pueden derivarse el derecho a la vida, al honor y a la libertad.
\end{description}

\cornell{La persona humana}{
\cq{¿Qué significa etimológicamente «persona»?}{Del latín: la máscara teatral que amplificaba la voz (\emph{per-sonare}); pasó a designar el rol y luego al ser humano mismo.}
\cq{¿Cuál es la definición de Boecio?}{«Sustancia individual de naturaleza racional»: ser que existe por sí mismo, único e irrepetible, de esencia racional aunque no la ejerza en todo momento.}
\cq{¿Por qué la persona tiene dimensión jurídica?}{Es un ser en relación (alteridad): de ello surgen la justicia, los contratos y la sujeción a normas.}
\cq{¿Todo ser humano es persona? ¿Qué importa el verbo «reconocer»?}{Sí, sin excepción. La personalidad es un dato de la naturaleza que el derecho reconoce, no adjudica ni concede; no depende de ningún requisito estatal.}
\cq{¿Qué establecen los tratados?}{DUDH (1948), art.~6: derecho al reconocimiento de la personalidad jurídica; art.~1: libres e iguales en dignidad. Pacto de San José, art.~1.2: persona es todo ser humano.}
\cq{¿Qué es la dignidad ontológica?}{El merecimiento de respeto por el solo hecho de ser humano, independiente de cargo, conducta o cualidad; fundamento de los derechos humanos.}
\cq{¿Qué cambió el CCyC?}{Persona humana en lugar de persona física; sin definición de persona; abre con la persona humana (Título I) y consagra la inviolabilidad y dignidad (art.~51).}
}{Todo ser humano es persona por el solo hecho de serlo; el derecho reconoce, y no concede, esa personalidad. El CCyC suprime las definiciones de Vélez, coloca a la persona humana al inicio del Libro Primero y consagra en el artículo 51 su inviolabilidad y su dignidad ontológica, fundamento de todos los derechos.}

% ------------------------------------------------------------
\chapter{El comienzo de la existencia: la persona por nacer}
\label{cap:por-nacer}

La cuestión del comienzo de la existencia de la persona posee dimensiones biológicas, filosóficas, morales y jurídicas. Este capítulo se concentra exclusivamente en la \textbf{dimensión civil}, dejando de lado las repercusiones propias del derecho penal u otras ramas, y recorre su evolución desde el derecho romano hasta el debate actual sobre el artículo 19 del CCyC y los embriones no implantados.

\section{La tradición romana}

\begin{definicion}[\textit{Nasciturus pro iam nato habetur}]
Al que está concebido o por nacer se lo tiene como si ya hubiera nacido: se lo equipara jurídicamente al nacido. Es el antecedente más remoto de la protección jurídica del concebido.
\end{definicion}

\begin{ejemplo}[Herencia del hijo concebido antes de la muerte del padre]
Si un padre fallece mientras su hijo está concebido pero no ha nacido, el derecho romano concluyó que el hijo hereda igualmente, aunque nazca después de la muerte del causante. Se consideró injusto excluirlo, dado que ya era hijo por haber sido concebido antes de la muerte y puede haber hasta nueve meses entre concepción y nacimiento.
\end{ejemplo}

\begin{ejemplo}[Libertad del hijo de una mujer que cae en esclavitud]
Cuando una mujer concebía siendo libre y luego caía en esclavitud, el momento de la concepción ---y no el del nacimiento--- determinaba el estatus del hijo, que nacía libre.
\end{ejemplo}

\section{El Código de Vélez (1871)}

El principio llega al Código argentino a través de los artículos 70 y 63.

\begin{nota}[Art.~70 del Código de Vélez (texto histórico)]
Desde la concepción en el seno materno comienza la existencia de la persona física.
\end{nota}

Vélez recoge la tradición romanista y la lleva a su máxima expresión: mientras otros códigos ---como el chileno--- disponen que la persona existe desde que nace, aunque el concebido se equipare al nacido a efectos de su protección, Vélez declaró directamente que la persona existe desde la concepción. La categoría de «persona por nacer» proviene de Teixeira de Freitas.

\begin{nota}[Art.~63 del Código de Vélez (texto histórico)]
Son personas por nacer las que, no habiendo nacido, están concebidas en el seno materno.
\end{nota}

\section{La reforma constitucional de 1994 y los tratados}

En 1994 se incorporaron tratados de derechos humanos que refuerzan la protección del concebido:

\begin{itemize}
\item \textbf{Art.~75 inc.~23 CN}: el Congreso tiene potestad para dictar un régimen de seguridad social especial e integral en protección de la madre y el niño, \emph{desde el embarazo}, lo que revela que existe una etapa prenatal protegida.
\item \textbf{Ley 23.849} (ratificación de la Convención sobre los Derechos del Niño): declaración interpretativa con valor de reserva, según la cual niño es todo ser humano desde la concepción hasta los 18 años; integra las condiciones de vigencia de la Convención y por tanto es texto constitucional para interpretar el comienzo de la existencia.
\item \textbf{Pacto de San José, art.~4}: todo ser humano tiene derecho a la vida, en general, a partir del momento de la concepción; de la interpretación conjunta resulta que persona es todo ser humano.
\item \textbf{Pacto Internacional de Derechos Civiles y Políticos}: prohíbe aplicar la pena de muerte a mujeres embarazadas, lo que implica un reconocimiento implícito de la personalidad del no nacido.
\end{itemize}

\section{El Código Civil y Comercial}

\subsection{El debate de elaboración y el texto del artículo 19}

Importa distinguir el proyecto elevado por el Poder Ejecutivo del texto sancionado. El proyecto establecía que la existencia comenzaba con la concepción en el seno materno y que, en las técnicas de reproducción humana asistida, comenzaba con la implantación del embrión en la mujer: preveía \textbf{dos momentos distintos}, de modo que el mismo embrión sería persona dentro del cuerpo de la mujer pero no fuera de él. Esa solución generó un intenso debate en la Comisión Bicameral y el criterio de la implantación fue eliminado en el debate parlamentario.

\begin{regla}[Art.~19 CCyC (Ley 26.994)]
La existencia de la persona humana comienza con la concepción.
\end{regla}

Se eliminó la expresión «en el seno materno» del código anterior. La postura que ubica el comienzo en el nacimiento resultó minoritaria en el derecho argentino y no hubo proyecto de ley en tal sentido.

\begin{importante}
El problema abierto por el artículo 19 no es si el concebido es persona por nacer ---esa tradición se mantiene sin discusión---, sino qué ocurre con los \textbf{embriones no implantados}, generados \textit{in vitro}, fuera del seno materno, mientras no son transferidos.
\end{importante}

\subsection{Concepción: ¿fecundación o implantación?}

\paragraph{Primera postura: concepción equivale a fecundación.} Identifica la concepción con la unión del óvulo y el espermatozoide, que forma un nuevo organismo individual (el cigoto). La sostienen la Cámara Federal de Salta, la Suprema Corte de Mendoza, las Jornadas de Derecho Civil de 2003, el artículo 17 del CCyC y profesores como Sambrizzi y Ambrosi.

\begin{regla}[Art.~17 CCyC: cuerpo humano (según la explicación de las fuentes)]
Existe una protección del cuerpo humano desde el momento en que este se inicia, que es la fecundación. El cuerpo es el soporte de la persona y se forma cuando se fusionan el óvulo y el espermatozoide: surge un nuevo cuerpo, distinto del de los padres y del de los gametos, aunque la fusión sea \textit{in vitro}.
\end{regla}

\paragraph{Segunda postura: concepción equivale a implantación.} Se funda en el fallo \textit{Artavia Murillo}, defendida también por Marisa Herrera y Andrés Gil Domínguez.

\begin{ejemplo}[Caso Artavia Murillo y otros c.\ Costa Rica (CIDH, noviembre de 2012)]
Matrimonios con infertilidad, sin gametos de terceros, solicitaron acceder a la fecundación \textit{in vitro}, que Costa Rica no permitía por considerar que ponía en riesgo a embriones y que en el sistema interamericano «concepción» equivalía a fecundación. La Corte Interamericana entendió que Costa Rica se había extralimitado, sostuvo que para esas circunstancias el concepto debía equipararse a «implantación» y condenó al país.
\end{ejemplo}

\paragraph{Consideraciones críticas.} El contexto de \textit{Artavia Murillo} es distinto del argentino: en Costa Rica la fecundación \textit{in vitro} estaba totalmente prohibida y en la Argentina no. Las sentencias de la Corte Interamericana obligan solo a las partes del litigio; para los demás países son una pauta de interpretación de la que se puede apartar con argumentación fundada, en el marco del control de convencionalidad. El fallo es una sentencia, no una ley; la norma sigue empleando «concepción» y lo discutido es su interpretación. Interpretar «concepción» como «implantación» resulta, según esta crítica, contrario al principio \textit{pro homine}, pues la palabra remite semánticamente al inicio de algo; y el artículo 1 del CCyC excluyó deliberadamente a la jurisprudencia como fuente.

\begin{importante}
La Corte Suprema argentina ha señalado que los pronunciamientos de la Corte Interamericana son una pauta de interpretación, de la que cabe apartarse cuando el ordenamiento, las circunstancias o el contexto son distintos. Usar \textit{Artavia Murillo} ---referido a la creación de embriones mediante fecundación asistida--- para autorizar la destrucción de embriones ya existentes sería trasladar sus conclusiones a una circunstancia distinta de la litigada. En los precedentes de la Cámara Federal de Salta se discutieron cuestiones de cobertura y de diagnóstico genético preimplantatorio que \textit{Artavia Murillo} no abordó.
\end{importante}

\subsection{Otras cuestiones sobre embriones}

El proyecto original contemplaba la inseminación artificial \textit{post mortem}, pero la previsión fue eliminada del texto sancionado, entre otras razones por los problemas éticos que plantea. Si el embrión ya existe, existe también una persona ---según esta postura--- y la cuestión pasa a ser garantizar su derecho a la vida. El destino de los embriones sobrantes de la fecundación \textit{in vitro} es un problema abierto que las fuentes no desarrollan en profundidad: se señala que lo habitual es procurar no generar embriones congelados, que una vez generados no hay solución éticamente indiscutida, y se afirma, sin mayor desarrollo, que su destrucción no es una solución adecuada.

\section{La época de la concepción (art.~20 CCyC)}

El artículo 20 retoma los artículos 76 y 77 del código anterior.

\begin{regla}[Art.~20 CCyC]
Época de la concepción es el lapso entre el máximo y el mínimo fijados para la duración del embarazo. Se presume, excepto prueba en contrario, que el máximo de tiempo del embarazo es de trescientos días y el mínimo de ciento ochenta, excluyendo el día del nacimiento. Este plazo se cuenta retroactivamente desde el día de ese nacimiento.
\end{regla}

Proviene de una larga tradición legislativa; el máximo y el mínimo se fijan para resolver situaciones vinculadas principalmente a la determinación de la paternidad y a conflictos sucesorios.

\begin{ejemplo}[El soldado en guerra]
Un soldado casado parte a la guerra el 1.º de enero de 2020 y regresa el 30 de junio de 2021. Su esposa le presenta un hijo nacido en mayo de 2021. Contando nueve meses hacia atrás desde el nacimiento, la concepción habría ocurrido mientras él estaba en la guerra, por lo que no podría reconocerse la paternidad. Antes de las pruebas genéticas, este cálculo era la principal herramienta para determinarla.
\end{ejemplo}

\begin{ejemplo}[El heredero inventado]
Fallece una persona el 1.º de octubre de 2020 sin hijos y sus padres inician la sucesión. Se presenta una mujer que dice tener un hijo concebido por el fallecido, nacido el 31 de diciembre de 2020. Contando retroactivamente 300 días desde esa fecha se llega a una época en la que el causante ya había fallecido, por lo que no puede presumirse a ese hijo como heredero, sin perjuicio de que la presunción admite prueba biológica en contrario.
\end{ejemplo}

El nuevo código eliminó los artículos 65, 66, 67, 68 y 78 del Código de Vélez, referidos a la postergación de compromisos y a los controles admitidos durante el embarazo.

\section{El nacimiento con vida (art.~21 CCyC)}

El artículo 21 subsume los antiguos artículos 71 a 75.

\begin{regla}[Art.~21 CCyC]
Los derechos y obligaciones del concebido o implantado en la mujer quedan irrevocablemente adquiridos si nace con vida. Si no nace con vida, se considera que la persona nunca existió. El nacimiento con vida se presume.
\end{regla}

Esta condición ya existía en el código anterior y no implica que el concebido sea «menos persona»: su finalidad histórica es evitar fraudes.

\begin{ejemplo}[Heredero simulado por embarazo ficticio]
Fallece un hombre sin descendencia y una mujer dice estar embarazada de él. No pueden practicarse estudios genéticos a la mujer embarazada sin su consentimiento, por ser una persona con derechos propios; el embarazo puede acreditarse con su simple declaración. Si el hijo por nacer falleciera durante el embarazo, ella lo heredaría y excluiría a los padres del causante. La condición del nacimiento con vida ---si nace sin vida se considera que nunca existió--- evita este tipo de conflictos sucesorios simulados.
\end{ejemplo}

\subsection{Qué significa nacer con vida}

Nacer con vida implica quedar completamente separado de la madre. El código anterior admitía la prueba por amplitud de medios; el vigente presume que se nace con vida, y quien sostenga lo contrario debe probarlo: por ejemplo, si los pulmones llegaron a expandirse, si alguien escuchó llorar al recién nacido o si se cortó el cordón umbilical. Si el niño muere dentro del útero, la figura penal aplicable es la de aborto; si muere fuera, la de homicidio. A los fines civiles, la cuestión es relevante centralmente en materia sucesoria.

\begin{ejemplo}[Caso Píparo]
En una salidera bancaria se dispara contra una mujer embarazada; el hijo por nacer nace y fallece poco después. Al haber nacido con vida ---aunque brevemente---, el hecho fue calificado como homicidio y no solo como lesión a la madre. Ilustra, en el plano penal, las consecuencias del artículo 21, aunque excede el ámbito civil.
\end{ejemplo}

\section{El régimen jurídico de la persona por nacer}

\begin{longtable}{@{}B{0.20\textwidth}L{0.70\textwidth}@{}}
\toprule
\textbf{Norma} & \textbf{Contenido} \\
\midrule
\endfirsthead
\toprule
\textbf{Norma} & \textbf{Contenido} \\
\midrule
\endhead
\bottomrule
\endfoot
Art.~24 CCyC & Las personas por nacer son incapaces de ejercicio, pero capaces de derecho: no pueden actuar por sí mismas, pero sí ser titulares de derechos y adquirirlos, razón por la cual necesitan representación. \\
Art.~101 CCyC & Sus representantes son sus padres, de la misma manera en que representan a la persona ya nacida menor de edad. \\
Art.~2279 CCyC & Pueden ser herederas las personas concebidas que nazcan con vida; también las concebidas mediante técnicas de fecundación \textit{in vitro}, si se cumple el requisito del consentimiento del art.~561. Los bienes heredados no se inscriben a su nombre hasta que nazca con vida, para evitar conflictos sucesorios. \\
Art.~57 CCyC & Se prohíben las prácticas destinadas a modificar genéticamente el embrión cuando la alteración se transmita a su descendencia; reconoce plenamente al embrión como persona, lo que es relevante para interpretar «concepción». \\
Art.~592 CCyC & Puede imponerse preventivamente la filiación del hijo por nacer. \\
Art.~574 CCyC & Puede reconocerse prenatalmente al hijo por nacer. \\
Alimentos & El hijo por nacer tiene derecho a alimentos, desde la concepción; la mujer embarazada puede reclamarlos al progenitor. Consisten en una suma de dinero destinada al mantenimiento y la crianza de los hijos. \\
\end{longtable}

\section{Evolución del reconocimiento de la persona por nacer}

\begin{table}[H]
\centering
\small
\begin{tabularx}{\textwidth}{B{0.23\textwidth} Y B{0.26\textwidth}}
\toprule
\textbf{Sistema / norma} & \textbf{¿Qué establece?} & \textbf{¿Desde cuándo hay persona?} \\
\midrule
Derecho romano & \textit{Nasciturus pro iam nato habetur} & Desde la concepción (para efectos específicos) \\
Código de Vélez (art.~70) & Existencia desde la concepción en el seno materno & Desde la concepción en el seno materno \\
Código Civil chileno & El concebido se equipara al nacido & Desde el nacimiento (con protección prenatal) \\
Art.~75 inc.~23 CN (1994) & Seguridad social especial desde el embarazo & Reconoce al niño prenatal \\
Ley 23.849 (CDN) & Niño: todo ser humano desde la concepción hasta los 18 años & Desde la concepción \\
Pacto de San José (art.~4) & Derecho a la vida, en general, desde la concepción & Desde la concepción (con matiz «en general») \\
Art.~19 CCyC (2015) & La existencia comienza con la concepción & Desde la concepción (debate: fecundación o implantación) \\
\textit{Artavia Murillo} (CIDH, 2012) & Concepción = implantación, para el sistema interamericano & Desde la implantación (pauta interpretativa) \\
\bottomrule
\end{tabularx}
\caption{Evolución normativa del reconocimiento de la persona por nacer}
\end{table}

\cornell{El comienzo de la existencia}{
\cq{¿Qué significa \textit{nasciturus pro iam nato habetur} y qué aplicaciones tuvo?}{Que el concebido se tiene por nacido: hereda aunque nazca tras la muerte del padre y su libertad depende de la concepción, no del nacimiento.}
\cq{¿Cómo recogió Vélez el principio?}{Arts.~70 y 63: la existencia comienza desde la concepción en el seno materno; categoría de «persona por nacer» tomada de Freitas.}
\cq{¿Qué tratados refuerzan la protección en 1994?}{Art.~75 inc.~23 CN, Ley 23.849 (niño desde la concepción), Pacto de San José (art.~4) y Pacto de Derechos Civiles y Políticos.}
\cq{¿Qué dice el art.~19 y qué problema deja abierto?}{La existencia comienza con la concepción; el problema son los embriones no implantados: ¿concepción es fecundación o implantación?}
\cq{¿Qué regulan los arts.~20 y 21?}{Época de la concepción (entre 180 y 300 días antes del nacimiento) y condición del nacimiento con vida, que se presume; si no nace con vida, se considera que la persona nunca existió.}
\cq{¿Cuál es el régimen de la persona por nacer?}{Incapaz de ejercicio y capaz de derecho (art.~24), representada por sus padres (art.~101), puede heredar (art.~2279), tiene alimentos, reconocimiento prenatal (art.~574) y protección frente a la manipulación genética (art.~57).}
}{La existencia de la persona humana comienza con la concepción, tradición que va del derecho romano al art.~19 CCyC. El debate actual gira en torno a si concepción es fecundación o implantación. La persona por nacer es capaz de derecho pero incapaz de ejercicio, y sus derechos quedan irrevocablemente adquiridos si nace con vida.}

% ============================================================
\part{Los atributos de la persona humana}
% ============================================================

\chapter{Los atributos de la persona humana}
\label{cap:atributos}

Tras estudiar la persona humana y el comienzo de su existencia, el derecho civil se ocupa de sus atributos: los elementos que permiten que la personalidad jurídica se despliegue. Los atributos posibilitan que las personas, en tanto sujetos en relación, establezcan relaciones y situaciones jurídicas y puedan identificarse dentro de ellas.

\section{Enumeración de los atributos}

La doctrina clásica reconoce cinco atributos de la persona humana: el nombre, el domicilio, el estado, la capacidad y el patrimonio.

\begin{table}[H]
\centering
\begin{tabularx}{\textwidth}{B{0.20\textwidth} Y}
\toprule
\textbf{Atributo} & \textbf{Función} \\
\midrule
Nombre & Permite identificar a la persona. \\
Domicilio & Permite ubicar a la persona. \\
Estado & Permite establecer frente a qué relaciones jurídicas se encuentra la persona. \\
Capacidad & Permite determinar qué aptitud tiene la persona para entablar relaciones jurídicas. \\
Patrimonio & Permite conocer qué bienes tiene una persona. \\
\bottomrule
\end{tabularx}
\caption{Atributos de la persona humana y su función}
\end{table}

\begin{nota}
Algunos autores discuten la clasificación clásica, en particular la inclusión del patrimonio como atributo. Estas notas siguen la distinción clásica de nombre, domicilio, estado, capacidad y patrimonio. Los capítulos siguientes estudian el nombre (capítulo~\ref{cap:nombre}), el domicilio (capítulo~\ref{cap:domicilio}), el estado civil (capítulo~\ref{cap:estado}) y la capacidad (capítulo~\ref{cap:capacidad}).
\end{nota}

% ------------------------------------------------------------
\chapter{El nombre}
\label{cap:nombre}

El nombre es el primer atributo en el orden de exposición: permite identificar a la persona. Este capítulo define el nombre, expone su régimen legal, las reglas sobre prenombre y apellido según el tipo de filiación, el cambio de nombre y sus acciones de protección, y explica la filiación como presupuesto de la atribución del apellido.

\section{Concepto y caracteres}

\begin{definicion}[Nombre]
Es la identificación exclusiva que corresponde a una persona.
\end{definicion}

El nombre es un concepto intuitivo, regulado de manera clara en el CCyC. Sus caracteres son: \textbf{necesario, obligatorio, único, inalienable, inembargable, imprescriptible, inmutable e indivisible}.

Está compuesto por dos elementos:

\begin{itemize}
\item el \textbf{nombre individual}, de pila o familiar, que el Código denomina \textbf{prenombre}: la denominación que permite distinguir a la persona dentro de su familia;
\item el \textbf{apellido}: la denominación común a los integrantes de una misma familia.
\end{itemize}

Coloquialmente se habla de «nombre y apellido» por practicidad, pero jurídica y técnicamente ambos integran el concepto único de \textbf{nombre}. Al prenombre se lo conoce como «nombre de pila» por la pila bautismal, tradición de raíz cristiana; el código actual utiliza la palabra técnica «prenombre».

\section{Régimen legal: evolución legislativa}

La regulación del nombre se ubica en el Libro Primero, Título 1 (La persona humana), Capítulo 4 (El nombre), artículos 62 a 72 del CCyC, y deroga la Ley 18.248, conocida como «ley del nombre».

\begin{table}[H]
\centering
\begin{tabularx}{\textwidth}{B{0.30\textwidth} Y}
\toprule
\textbf{Etapa} & \textbf{Régimen} \\
\midrule
Código de Vélez (Ley 340) & Sin regulación específica del nombre; se regía por la costumbre. \\
Ley 18.248 (1969) & Primera «ley del nombre», complementaria del Código de Vélez. \\
CCyC vigente & Incorpora la regulación en su texto (arts.~62 a 72) y deroga la Ley 18.248. \\
\bottomrule
\end{tabularx}
\caption{Evolución legislativa de la regulación del nombre}
\end{table}

\section{El nombre como derecho y como deber (art.~62)}

\begin{regla}[Art.~62 CCyC]
La persona humana tiene el derecho y el deber de usar el prenombre y el apellido que le corresponde.
\end{regla}

La doctrina mayoritaria sostiene que el uso del nombre es simultáneamente un derecho y un deber. Es un \textbf{derecho subjetivo}, porque la persona puede exigir que se la reconozca con un determinado nombre. Es también un \textbf{deber}, porque existe un interés público en identificar a las personas; como correlato, la persona debe identificarse con su propio nombre, sin usar un nombre falso ni mentir al respecto.

\section{El prenombre (art.~63)}

\subsection{Adquisición y orden de elección}

\begin{regla}[Art.~63 inc.~a) CCyC]
El prenombre se adquiere por la inscripción en el acta de nacimiento, a pedido de quien ostente el derecho a elegirlo.
\end{regla}

El orden de prelación de quienes eligen el prenombre es el siguiente:

\begin{enumerate}
\item Los \textbf{padres}, en conjunto o uno solo de ellos, sobreentendiéndose en este caso la autorización del otro (art.~641 CCyC, sobre responsabilidad parental).
\item Los padres pueden \textbf{delegar} la elección en un tercero, otorgándole autorización.
\item Si uno de los padres no está presente, elige el otro (por ejemplo, una persona soltera).
\item Si no hay padre conocido (por ejemplo, un niño abandonado), eligen los guardadores designados judicialmente, el Ministerio Público o el funcionario del Registro Civil y de Capacidad de las Personas.
\end{enumerate}

\subsection{Reglas y límites}

Rige un principio de \textbf{libertad de elección}, que incluso permite inscribir nombres aborígenes o derivados de voces aborígenes, autóctonas o latinoamericanas. Los límites son:

\begin{importante}
No pueden inscribirse más de tres prenombres, ni apellidos como prenombres, ni primeros prenombres idénticos a los de hermanos vivos, ni prenombres extravagantes.
\end{importante}

\begin{ejemplo}[Apellidos como prenombres]
No es posible inscribir a un hijo con el prenombre «Messi» o «Maradona», porque son apellidos.
\end{ejemplo}

La prohibición de repetir el primer prenombre entre hermanos vivos no se aplica a los nombres compuestos, como María de las Mercedes, María del Rosario, María del Pilar, Juan Bautista, Juan Segundo, Juan Manuel o Juan Ignacio.

Respecto de los nombres extravagantes, la ley anterior a 1948 era más amplia: no podían inscribirse nombres extravagantes o ridículos, que expresaran tendencias políticas o ideológicas, o que fueran equívocos respecto del sexo de quien los llevara. Hoy ese conjunto se sintetiza en la expresión «extravagantes», que deja un margen de discusión: un nombre como «Revolución» probablemente lo sería, y lo mismo los nombres de marcas famosas como «Amazon» o «Google», práctica que existió en Estados Unidos. Quien define en la práctica si un nombre es extravagante es el Registro del Estado Civil; ante un conflicto, la persona interesada puede recurrir a la vía judicial.

\section{El apellido}

\subsection{Filiación matrimonial (art.~64)}

El CCyC innovó respecto de la Ley 18.248, que establecía que el hijo llevaba necesariamente el apellido paterno.

\begin{regla}[Art.~64 CCyC: apellido de los hijos matrimoniales]
El hijo matrimonial lleva el primer apellido de alguno de los cónyuges; en caso de no haber acuerdo sobre qué apellido llevará, se determina por sorteo realizado en el Registro del Estado Civil y Capacidad de las Personas. A pedido de los padres, o del interesado con edad y madurez suficiente, se puede agregar el apellido del otro cónyuge. Todos los hijos comunes deben llevar el mismo apellido y, en el mismo orden, que el primero de los hijos.
\end{regla}

El orden de apellidos fijado para el primer hijo ---por acuerdo o por sorteo--- queda establecido para todos los hijos comunes posteriores. Para agregar el apellido del otro progenitor a pedido del interesado, el Código usa el estándar de \textbf{edad y madurez suficiente}, que evalúa en concreto el Registro Civil; como criterio orientativo, se señala que debería contar al menos con trece años, edad en que el Código presume discernimiento para los actos lícitos, sin perjuicio del análisis de cada caso. El artículo 64 solo permite agregar apellidos, no suprimirlos.

\subsection{Filiación extramatrimonial}

\begin{enumerate}
\item \textbf{Un solo progenitor reconoce al hijo} (por ejemplo, solo la madre): el hijo lleva el apellido de quien lo reconoce.
\item \textbf{El segundo progenitor reconoce al hijo con posterioridad:} deben ponerse de acuerdo; si no lo logran, decide el juez teniendo en cuenta el interés superior del niño.
\item \textbf{Ambos reconocen al mismo tiempo:} rigen las reglas de la filiación matrimonial; si no hay acuerdo, sorteo.
\end{enumerate}

\subsection{Persona sin filiación determinada (art.~66)}

Una persona puede carecer de filiación determinada, lo que suele ocurrir en zonas marginales, en familias sin contacto con instituciones estatales o de protección, donde nunca fue inscripta aunque la comunidad la conozca por un apellido. No se puede excluir a un niño de la educación escolar por no contar con DNI. En estos casos, el oficial del Registro Civil anota la circunstancia utilizando el apellido que ya use la persona o asignándole un apellido común.

\begin{regla}[Art.~66 CCyC]
La persona, con edad y grado de madurez suficiente, que carezca de apellido inscripto, puede pedir la inscripción del que está usando.
\end{regla}

\subsection{Apellido de los cónyuges, viudez y adopción}

\paragraph{Durante el matrimonio.} El Código consagra la libertad de que cualquiera de los cónyuges opte por usar el apellido del otro, con la partícula «de» o sin ella, como apellido adicional al propio. Es un derecho indistinto, a diferencia del régimen anterior en que era la mujer quien tomaba el apellido del marido; el texto actual está redactado en términos neutros.

\paragraph{Divorcio y nulidad.} Quien adoptó el apellido del otro deja de usarlo, salvo que por motivos razonables un juez autorice a conservarlo.

\begin{ejemplo}[Motivo razonable: uso comercial]
Un cónyuge que usaba el apellido adquirido por matrimonio para identificar a su empresa continuó usándolo tras el divorcio por razones económicas, dado que era el nombre con el que su empresa era reconocida en el mercado.
\end{ejemplo}

\paragraph{Viudez.} El cónyuge viudo puede seguir usando el apellido del otro mientras no contraiga nuevas nupcias ni constituya una nueva unión convivencial; su uso no es obligatorio.

\paragraph{Adopción.} El régimen varía según el tipo:

\begin{table}[H]
\centering
\begin{tabularx}{\textwidth}{B{0.22\textwidth} Y}
\toprule
\textbf{Tipo de adopción} & \textbf{Apellido del adoptado} \\
\midrule
Plena & Lleva el apellido del adoptante; si la adopción es conjunta, rigen las reglas de la filiación matrimonial sobre el orden. Excepcionalmente, fundado en el derecho a la identidad y a pedido del interesado, puede agregarse o anteponerse el apellido de origen (identidad en su dimensión estática). \\
Simple & No corta el vínculo con la familia de origen. El adoptado puede solicitar mantener su apellido de origen adicionándole el del adoptante; si no manifiesta voluntad, se aplica la regla de la adopción plena. La consulta sobre mantener el apellido de origen es la regla, mientras que en la plena es excepcional. \\
De integración & Adopción del hijo del cónyuge o conviviente: se aplican las reglas del tipo de adopción de que se trate (plena o simple). \\
\bottomrule
\end{tabularx}
\caption{Apellido del hijo adoptado}
\end{table}

\section{Cambio de nombre y acciones de protección}

\begin{importante}
El nombre es inmutable. El cambio de nombre y apellido solo procede si existen justos motivos, y puede tramitarse por vía judicial o administrativa.
\end{importante}

\subsection{Vía judicial: justos motivos (art.~69)}

Debe alegarse y acreditarse un justo motivo. El artículo 69 enumera algunos, a título orientativo:

\begin{regla}[Art.~69 CCyC: justos motivos (enumeración no taxativa)]
El seudónimo, cuando hubiere adquirido notoriedad; la raigambre cultural, étnica o religiosa; la afectación de la personalidad de la persona interesada, cualquiera sea su causa, siempre que se encuentre acreditada.
\end{regla}

El \textbf{seudónimo notorio} se da habitualmente entre artistas. La \textbf{raigambre cultural, étnica o religiosa} se refiere a quien descubre una raíz propia y desea reflejarla en su nombre. La \textbf{afectación de la personalidad} se vincula con nombres que han provocado bullying, burlas o vergüenza social. Se trata de una lista enunciativa: puede haber otros justos motivos con fuerza suficiente para convencer a un juez, en el proceso judicial específico.

\subsection{Vía administrativa}

También prevista en el artículo 69, se aplica en dos supuestos en que la ley autoriza el cambio sin acudir a un juez: por razón de identidad de género, conforme a la Ley 26.743; y por haber sido víctima de desaparición forzada, sustitución de identidad, alteración o supresión del estado civil, o de un delito que afecte la personalidad. Se recurre directamente al Registro Civil.

\subsection{Procedimiento y acciones de protección}

El procedimiento está regulado en el artículo 70 y las acciones de protección en el artículo 71:

\begin{table}[H]
\centering
\begin{tabularx}{\textwidth}{B{0.26\textwidth} Y Y}
\toprule
\textbf{Acción (art.~71)} & \textbf{Procede cuando\ldots} & \textbf{Ejemplo} \\
\midrule
Reclamación & Se desconoce a una persona el uso de su propio nombre. & Una universidad anota incorrectamente el nombre; la persona reclama que se lo anote bien. \\
Impugnación o usurpación & Otra persona usa el nombre de quien reclama, causándole un perjuicio. & Alguien abre una cuenta en una red social con el nombre de otra y actúa en su nombre. \\
Defensa o supresión & El uso del nombre por un tercero causa un perjuicio a quien lo lleva. & El nombre aparece como nombre de fantasía en una serie o telenovela. \\
\bottomrule
\end{tabularx}
\caption{Acciones de protección del nombre}
\end{table}

\begin{definicion}[Seudónimo]
Nombre ficticio empleado por una persona, generalmente en el ámbito artístico, en el ámbito de despliegue de su actividad. Su uso es lícito. Si goza de notoriedad, tiene la misma protección que el nombre y le corresponden las acciones del artículo 71.
\end{definicion}

\section{La filiación y el nombre}
\label{sec:filiacion}

La regulación del apellido depende de la filiación, que conviene precisar.

\begin{definicion}[Filiación]
Es el estatuto de hijo: un vínculo jurídico entre progenitores e hijos, que tiene su origen en la naturaleza, en un reconocimiento jurídico por el cual una persona reconoce a su hijo extramatrimonial, en la fecundación \textit{in vitro} (con regulación propia dentro del derecho de familia) o en la adopción.
\end{definicion}

\paragraph{Filiación matrimonial.} Cuando dos personas se casan, los hijos que tengan se presumen del matrimonio. La presunción puede no coincidir con la realidad biológica en todos los casos, pero la ley la establece por razones de orden social: el matrimonio es ordenador de la vida social, los cónyuges se comprometen a vivir juntos y a educar a los hijos que nazcan, y la sociedad reconoce ese compromiso otorgando a esos hijos el apellido fijado y presumido legalmente.

\paragraph{Filiación por naturaleza.} Cuando una mujer da a luz, la maternidad queda determinada por el parto. Si la madre está casada, la ley presume la paternidad del cónyuge; si no lo está, la persona tiene en principio un solo progenitor determinado, salvo que otra persona la reconozca. En estos casos la ley sigue un modelo binario que sigue a la biología.

\paragraph{Filiación extramatrimonial.} Puede ser fruto de un reconocimiento voluntario ---técnicamente un acto jurídico de reconocimiento--- o de una acción judicial.

\begin{definicion}[Acción de filiación forzosa]
Acción judicial por la cual se obliga al padre que no reconoce voluntariamente a su hijo a reconocerlo y a pasar alimentos; al asignarse el vínculo, el hijo pasa a tener determinados a ambos progenitores.
\end{definicion}

Frente a la filiación determinada existe la filiación no determinada, en la que se desconoce el vínculo y se usa el apellido que ya emplea el hijo (art.~66).

\paragraph{El interés superior del niño.} Cuando el juez resuelve atendiendo al interés superior del niño ---por ejemplo, ante el desacuerdo sobre el apellido--- se remite centralmente a la Convención sobre los Derechos del Niño, sin perjuicio de la Ley 26.061. Es un concepto complejo con un componente objetivo ---determinar qué es lo mejor para ese niño en concreto--- que también tiene en cuenta la opinión del niño según su edad y madurez.

\cornell{El nombre}{
\cq{¿Qué es el nombre y cuáles son sus caracteres?}{La identificación exclusiva de una persona. Es necesario, obligatorio, único, inalienable, inembargable, imprescriptible, inmutable e indivisible; se compone de prenombre y apellido.}
\cq{¿Por qué es derecho y deber (art.~62)?}{Es derecho subjetivo a ser reconocido con un nombre y deber por el interés público en identificar a las personas.}
\cq{¿Qué límites tiene la elección del prenombre?}{No más de tres, no apellidos, no el primer prenombre de hermanos vivos (salvo compuestos) y no extravagantes; decide el Registro y cabe vía judicial.}
\cq{¿Cómo se determina el apellido del hijo?}{Matrimonial: primer apellido de cualquiera de los cónyuges, por acuerdo o sorteo, igual orden para los hermanos. Extramatrimonial: según quién reconoce y cuándo. Adopción: según sea plena, simple o de integración.}
\cq{¿Cuándo procede el cambio de nombre?}{Solo con justos motivos (art.~69): vía judicial (seudónimo notorio, raigambre cultural, étnica o religiosa, afectación de la personalidad) o administrativa (identidad de género, víctimas de ciertos delitos).}
\cq{¿Qué acciones protegen el nombre?}{Reclamación, impugnación o usurpación, y defensa o supresión (art.~71).}
}{El nombre, compuesto por prenombre y apellido, identifica a la persona y es un derecho y un deber regulados en los arts.~62 a 72 CCyC. La libertad de elección del prenombre tiene límites; el apellido depende de la filiación; el cambio exige justos motivos y el nombre se protege con acciones específicas.}

% ------------------------------------------------------------
\chapter{El domicilio}
\label{cap:domicilio}

El domicilio es el atributo que permite ubicar a la persona. Este capítulo lo define y lo distingue del domicilio constitucional y de la residencia, expone su clasificación en general y especial, desarrolla el domicilio real y el legal, los tipos de domicilio especial, las reglas del domicilio desconocido y los efectos jurídicos del domicilio.

\section{Concepto, finalidad y ubicación normativa}

El domicilio es un concepto jurídico civil específico, que no debe confundirse con el domicilio que figura en el documento nacional de identidad (DNI). Este último es una forma de registración, mientras que el domicilio depende de reglas propias del derecho civil. Está regulado en el Libro Primero, Título Primero (La persona humana), Capítulo Quinto (El domicilio), artículos 73 a 78 del CCyC.

\begin{definicion}[Domicilio]
Es el asiento jurídico de la persona: el lugar donde, a los efectos jurídicos correspondientes, se considera que se encuentra una persona.
\end{definicion}

El domicilio resulta relevante para:

\begin{itemize}
\item \textbf{determinar la ley aplicable} a una situación (por ejemplo, la ley aplicable a un inmueble es la del territorio donde se encuentra);
\item \textbf{determinar la jurisdicción del juez competente} (en una sucesión, el juez del último domicilio de la persona fallecida);
\item \textbf{practicar válidamente las notificaciones};
\item \textbf{determinar el lugar de cumplimiento de las obligaciones} (en un préstamo, el deudor devuelve el dinero en el domicilio del acreedor).
\end{itemize}

\subsection{Domicilio civil y domicilio constitucional}

Cuando la Constitución se refiere a la inviolabilidad del domicilio protege la vida privada, la vivienda y la imposibilidad de realizar allanamientos sin orden judicial, junto con una serie de recaudos vinculados a la concepción política constitucional.

\begin{importante}
El domicilio civil sirve a los fines de las relaciones jurídicas ---cumplimiento de obligaciones, ley aplicable, competencia judicial y notificaciones---; el domicilio constitucional protege la inviolabilidad de la vivienda y la vida privada frente a injerencias como los allanamientos sin orden judicial.
\end{importante}

El domicilio civil surge de la ley: a veces coincide con la residencia real de la persona y otras veces la ley lo asigna en función de una determinada función que la persona desempeña en una situación dada.

\subsection{Domicilio y residencia}

La \textbf{residencia} es el lugar donde, en los hechos, se encuentra una persona, y puede ser habitual o actual. Es un concepto de hecho que puede o no coincidir con el domicilio jurídico.

\begin{ejemplo}[Residencia actual o habitación]
Quien se va de vacaciones y habita transitoriamente una casa en Mar del Plata tiene allí su residencia actual (también llamada habitación), aunque su domicilio continúe en la ciudad donde reside habitualmente. La residencia actual no implica, por sí sola, un cambio de domicilio.
\end{ejemplo}

\section{Clasificación: domicilio general y especial}

\begin{definicion}[Domicilio general]
Rige para la generalidad de las relaciones o situaciones jurídicas de una persona. Debe ser uno solo y puede ser real o legal.
\end{definicion}

\begin{definicion}[Domicilio especial]
Tiene efectos para una o algunas relaciones jurídicas particulares determinadas. Una persona puede tener varios.
\end{definicion}

Dentro del domicilio general, el \textbf{domicilio real} es el que la persona elige y donde tiene el asiento jurídico de su familia y sus negocios; el \textbf{domicilio legal} es el que la ley asigna en ciertos casos.

\begin{advertencia}
No es correcto identificar el «domicilio legal» con el domicilio que figura en el DNI. Normalmente toda persona tiene un domicilio real, que es voluntario; solo en circunstancias previstas expresamente por la ley corresponde asignarle un domicilio legal.
\end{advertencia}

\begin{ejemplo}[Múltiples domicilios especiales]
Una persona puede celebrar un préstamo fijando un domicilio y, al día siguiente, otro préstamo fijando uno distinto, pues puede constituir tantos domicilios especiales como quiera. La multiplicidad puede complicar en la práctica la administración de los efectos jurídicos, ya que cada contraparte enviará notificaciones o pagos al domicilio fijado en su contrato.
\end{ejemplo}

\begin{importante}
El domicilio general es \textbf{necesario} ---ninguna persona puede carecer de él--- y \textbf{único}: más de uno generaría confusión sobre dónde ubicar a la persona a todos los efectos jurídicos.
\end{importante}

\section{El domicilio real (art.~73)}

\begin{definicion}[Domicilio real (art.~73, primer párrafo, CCyC)]
Es el lugar donde la persona tiene el lugar de su residencia habitual.
\end{definicion}

El artículo 73 contempla dos supuestos distintos: el primer párrafo regula el domicilio real y el segundo, un domicilio especial (el profesional, tratado más abajo). Son instituciones distintas, reguladas en el mismo artículo.

\paragraph{Caracteres.} El domicilio real es \textbf{voluntario}: depende de la persona constituirlo, mantenerlo, exigirlo o modificarlo, en virtud del principio de libre elección del domicilio. La facultad de cambiar de domicilio no puede ser coartada por contrato ni por una imposición de última voluntad.

\begin{ejemplo}[Imposibilidad de coartar el cambio de domicilio]
Unos padres regalan a su hijo una casa de casamiento bajo la condición de que nunca se mude de allí. La condición no puede disponerse válidamente: no es posible obligar a una persona a no cambiar nunca su domicilio.
\end{ejemplo}

\paragraph{Elementos: corpus y ánimo.}

\begin{definicion}[Corpus]
La residencia efectiva en un lugar: la materialidad de estar viviendo en un lugar determinado.
\end{definicion}

\begin{definicion}[Ánimo]
La intención de permanecer en un lugar y de constituir allí el domicilio; el elemento subjetivo por el cual la persona quiere estar en ese lugar.
\end{definicion}

\begin{ejemplo}[El estudiante de otra provincia]
Una persona de Santa Rosa que se traslada a Buenos Aires para cursar una carrera reside allí transitoriamente, sin intención de cambiar su domicilio de origen: conserva su domicilio real en Santa Rosa, aunque su habitación esté en Buenos Aires.
\end{ejemplo}

\paragraph{Constitución, mantenimiento y extinción.} El domicilio real se constituye cuando concurren corpus y ánimo.

\begin{regla}[Art.~77 CCyC: permanencia y cambio]
El domicilio real se mantiene mientras subsista alguno de sus dos elementos (corpus o ánimo). El cambio de domicilio se produce instantáneamente por el hecho de trasladar la residencia de un lugar a otro con ánimo de permanecer en ella.
\end{regla}

Una persona puede mudarse físicamente conservando el ánimo de regresar al domicilio anterior; para el cambio efectivo se requieren ambos elementos. El domicilio real se extingue por la constitución de un nuevo domicilio o porque la ley impone un domicilio legal.

\section{El domicilio legal (art.~74)}

El domicilio legal, a diferencia del real, es \textbf{forzoso}: la ley lo impone.

\begin{definicion}[Domicilio legal (art.~74 CCyC)]
Es el lugar donde la ley presume, sin admitir prueba en contra, que la persona reside de manera permanente para el ejercicio de sus derechos y el cumplimiento de sus obligaciones.
\end{definicion}

Se trata de una presunción \textit{iuris et de iure}: se impone a la persona y no le permite demostrar lo contrario. Sus caracteres son:

\begin{itemize}
\item \textbf{forzoso}: lo impone la ley;
\item \textbf{de interpretación restrictiva}: no cabe la analogía para crear supuestos no previstos;
\item \textbf{excepcional}: en principio toda persona tiene domicilio real;
\item \textbf{único}, por ser una especie de domicilio general.
\end{itemize}

\begin{regla}[Art.~74 CCyC: supuestos de domicilio legal]
\begin{itemize}
\item Los funcionarios públicos tienen su domicilio en el lugar donde deben cumplir sus funciones, sea que las cumplan de manera permanente, periódica o por comisión.
\item Los militares en servicio activo, en el lugar en que están prestando servicio.
\item Los transeúntes o las personas que no tienen domicilio conocido, en el lugar de su residencia actual.
\item Las personas incapaces, en el domicilio de su representante.
\end{itemize}
\end{regla}

\begin{ejemplo}[Los cuatro supuestos]
\textbf{Funcionarios}: un ministro de la Corte Suprema o del Poder Ejecutivo tiene domicilio donde cumple sus funciones. \textbf{Militares}: en el lugar donde prestan servicio activo. \textbf{Transeúntes}: personas cuya profesión o manera de vida implica traslados constantes o en situación de calle; por ejemplo, los integrantes de un circo tienen domicilio donde residan en cada momento. \textbf{Incapaces}: los menores tienen su domicilio en el de su representante; si un menor tiene un derecho en una sucesión, las notificaciones se dirigen al domicilio de sus representantes.
\end{ejemplo}

El régimen anterior contemplaba un caso adicional, el del servicio doméstico agregado a la casa de sus empleadores, que fue eliminado del Código actual.

\section{El domicilio especial}

\subsection{Domicilio convencional o de elección}

\begin{definicion}[Domicilio especial convencional]
Es el que las partes de un contrato eligen para el ejercicio de los derechos y obligaciones que de él emanan. Las notificaciones enviadas a ese domicilio son válidas aunque la persona ya no resida allí, salvo que se modifique el contrato o se notifique y acepte el cambio.
\end{definicion}

\begin{ejemplo}[La cláusula de domicilio en un contrato]
En el encabezado de un contrato se identifica a cada parte con un domicilio donde serán válidas todas las notificaciones. Ante una discusión derivada del contrato, la notificación enviada a ese domicilio tendrá efecto aunque la persona se haya mudado y alegue que ya no vive allí.
\end{ejemplo}

El fundamento reside en la fuerza vinculante de los contratos: el acuerdo de respetar como válidas las notificaciones al domicilio pactado debe respetarse mientras no se modifique el contrato.

\subsection{Domicilio procesal}

El domicilio procesal, o \textit{ad litem}, presenta hoy una modalidad electrónica, utilizada a través del sistema de notificaciones del Poder Judicial, que coincide con el número de CUIT del profesional interviniente. Existe además el domicilio constituido, que generalmente coincide con el del abogado interviniente.

\begin{ejemplo}[Domicilio constituido en distintos procesos]
Una persona que vive en un domicilio determinado constituye, para un proceso concreto, domicilio en el estudio de su abogado; si tiene otro abogado en un juicio distinto, tendrá en ese otro proceso un domicilio especial diferente.
\end{ejemplo}

\subsection{Domicilio profesional (art.~73, segundo párrafo)}

\begin{definicion}[Domicilio profesional]
La persona que ejerce actividad profesional o económica tiene su domicilio en el lugar donde la desempeña, para el cumplimiento de las obligaciones emergentes de dicha actividad.
\end{definicion}

Comprende las profesiones liberales: contador, abogado, ingeniero, psicólogo, psiquiatra, médico.

\begin{ejemplo}[El abogado con un problema de familia]
Un abogado con un problema sobre el régimen de cuidado de sus propios hijos no es notificado donde ejerce su profesión sino en su domicilio real, porque el problema no se vincula con su actividad profesional.
\end{ejemplo}

\section{Domicilio desconocido (art.~76)}

\begin{regla}[Art.~76 CCyC]
La persona cuyo domicilio no es conocido lo tiene en el lugar donde se encuentra; y si este también se ignora, en su último domicilio conocido.
\end{regla}

La regla es relevante en la práctica, especialmente con deudores que procuran no ser ubicados; adquiere importancia el domicilio registrado en el DNI o en los trámites electorales.

\begin{ejemplo}[El deudor que no aparece]
Si se desconoce el domicilio de un deudor y las cédulas enviadas al último domicilio conocido son rechazadas, corresponde solicitar un oficio al Registro de las Personas para que informe el último domicilio conocido y notificar allí.
\end{ejemplo}

Frente a esa notificación, la persona puede alegar que ya no vivía allí y pedir que se la declare inválida junto con las actuaciones derivadas; se abre entonces un proceso para determinar si la persona realmente no residía allí o si quien notificó conocía el verdadero domicilio.

\begin{importante}
Distribución de la carga de la prueba: si quien notificó afirma que no conocía el verdadero domicilio, se aplica el artículo 76 y corresponde a la otra parte demostrar la mala fe de quien notificó; si no lo logra, se considera debidamente notificado.
\end{importante}

En principio, una notificación que llega a un domicilio donde la persona ya no vive no surte efectos, pues debe llegar efectivamente al lugar donde reside. Cuando se han intentado distintos domicilios sin éxito, se procura notificar en el del DNI y, si tampoco resulta posible, se notifica bajo responsabilidad de la parte actora para continuar válidamente con el proceso.

\section{Efectos jurídicos del domicilio (art.~78)}

\begin{regla}[Art.~78 CCyC]
El domicilio determina la competencia de las autoridades en las relaciones jurídicas. Allí debe notificarse a la persona, intimarla, demandarla o requerirla extrajudicialmente, sin perjuicio de lo dispuesto en leyes especiales.
\end{regla}

Incluye la elección de domicilio respecto de la competencia: si una persona constituye un domicilio en un contrato, ello produce efectos frente a la competencia del juez ante el que podrá ser demandada.

\section{Precisiones adicionales}

\begin{description}[leftmargin=0pt,style=unboxed]
\item[Personas privadas de su libertad.] La situación depende del caso: en principio tendrán el domicilio donde residen actualmente, aunque la cuestión se rige más por las reglas penitenciarias y de ejecución de la pena. Cuando el traslado fue contrario a su voluntad no se configuran, en rigor, corpus y ánimo, por ausencia de voluntad libre de constituir allí su residencia. % TODO: verificar consistencia entre fuentes originales: la fuente comienza señalando que el domicilio especial no equivale al de una persona detenida, sin precisar el supuesto.
\item[Argentinos residentes en el exterior.] Pueden trasladar su domicilio al exterior; para la ley argentina dejan de tener domicilio en el país. Si el último domicilio de una persona fallecida estaba en el exterior, la sucesión se tramita allí. La materia se vincula con el derecho internacional privado, donde el domicilio es un punto de conexión.
\item[Domicilio procesal y cambio de domicilio.] Constituido domicilio en un proceso, las cédulas siguen dirigiéndose allí mientras no se informe un cambio en el expediente, momento en que se constituye uno nuevo. Distinto es el caso de una notificación ---como una carta documento--- enviada antes de constituir domicilio en un proceso: si se dirige a un lugar donde la persona ya no vive, en principio no surte efectos.
\end{description}

\cornell{El domicilio}{
\cq{¿Qué es el domicilio y para qué sirve?}{El asiento jurídico de la persona (arts.~73 a 78). Determina ley aplicable, competencia judicial, notificaciones y lugar de cumplimiento de las obligaciones.}
\cq{¿En qué se distingue de la residencia y del domicilio constitucional?}{La residencia es un hecho (habitual o actual); el domicilio constitucional protege la vivienda y la vida privada; el DNI es solo una forma de registración.}
\cq{¿Cuál es la clasificación del domicilio?}{General (único y necesario; real o legal) y especial (varios; convencional, procesal, profesional).}
\cq{¿Qué elementos tiene el domicilio real y cómo se mantiene?}{Corpus (residencia efectiva) y ánimo (intención de permanecer); se mantiene mientras subsista alguno y cambia cuando se trasladan ambos (art.~77).}
\cq{¿Cuáles son los cuatro supuestos de domicilio legal?}{Funcionarios públicos, militares en servicio activo, transeúntes e incapaces (en el domicilio de su representante); presunción \textit{iuris et de iure}, interpretación restrictiva.}
\cq{¿Qué ocurre si el domicilio se ignora?}{Se tiene por domicilio el lugar donde la persona se encuentra y, si también se ignora, el último domicilio conocido (art.~76).}
}{El domicilio es el asiento jurídico de la persona. El general es único y puede ser real (voluntario: corpus y ánimo) o legal (forzoso, en cuatro supuestos); el especial puede ser múltiple (convencional, procesal, profesional). Si se ignora, rige el último domicilio conocido, y sus efectos alcanzan a la competencia, las notificaciones y el cumplimiento de obligaciones.}

% ------------------------------------------------------------
\chapter{El estado civil}
\label{cap:estado}

El estado civil es el atributo que permite establecer frente a qué relaciones jurídicas se encuentra la persona. Este capítulo define el estado civil, repasa sus antecedentes romanos y su sentido en la era de los derechos humanos, estudia las acciones de estado, la posesión de estado con sus cuatro aplicaciones y la prueba del estado civil mediante las partidas del Registro Civil.

\section{Concepto de estado civil}

El término «estado» es polisémico: habitualmente remite a la organización política de una nación, pero esa es solo una de sus acepciones.

\begin{definicion}[Estado civil]
En derecho civil, el \textbf{estado} es la posición que la persona ocupa en la sociedad y en la familia. Se lo denomina estado civil y se asocia con las categorías de casado, soltero o divorciado ---más vinculadas a lo matrimonial---, pero también comprende la \textbf{filiación} (la relación entre padres e hijos: paternidad y maternidad) y el \textbf{parentesco}.
\end{definicion}

\subsection{Antecedentes en el derecho romano}

En Roma el estatus determinaba el plexo de derechos de una persona. Existían tres:

\begin{table}[H]
\centering
\begin{tabularx}{\textwidth}{B{0.22\textwidth} Y}
\toprule
\textbf{Estatus} & \textbf{Contenido} \\
\midrule
\textit{Status libertatis} & Definía si la persona era libre o esclava; el libre tenía derechos que no tenía el esclavo. \\
\textit{Status civitatis} & Estado en relación con la ciudad: ciudadano romano o extranjero; los ciudadanos tenían derechos que los extranjeros no tenían. \\
\textit{Status familiae} & Estado en relación con la familia; la cabeza era el \textit{paterfamilias}, con derechos que no tenían los demás miembros. \\
\bottomrule
\end{tabularx}
\caption{Los tres estatus del derecho romano}
\end{table}

El sistema vinculado al \textit{status libertatis} ---la distinción entre libres y esclavos--- ha sido abolido.

\subsection{El estado en la era de los derechos humanos}

Con el tiempo la división en estatus que determinaba más o menos derechos se desvaneció: en la era de los derechos humanos no pueden hacerse distinciones de ese tipo, pues todos los seres humanos son iguales en derechos y condiciones. Por eso el estado pierde relevancia para determinar derechos.

\begin{importante}
A lo sumo reaparecen formas asimilables a estados particulares, que apuntan a una mayor protección de los sectores vulnerables, como ocurre con la Convención sobre los Derechos de las Personas con Discapacidad: quien presenta una discapacidad tiene un conjunto mayor de protecciones y derechos, para favorecer su inclusión y su desarrollo. No es, sin embargo, el concepto tradicional de estado: es una noción distinta. La noción de estado civil se reduce a la posición de la persona en la familia y en la sociedad.
\end{importante}

\subsection{Importancia e incidencia del estado civil}

El estado civil incide en otros institutos y atributos:

\begin{itemize}
\item \textbf{Domicilio}: los incapaces, especialmente los hijos mientras dura la responsabilidad parental, tienen su domicilio en el de sus padres (capítulo~\ref{cap:domicilio}).
\item \textbf{Capacidad}: por ejemplo, los cónyuges no pueden contratar entre sí bajo el régimen de comunidad.
\item \textbf{Apellido de los hijos}: los matrimoniales llevan el de sus padres; si la filiación es extramatrimonial, el del progenitor que los reconoció, si tienen una sola filiación determinada (capítulo~\ref{cap:nombre}).
\item \textbf{Alimentos}: la determinación del estado incide en la obligación alimentaria entre parientes, especialmente en la responsabilidad parental.
\item Otras obligaciones, como la de informar ciertas situaciones cuando existe una relación de familia.
\end{itemize}

\begin{regla}[Art.~344 CCyC]
Está prohibida la condición que afecte de modo grave la libertad de las acciones de la persona, como la de elegir un domicilio o una religión, o decidir sobre su estado civil ---por ejemplo, no casarse, o casarse o no con determinada persona---.
\end{regla}

No se puede imponer como condición de un contrato o de una donación que alguien cambie o no cambie su estado civil, ni exigir como condición para donar una casa que nunca contraiga matrimonio: sería contrario a la libertad básica de las personas.

\section{Las acciones de estado}

\begin{definicion}[Acción de estado]
Posibilidad técnica de iniciar reclamos en sede judicial orientados a la \textbf{reclamación} o a la \textbf{impugnación} de un estado determinado. El caso típico es el de la \textbf{filiación}.
\end{definicion}

\paragraph{Acción de reclamación.} La persona busca ser emplazada en un estado determinado, como ocurre cuando no ha sido reconocida por uno de sus progenitores.

\begin{ejemplo}[Reclamación de filiación]
Quien no ha sido reconocido por uno de sus progenitores presenta pruebas y ofrece un estudio genético. La parte reclamada contesta y puede negar el vínculo; se genera un contradictorio en el que se ofrecen y controvierten pruebas, se realiza el estudio genético y una sentencia determina si la persona es o no hija.
\end{ejemplo}

\paragraph{Acción de impugnación.} Busca desplazar a una persona de un estado determinado. Cuando alguien siente que quien pasa por su padre o madre no es su verdadero progenitor, puede iniciar una acción de impugnación de la filiación; también la puede iniciar el propio interesado.

\begin{importante}
A veces una persona reconoce a un hijo creyendo hacerle un bien, lo cual no es así: es el \textbf{reconocimiento complaciente}. Es problemático porque no es la vía jurídica para establecer la filiación, que es muy importante para la identidad y hace al conocimiento de los orígenes.
\end{importante}

El Código regula distintas variantes que exceden esta presentación. Como en todo proceso, las acciones pueden prosperar o no; si el juez de primera instancia no hace lugar a las pruebas, siempre cabe apelar a segunda instancia, donde interviene un tribunal colegiado (al menos tres jueces en la Capital Federal). La prueba biológica y genética es especialmente decisiva en la filiación por naturaleza.

\begin{importante}
Si se parte de un principio de verdad biológica, la ley presupone que el vínculo filiatorio se establece por el origen genético. Pero la filiación es un instituto complejo, compuesto también de elementos estáticos, como la posesión de estado.
\end{importante}

\section{La posesión de estado}

La posesión de estado es análoga a la posesión de las cosas. En el CCyC la posesión es la situación de quien detenta una cosa con espíritu de dueño, con independencia de que tenga o no el título; el Código denomina a estas situaciones relaciones de poder.

\begin{ejemplo}[Posesión de cosas y prescripción adquisitiva]
El poseedor que durante veinte años posee un inmueble como si fuera el dueño puede iniciar una acción para ser tenido por dueño ante el Registro de la Propiedad Inmueble. El Código denomina a esta figura \textbf{prescripción adquisitiva}.
\end{ejemplo}

\begin{definicion}[Posesión de estado]
Es el goce de hecho de un estado determinado, con independencia de si la persona detenta o no el título correspondiente.
\end{definicion}

\begin{importante}
Título y posesión coinciden en la mayoría de los casos, pero son cosas distintas. En el estado de hijo, el \textbf{título} es estar anotado como hijo de una persona en el Registro Civil; la \textbf{posesión} es ser tratado como hijo, en los hechos, por quienes figuran como sus padres.
\end{importante}

\subsection{Elementos históricos: \textit{nomen}, \textit{tractatus} y \textit{fama}}

\begin{table}[H]
\centering
\begin{tabularx}{\textwidth}{B{0.20\textwidth} Y}
\toprule
\textbf{Elemento} & \textbf{Contenido} \\
\midrule
\textit{Nomen} & El nombre: ser llamado como hijo (en el caso de la filiación). \\
\textit{Tractatus} & El trato: ser tratado como hijo. \\
\textit{Fama} & Ser conocido por los demás como hijo. \\
\bottomrule
\end{tabularx}
\caption{Elementos romanos de la posesión de estado}
\end{table}

Hoy el elemento más importante es el \textbf{trato}.

\begin{ejemplo}[Cómo se acredita el trato como hijo]
La fotografía de un cumpleaños, el pago de la obra social, de la cuota del colegio, del club o de un profesor particular, o llevar a la persona de vacaciones con la fotografía correspondiente.
\end{ejemplo}

\subsection{Las cuatro aplicaciones en el CCyC}

\paragraph{Primera: posesión de estado matrimonial (art.~423).} Por regla, no acredita el matrimonio, que es un acto jurídico solemne con consentimiento expresado ante el oficial del Registro Civil. Quienes simplemente conviven no acceden por ello a los beneficios del matrimonio: el Código regula la unión convivencial como vinculación distinta. La excepción se presenta cuando hubo un problema de forma en el acta de celebración.

\begin{regla}[Art.~423 CCyC (según la explicación de las fuentes)]
La posesión de estado matrimonial no acredita el matrimonio, que es un acto jurídico solemne. Sin embargo, si existió posesión de estado y se presenta un acta de celebración, los defectos formales del acta no pueden alegarse contra su existencia.
\end{regla}

\begin{ejemplo}[Error formal en el acta]
Al fallecer un cónyuge, el sobreviviente se presenta al Registro Civil para iniciar la sucesión y se le informa que hay un error en el acta de un matrimonio celebrado veinticinco años atrás, por lo que sería nulo. Como hubo posesión del estado de cónyuge durante todo ese tiempo, el vicio formal queda saneado y el matrimonio se considera válidamente celebrado.
\end{ejemplo}

\paragraph{Segunda: reconocimiento de la filiación (art.~584).} Es la aplicación más típica.

\begin{regla}[Art.~584 CCyC]
La posesión de estado debidamente acreditada en juicio tiene el mismo valor que el reconocimiento, siempre que no sea desvirtuada por prueba en contrario sobre el nexo genético.
\end{regla}

Si quien reclama prueba que el progenitor le pagaba la escuela, lo trataba, convivía con él y lo llevaba de vacaciones, acredita la posesión de estado; pero la prueba genética prevalece sobre la posesión de estado cuando entran en conflicto.

\paragraph{Tercera: límite a la negación de la filiación presumida (art.~591).} Si una persona está casada se presume que los hijos nacidos durante el matrimonio son de ambos cónyuges. El cónyuge que no dio a luz puede negar esa filiación en un plazo muy inmediato al nacimiento ---el caso habitual es alegar una infidelidad---, pero no puede negarla si hubo posesión de estado, es decir, si trató al hijo como propio.

\paragraph{Cuarta: adopción de personas mayores de edad (art.~597).} En principio se adopta a menores, pero el Código permite adoptar excepcionalmente a un mayor cuando hubo posesión de estado de hijo durante su minoridad. Puede tener un efecto sucesorio importante, porque el hijo adoptivo podría suceder a quien lo trató como padre o madre.

\begin{table}[H]
\centering
\begin{tabularx}{\textwidth}{B{0.14\textwidth} B{0.38\textwidth} Y}
\toprule
\textbf{Artículo} & \textbf{Materia} & \textbf{Efecto de la posesión de estado} \\
\midrule
423 & Matrimonio & No acredita el matrimonio, pero sanea los defectos formales del acta. \\
584 & Reconocimiento de la filiación & Vale como reconocimiento, salvo prueba genética en contrario. \\
591 & Negación de la filiación presumida & Impide negar la filiación si se trató al hijo como propio. \\
597 & Adopción de mayores de edad & Habilita la adopción excepcional del mayor tratado como hijo en su minoridad. \\
\bottomrule
\end{tabularx}
\caption{Aplicaciones de la posesión de estado en el CCyC}
\end{table}

\section{La prueba del estado civil}

\subsection{Las partidas}

\begin{regla}[Regla general sobre la prueba del estado civil]
El nacimiento, sus circunstancias de tiempo y lugar, el sexo, el nombre y la filiación de las personas nacidas se prueban con las partidas del Registro Civil.
\end{regla}

El Registro Civil se denomina técnicamente Registro del Estado Civil y Capacidad de las Personas. La forma de registración está regida por una ley nacional ---la Ley 26.413, que uniforma la registración de los hechos y actos relativos al estado civil en todo el país---, pero la organización de cada registro depende de cada provincia, por la organización federal.

La exigencia de registración se funda en que el estado civil está constituido a partir de hechos: el nacimiento, su tiempo y lugar y el sexo al nacer son hechos; el nombre se adjudica y la filiación se determina jurídicamente. Existe un interés público sustantivo en que la sociedad registre estos eventos. Los grandes hechos y actos a registrar son la adopción, el matrimonio, el divorcio y el fallecimiento, que cierra el ciclo; de allí los tres grandes tipos de partida: de nacimiento, de matrimonio y de defunción.

\begin{definicion}[Partida]
Asiento en el libro del Registro del Estado Civil y Capacidad de las Personas en el que constan los hechos y actos constitutivos del estado: nacimiento, matrimonio y defunción.
\end{definicion}

\begin{importante}
Las partidas son \textbf{instrumentos públicos}: tienen valor probatorio propio y hacen plena fe de su contenido.
\end{importante}

\begin{nota}
La partida no se confunde con el certificado médico de nacimiento o de defunción: este lo expide el médico y sirve para concurrir al Registro Civil, donde se confecciona la partida. Lo que siempre debe ofrecerse como prueba es la partida; solo si no existiera corresponde recurrir a otros medios.
\end{nota}

\subsection{La prueba supletoria}

La ausencia de partidas puede deberse a que no hay registro, a que falta la partida o el asiento, o a que este es nulo.

\begin{ejemplo}[El incendio de un Registro Civil]
Si un incendio destruyera los libros de un período, la persona puede acompañar otros certificados ---por ejemplo, de un hecho religioso como el bautismo, o una constancia de inscripción escolar que consigne la fecha de nacimiento--- para acreditar su fecha de nacimiento, o la de defunción de otra persona, a los fines de elaborar la partida.
\end{ejemplo}

\subsection{Fallecimiento sin hallazgo o identificación del cadáver (art.~98)}

El artículo 98 contempla, en su segunda parte, el caso en que fallece una persona pero no se encuentra el cadáver, o no puede identificarse, aunque no haya duda de que murió.

\begin{regla}[Art.~98 CCyC]
Si el cadáver de una persona no es hallado, o no puede ser identificado, el juez puede tener por comprobada la muerte y disponer la pertinente inscripción en el registro, siempre que la desaparición se hubiera producido en circunstancias tales que la muerte deba ser tenida como cierta.
\end{regla}

\begin{ejemplo}[El alpinista]
Un alpinista cae por un precipicio en un lugar inaccesible ante la vista de cinco compañeros; nunca se encontrará el cadáver. No puede declarárselo fallecido directamente, porque se necesita un médico que certifique la muerte, ni corresponde el proceso de presunción de fallecimiento, mucho más largo: se sigue un procedimiento breve ante un juez para disponer la inscripción de la muerte.
\end{ejemplo}

\begin{ejemplo}[El accidente aéreo]
Fallecen todos los ocupantes de una aeronave y los cadáveres pueden no identificarse individualmente. Si se prueba que una persona estaba en ese vuelo, se la puede anotar como fallecida y iniciar los trámites sucesorios.
\end{ejemplo}

No es un fallecimiento presunto, sino indudable: las circunstancias deben hacer tener la muerte por cierta. Acreditado el fallecimiento, se expide una partida de defunción común.

\cornell{El estado civil}{
\cq{¿Qué es el estado civil?}{La posición de la persona en la sociedad y en la familia: estado matrimonial, filiación y parentesco.}
\cq{¿Por qué pierde relevancia en la era de los derechos humanos?}{Porque ya no se admiten estatus que determinen más o menos derechos (a diferencia de Roma: \textit{status libertatis}, \textit{civitatis} y \textit{familiae}); la protección de vulnerables no es el estado tradicional.}
\cq{¿Qué condiciones están prohibidas?}{Las que afectan gravemente la libertad de decidir sobre el estado civil (art.~344): casarse o no, o hacerlo con determinada persona.}
\cq{¿Qué son las acciones de estado?}{Acciones judiciales de reclamación (emplazar en un estado) e impugnación (desplazar de un estado), típicas en la filiación.}
\cq{¿Qué es la posesión de estado y cuáles son sus cuatro aplicaciones?}{El goce de hecho de un estado con independencia del título; elementos \textit{nomen}, \textit{tractatus} (el más importante) y \textit{fama}. Arts.~423 (matrimonio), 584 (reconocimiento), 591 (límite a la negación) y 597 (adopción de mayores).}
\cq{¿Cómo se prueba el estado civil?}{Con las partidas del Registro Civil (instrumentos públicos de plena fe: nacimiento, matrimonio, defunción); supletoriamente por otros medios; art.~98 para la muerte cierta sin cadáver.}
}{El estado civil es la posición de la persona en la familia y la sociedad. Se discute mediante acciones de reclamación e impugnación, la posesión de estado tiene cuatro aplicaciones en el CCyC y se prueba con las partidas del Registro Civil, que son instrumentos públicos.}

% ============================================================
\part{La capacidad de la persona humana}
% ============================================================

\chapter{La capacidad jurídica}
\label{cap:capacidad}

La capacidad es el atributo que permite determinar qué aptitud tiene la persona para entablar relaciones jurídicas. Es el punto de partida del derecho civil porque define quién puede ser titular de derechos y quién puede ejercerlos por sí mismo, con incidencia en los contratos, las sucesiones, la representación judicial y la protección de menores y personas con discapacidad.

Un caso guía atraviesa el capítulo: una niña de corta edad hereda un inmueble de su padre fallecido. Es titular del derecho sobre el bien (capacidad de derecho), pero no puede ejercerlo por sí misma (carece de capacidad de ejercicio), por lo que la representan sus padres o un tutor. Si, ya adulta, un juez determinara que necesita apoyo para administrar bienes, se le designaría un asistente y no un representante. En todo el recorrido, el Ministerio Público vela por sus intereses.

\section{Concepto de capacidad jurídica}

La capacidad es una aptitud para entablar relaciones jurídicas que presenta dos componentes: uno de \textbf{derecho} (aptitud para ser titular de derechos y obligaciones) y otro \textbf{de hecho o de ejercicio} (aptitud para ejercer esos derechos). El Código los regula bajo el concepto común de \textbf{capacidad jurídica}.

\begin{regla}[Art.~22 CCyC: regla general de la capacidad]
Toda persona humana goza de la aptitud para ser titular de derechos y deberes jurídicos. La ley puede privar o limitar esta capacidad respecto de hechos, simples actos o actos jurídicos determinados.
\end{regla}

La razón de fondo es existencial: si los seres humanos se desarrollan entablando relaciones jurídicas, negar esa aptitud equivaldría a impedirles satisfacer sus necesidades básicas y desplegar su personalidad. Ese despliegue requiere ser titular de derechos y deberes, comenzando por los derechos humanos, que no son concedidos sino reconocidos: libertad, libertad de contrato, libertad de asociación y otras libertades básicas.

\subsection{Capacidad de derecho}

Algunos autores identifican la capacidad de derecho con la condición de persona; un ser humano con incapacidad absoluta de derecho sería un muerto civil. Por eso se reconoce a toda persona; solo pueden existir limitaciones o privaciones ---las llamadas incapacidades de derecho relativas--- para ciertas circunstancias.

\begin{importante}
Las incapacidades de derecho son siempre \textbf{relativas}, se aplican a \textbf{actos puntuales} y se establecen para proteger el \textbf{interés general}. No pueden suplirse por medio de un representante y su violación genera \textbf{nulidad absoluta}.
\end{importante}

\subsection{Capacidad de ejercicio}

Es la aptitud para desplegar y gozar los derechos por sí mismo; la clave es la actuación por sí. A diferencia de las limitaciones a la capacidad de derecho, las de ejercicio se establecen para \textbf{proteger a la propia persona}.

\begin{ejemplo}[Menor de edad]
Un menor es incapaz de ejercicio no porque carezca de derechos, sino porque, por su inmadurez, se le proveen representantes y salvaguardias que lo tutelan en su vulnerabilidad.
\end{ejemplo}

\section{Incapacidades de derecho: el artículo 1002}

El ejemplo más ilustrativo son las inhabilidades para contratar.

\begin{regla}[Art.~1002 CCyC: inhabilidades para contratar]
No pueden contratar en interés propio: los funcionarios públicos respecto de bienes cuya administración o enajenación está a su cargo; los jueces, funcionarios y auxiliares de la justicia, árbitros, mediadores y sus auxiliares, respecto de bienes relacionados con procesos en los que intervienen o han intervenido; los abogados y procuradores respecto de bienes litigiosos en procesos en que intervienen o han intervenido; los cónyuges bajo el régimen de comunidad entre sí; los albaceas que no son herederos respecto de bienes de las testamentarías a su cargo.
\end{regla}

La razón es de interés público: no se admite que el funcionario que administra un bien lo compre, que un árbitro adquiera los bienes en disputa o que un juez compre la casa en remate en su propio juzgado, pues habría conflicto de intereses.

\begin{table}[H]
\centering
\begin{tabularx}{\textwidth}{B{0.26\textwidth} Y Y}
\toprule
\textbf{Característica} & \textbf{Incapacidad de derecho} & \textbf{Incapacidad de ejercicio} \\
\midrule
¿A quién protege? & Al interés general o público & A la propia persona \\
¿Puede ser absoluta? & Nunca (siempre relativa) & Sí, en casos extremos (art.~32) \\
¿Se suple con representante? & No & Sí \\
Alcance & Actos puntuales determinados & Puede ser amplio \\
Regulación & A lo largo de todo el Código & Art.~24 y siguientes \\
Sanción por violación & Nulidad absoluta & Acto anulable (según el caso) \\
\bottomrule
\end{tabularx}
\caption{Incapacidad de derecho e incapacidad de ejercicio}
\end{table}

\section{Incapaces de ejercicio}

\begin{regla}[Art.~24 CCyC]
Son incapaces de ejercicio: a) la persona por nacer; b) la persona que no cuenta con la edad y grado de madurez suficiente, con el alcance dispuesto en la Sección 2.ª de este Capítulo; c) la persona declarada incapaz por sentencia judicial, en la extensión dispuesta en esa decisión.
\end{regla}

\paragraph{Personas por nacer.} Conforme al artículo 19, están concebidas hasta el nacimiento. Son personas, las representan sus padres (art.~101) y pueden recibir bienes, pero la adquisición de derechos patrimoniales queda sujeta al nacimiento con vida (art.~21), como se vio en el capítulo~\ref{cap:por-nacer}.

\paragraph{Personas menores de edad.} Definida en el artículo 25 como quien no cuenta con la edad y grado de madurez suficiente, es por regla incapaz de ejercicio, pero adquiere una capacidad progresiva a medida que crece.

\begin{importante}
Desde que nace hasta los 18 años ---mayoría de edad--- la persona evoluciona. La representación que corresponde primero a sus padres (responsabilidad parental) decrece a medida que crecen sus aptitudes, hasta que se hace plenamente capaz.
\end{importante}

El Estado reconoce a los padres autoridad para educar porque los considera en las mejores condiciones para criar a sus hijos; en casos muy graves pueden ser privados de la responsabilidad parental, pero la regla es que ellos representan al hijo.

\paragraph{Personas declaradas incapaces por sentencia (art.~32, último párrafo).} Son casos excepcionales en que la persona está absolutamente imposibilitada de interaccionar con su entorno y no puede expresar su voluntad de ningún modo, y el sistema de apoyo resulta ineficaz.

\begin{regla}[Art.~32, último párrafo, CCyC]
Por excepción, cuando la persona se encuentre absolutamente imposibilitada de interaccionar con su entorno y expresar su voluntad por cualquier modo, medio o formato adecuado, y el sistema de apoyos resulte ineficaz, el juez puede declarar la incapacidad y designar un curador.
\end{regla}

El juez debe determinar en qué actos y funciones puede actuar la persona y se nombra un curador; son casos muy extremos, generalmente de discapacidad severa.

\begin{ejemplo}[Caso de Juana Fernández]
Juana Fernández tiene 3 años; su padre fallece y es única heredera de un inmueble y dinero en una cuenta bancaria. Las decisiones sobre esos bienes no las tomará una niña de tres años: se nombra un tutor o representante. Los bienes están a su nombre (capacidad de derecho), pero no puede ejercer ese derecho por sí misma (capacidad de ejercicio); el tutor interviene en esto último.
\end{ejemplo}

\section{Capacidad restringida y sistemas de apoyo}

La persona con capacidad restringida es capaz de ejercicio, pero una sentencia judicial establece que ciertos actos no puede realizarlos sola y necesita apoyo. Es hoy prácticamente la regla del sistema, especialmente a partir de la Convención sobre los Derechos de las Personas con Discapacidad, según la cual deben integrarse a la sociedad y pueden expresar sus preferencias aun necesitando apoyos.

\begin{ejemplo}[Restricción para actos patrimoniales]
Una persona con discapacidad mental puede manejarse con pequeñas sumas pero no comprende el valor de un inmueble. El juez restringe su capacidad: en la vida cotidiana actúa libremente, pero para vender inmuebles necesita el apoyo de un familiar.
\end{ejemplo}

El apoyo no ocupa el lugar de la persona ni ejerce los derechos en su nombre: es un caso de \textbf{asistencia}, en el que concurren ambas voluntades, con las funciones que determina el juez. Si el apoyo advierte que se hace un mal negocio ---por ejemplo, vender muy barato en un momento de caída de precios--- puede negarse: existe para dar seguridad, aconsejar y evitar que se aproveche a quien no comprende bien el valor del dinero.

\begin{table}[H]
\centering
\begin{tabularx}{\textwidth}{B{0.22\textwidth} Y Y}
\toprule
\textbf{Aspecto} & \textbf{Representación} & \textbf{Asistencia (apoyos)} \\
\midrule
¿Para quién? & Incapaces de ejercicio & Personas con capacidad restringida \\
¿Quién actúa? & El representante, en nombre del incapaz & Ambos: persona y apoyo \\
Voluntad del titular & No es necesaria & Es necesaria y principal \\
¿Quién designa? & La ley o el juez & El juez (proceso específico) \\
Finalidad & Sustituir al incapaz & Acompañar y proteger \\
\bottomrule
\end{tabularx}
\caption{Representación y asistencia}
\end{table}

\subsection{Personas inhabilitadas (pródigos)}

\begin{regla}[Art.~48 CCyC: inhabilitación por prodigalidad]
Pueden ser inhabilitados quienes, por la prodigalidad en la gestión de sus bienes, expongan a su cónyuge, conviviente o a sus hijos menores de edad o con discapacidad a la pérdida del patrimonio. Se considera que hay prodigalidad cuando, sin mediar razones atendibles, se dilapidan los bienes.
\end{regla}

Iniciada la causa, se restringe la capacidad de la persona y se le nombran apoyos para la administración de sus bienes.

\subsection{Representantes legales (art.~101)}

\begin{regla}[Art.~101 CCyC]
Son representantes: a) de las personas por nacer, sus padres; b) de las personas menores de edad no emancipadas, sus padres y, si faltan, son incapaces o fueron privados de la responsabilidad parental, un tutor; c) de las personas con capacidad restringida y de las declaradas incapaces, los apoyos (cuando tengan representación) y el curador, respectivamente.
\end{regla}

\section{El Ministerio Público (art.~103)}

Institución de fundamento constitucional, extrapoder, que en materia civil vela por los intereses de menores, incapaces y personas con capacidad restringida. Se divide en Ministerio Público de la Defensa y Ministerio Público Fiscal; aquí interesa el primero. Es una protección complementaria: a los incapaces se les provee representante o asistente, pero además el Ministerio Público vela siempre por sus intereses.

\begin{regla}[Art.~103 CCyC]
La actuación del Ministerio Público respecto de personas menores de edad, incapaces y con capacidad restringida, y de aquellas cuyo ejercicio de capacidad requiera de un sistema de apoyos, puede ser, en el ámbito judicial, complementaria o principal.
\end{regla}

\begin{table}[H]
\centering
\begin{tabularx}{\textwidth}{B{0.20\textwidth} Y Y}
\toprule
\textbf{Actuación} & \textbf{Cuándo} & \textbf{Ejemplo} \\
\midrule
Judicial complementaria & En todos los procesos con menores, incapaces o personas con capacidad restringida; se les da vista de las actuaciones, sin sustituir al representante. & En una sucesión con un heredero menor, de cada paso que toque sus intereses se da vista al asesor o defensor de menores e incapaces; la madre sobreviviente sigue siendo representante principal. \\
Judicial principal & Cuando los intereses están comprometidos y los representantes no accionan, para exigir a los representantes que cumplan sus deberes o cuando se carece de representante legal. & Un menor con un solo progenitor que no pasa alimentos: el Ministerio Público promueve el juicio. Una niña huérfana: actúa como si fuera el representante, sustituyendo a los del art.~101. \\
Extrajudicial & Extraordinaria; sin proceso, para preservar derechos económicos, sociales y culturales ante la ausencia, carencia o inacción de los representantes. & Un niño expuesto a un ambiente de drogas y peligroso, con padres inactivos o promotores de la situación. \\
\bottomrule
\end{tabularx}
\caption{Formas de actuación del Ministerio Público}
\end{table}

% ------------------------------------------------------------
\chapter{El régimen de capacidad de las personas menores de edad}
\label{cap:menores}

Este capítulo estudia el régimen de capacidad de las personas menores de edad en el CCyC: los principios rectores que lo estructuran, los actos que la ley autoriza según la edad y las consecuencias y modos de finalización de la incapacidad de ejercicio. Es un contenido de aplicación directa en las áreas de familia, niñez y adolescencia, salud, contratos laborales y responsabilidad civil: determina qué puede hacer un menor por sí mismo, cuándo intervienen sus representantes y cuáles son los límites de esa representación.

\begin{ejemplo}[El recorrido de la capacidad progresiva]
Los temas del capítulo pueden entenderse como un recorrido de menor a mayor autonomía. Juan nació sin capacidad de ejercicio y fue creciendo: a los 13 años pudo decidir sobre un tratamiento médico sencillo; a los 16 pudo trabajar, firmar contratos y cuidar su propio cuerpo como un adulto; al contraer matrimonio antes de los 18 años quedó emancipado, aunque con ciertas restricciones para proteger su patrimonio. Así se articula la tensión entre \textbf{autonomía y protección} que atraviesa todo el capítulo.
\end{ejemplo}

\section{Régimen general}

\subsection{Menor de edad y adolescente}

\begin{regla}[Art.~25 CCyC]
Menor de edad es la persona que no ha cumplido 18 años. Si hubiera cumplido 13 años, se la denomina adolescente.
\end{regla}

La categoría de adolescente ---el menor de entre 13 y 18 años--- ha sido criticada por no constituir un concepto plenamente técnico, aunque está consagrada en el Código.

\subsection{La incapacidad de ejercicio como medida de protección}

El punto de partida es una \textbf{incapacidad de ejercicio}. No es una disminución de la persona: el menor es un sujeto con dignidad, titular de derechos y con aptitud para adquirirlos. Toda persona nace en una familia, y por eso el ordenamiento confía primero en los padres como representantes: ofrecen el ámbito de educación, crianza y contención necesario para el desarrollo.

\begin{importante}
La incapacidad de ejercicio \textbf{no es una sanción} ni una disminución de la persona: es una herramienta jurídica que \textbf{protege} al menor, supliendo su falta de madurez con la actuación de sus representantes legales.
\end{importante}

\begin{ejemplo}[El niño de diez años que hereda un inmueble]
No sabe el valor de las cosas ni cómo negociar un contrato; puede hacer un negocio que lo perjudique o ser abusado por personas inescrupulosas. Por eso se le designa un representante que actúe en su nombre.
\end{ejemplo}

El derecho contempla a la vez el \textbf{derecho a ser oído}: la relación entre padres e hijos presupone una manera de ir educando y llevando hacia la madurez a una persona necesariamente inmadura.

\subsection{La mayoría de edad a través del tiempo}

\begin{table}[H]
\centering
\begin{tabularx}{\textwidth}{Y c Y}
\toprule
\textbf{Período} & \textbf{Edad de mayoría} & \textbf{Norma} \\
\midrule
Código Civil histórico & 22 años & Código de Vélez Sarsfield \\
Reforma intermedia & 21 años & Modificación legislativa \\
Ley 26.579 (2009) & 18 años & Ley 26.579 \\
Código Civil y Comercial (2015) & 18 años & Ratificado por el CCyC \\
\bottomrule
\end{tabularx}
\caption{Evolución de la mayoría de edad en la Argentina}
\end{table}
% TODO: verificar consistencia entre fuentes originales: la fuente atribuye 22 años al Código histórico de Vélez y 21 años a una «reforma intermedia» sin identificarla.

Fijar la mayoría en 18 años mantiene coherencia con la Convención sobre los Derechos del Niño, que define como niño a todo ser humano hasta esa edad. Se advierte contra reducir la mayoría de edad con espíritu demagógico: bajarla implica también dejar de brindar alimentos al joven; por eso la ley dispuso que, aunque se sea mayor a los 18, en algunos casos debe continuar la provisión de alimentos mientras la persona se forma y no tiene trabajo independiente.

\begin{nota}
El régimen alimentario posterior a la mayoría extiende la obligación hasta los 21 años con carácter general, y hasta los 25 si el joven se encuentra en proceso de formación educativa o universitaria.
\end{nota}

\subsection{¿Regla de capacidad o de incapacidad?}

Existe debate. Un sector sostiene que, por aplicación de las convenciones internacionales, el principio sería la capacidad; otra parte relevante de la doctrina argentina entiende que la regla es la \textbf{incapacidad}, que va admitiendo una capacidad progresiva.

\begin{regla}[Art.~26 CCyC, primera parte]
La persona menor de edad ejerce sus derechos a través de sus representantes legales. No obstante, la que cuenta con edad y grado de madurez suficiente puede ejercer por sí los actos que le son permitidos por el ordenamiento jurídico.
\end{regla}

\section{Principios rectores (arts.~26 y 639)}

El artículo 639, ubicado en la responsabilidad parental (Libro Segundo), enuncia tres principios:

\begin{description}[leftmargin=0pt,style=unboxed]
\item[Interés superior del niño.] Tiene como fuente la Convención sobre los Derechos del Niño de 1989, ratificada por la Ley 23.849 e incorporada con jerarquía constitucional por la reforma de 1994 (art.~75 inc.~22 CN). La palabra «interés» remite a un valor objetivo: qué es lo mejor para ese niño.
\item[Autonomía progresiva.] A medida que el niño crece en madurez y edad, disminuye la representación de los padres y aumenta su capacidad de ejercer por sí sus derechos. Se empalma con el artículo 26: la autonomía va acompañada de autorizaciones legales que el ordenamiento otorga progresivamente.
\item[Derecho a ser oído.] Proviene del artículo 12 de la Convención y tiene su correlato en el artículo 26: ser oído en todos los asuntos de su vida y participar en las decisiones fundamentales sobre la conducción de su persona. Por ejemplo, en un divorcio con hijos menores se los escucha sobre dónde vivirán o a qué colegio asistirán.
\end{description}

\subsection{Representación de los menores}

\begin{table}[H]
\centering
\begin{tabularx}{\textwidth}{B{0.24\textwidth} B{0.28\textwidth} Y}
\toprule
\textbf{Representante} & \textbf{Fuente legal} & \textbf{Carácter} \\
\midrule
Padres / progenitores & Art.~101 CCyC; responsabilidad parental & Principal y primario \\
Tutores & Art.~101 CCyC & Supletorio (a falta de padres, por fallecimiento, privación de la responsabilidad parental u otra causa) \\
Abogado del niño & Ley 26.061 & Técnico, en juicio: asesora al niño ante un conflicto de intereses que sus progenitores no pueden representar \\
Ministerio Público & Art.~103 CCyC & Genérico: judicial o extrajudicial, principal o complementario \\
\bottomrule
\end{tabularx}
\caption{Sistema de representación del menor de edad}
\end{table}

\subsection{Grado de madurez suficiente}

A veces la ley fija un límite de edad (por ejemplo, 10, 13 o 16 años); otras usa la fórmula amplia «grado de madurez suficiente», criticada por parte de la doctrina. Cuando la cuestión llega a un juez, este ---con el asesoramiento de su equipo interdisciplinario--- determina si el menor tiene la madurez exigida. Rivero critica esa indeterminación y sostiene que debería sustituirse por edades fijas; según las notas originales, existe un proyecto de reforma del CCyC elaborado por una comisión presidida por dicho autor que tiende a eliminarla.

\section{Actos que la ley autoriza realizar según la edad}

\subsection{Tratamientos médicos (art.~26)}

\begin{regla}[Art.~26 CCyC: tratamientos médicos]
Desde los 13 años se presume que la persona menor de edad posee aptitud para decidir sobre tratamientos no invasivos o que no comprometan su vida o su salud. Si el tratamiento es invasivo y compromete su salud, se requiere la asistencia de la persona menor de edad y de sus padres; en caso de discordancia se decide con base en la opinión médica, teniendo en cuenta el interés superior. Desde los 16 años el adolescente es considerado como un adulto para el cuidado de su propio cuerpo.
\end{regla}

\begin{table}[H]
\centering
\begin{tabularx}{\textwidth}{B{0.20\textwidth} Y Y}
\toprule
\textbf{Edad} & \textbf{Tipo de tratamiento} & \textbf{¿Quién decide?} \\
\midrule
0 a 12 años & Cualquier tratamiento & Los padres (o tutores) \\
13 a 15 años & No invasivo / sin riesgo de vida & El propio menor (presunción de capacidad) \\
13 a 15 años & Invasivo / con riesgo de vida & Menor y padres (asistencia conjunta) \\
16 años en adelante & Cuidado del propio cuerpo & El adolescente, como adulto \\
\bottomrule
\end{tabularx}
\caption{Decisiones médicas según la edad}
\end{table}

\subsection{Resumen de actos permitidos}

\paragraph{Desde temprana edad.} Pequeños contratos de menor cuantía (art.~684): quien compra algo en el kiosco o en la librería celebra un contrato que no es nulo.

\paragraph{Desde los 13 años.}

\begin{itemize}
\item Decidir sobre tratamientos no invasivos (art.~26).
\item Iniciar una acción para conocer sus orígenes (art.~596).
\item Reconocer hijos.
\item Ejercer la responsabilidad parental, con la supervisión del artículo 644: los abuelos pueden intervenir en las cuestiones más trascendentes ante la inacción de los progenitores adolescentes y vetar decisiones contrarias al interés del niño.
\item Participar en forma autónoma en juicio con autorización de sus progenitores (art.~677).
\end{itemize}

\paragraph{Desde los 16 años.}

\begin{itemize}
\item Ser equiparado a los adultos en el cuidado del propio cuerpo (art.~26, \textit{in fine}).
\item Celebrar contrato de trabajo: la Ley 26.390, en clave con el Convenio de la OIT, prohíbe el trabajo de menores de 16 años en sus distintas formas; desde los 16 se presume que pueden celebrarlo con autorización de sus padres, o solos si viven de forma independiente. % TODO: verificar consistencia entre fuentes originales: la fuente dice «solos si viven de forma independiente» y, más adelante, «con la autorización de sus padres, o bien solo si vive de forma independiente».
\item Ser donante de sangre con autorización de sus padres.
\end{itemize}

\subsection{El título habilitante (art.~30)}

\begin{regla}[Art.~30 CCyC]
La persona menor de edad que obtiene un título habilitante de profesión puede ejercerla por cuenta propia, celebrar todos los actos que la profesión le exija y administrar y disponer libremente de los bienes que adquiera con el producto de su trabajo.
\end{regla}

Quien obtuvo a los 16 años el título de profesor de piano puede ejercer y administrar lo que gane. La doctrina discute el caso del menor de 16 años con título: un sector entiende que el artículo no distingue edad; otro, que quien pudo obtener el título ---con validación del Ministerio de Educación y requisitos formales--- ya demostró madurez suficiente para ejercer la profesión.

\subsection{El matrimonio de personas menores de edad}

\begin{table}[H]
\centering
\begin{tabularx}{\textwidth}{B{0.20\textwidth} B{0.24\textwidth} Y}
\toprule
\textbf{Edad} & \textbf{¿Puede casarse?} & \textbf{Requisito} \\
\midrule
18 años o más & Sí & Ninguno adicional (regla general) \\
16 y 17 años & Sí & Autorización de los representantes legales; si la niegan, dispensa judicial \\
Menor de 16 años & Eventualmente & Siempre con dispensa judicial \\
\bottomrule
\end{tabularx}
\caption{Régimen matrimonial según la edad}
\end{table}

\section{Consecuencias jurídicas y fin de la incapacidad}

\subsection{Nulidad relativa de los actos celebrados sin representación}

Si no se respeta el régimen y el menor actúa sin representación, el acto es de nulidad relativa, porque está puesta para proteger al menor; si el acto lo beneficia, los padres pueden confirmarlo; y solo el menor o sus representantes pueden alegarla, no la contraparte.

\begin{ejemplo}[El cuadro heredado]
Una persona de 13 años heredó un cuadro valuado en 2.000 dólares y lo vendió en una plataforma en línea por 4.000. Sus padres, en lugar de anular el acto, lo confirmaron porque resultó un buen negocio.
\end{ejemplo}

\subsection{Responsabilidad civil por daños causados por menores}

\begin{table}[H]
\centering
\begin{tabularx}{\textwidth}{Y c c}
\toprule
\textbf{Situación} & \textbf{Responde el menor} & \textbf{Responden los padres} \\
\midrule
Menor de 10 años, acto ilícito & No (acto involuntario) & Sí, siempre (objetiva) \\
Mayor de 10 años, acto ilícito & Sí, con su patrimonio & Sí, también (objetiva) \\
\bottomrule
\end{tabularx}
\caption{Responsabilidad por daños según la edad del menor}
\end{table}

Los padres responden siempre, con responsabilidad objetiva fundada en el deber de vigilancia, aunque el menor de más de 10 años responda también con su propio patrimonio.

\subsection{Fin de la incapacidad: la mayoría de edad}

\begin{importante}
El día en que la persona cumple 18 años es automáticamente mayor de edad y plenamente capaz. No se requiere acto formal; basta acreditar la edad. Rige en el campo del derecho civil (contratos, familia, patrimonio) y no debe confundirse con la responsabilidad penal ni con el sufragio.
\end{importante}

\subsection{Fin de la incapacidad: la emancipación por matrimonio}

La emancipación confiere anticipadamente la plena capacidad a un menor. El CCyC mantiene un único tipo. Históricamente existían la emancipación por habilitación de edad ---entre los 18 y los 21 años, cuando la mayoría era a los 21--- y la emancipación por matrimonio; al bajar la mayoría a los 18, la primera desapareció y queda solo la matrimonial.

\begin{regla}[Art.~27 CCyC: emancipación]
La celebración del matrimonio antes de los 18 años emancipa a la persona menor de edad. La persona emancipada goza de plena capacidad de ejercicio con las limitaciones previstas en este Código. La emancipación es irrevocable. La nulidad del matrimonio no deja sin efecto la emancipación, excepto respecto del cónyuge de mala fe para quien el matrimonio fue anulado.
\end{regla}

\subsection{Restricciones a la capacidad del emancipado (arts.~28 y 29)}

\begin{regla}[Art.~28 CCyC: actos prohibidos, ni con autorización judicial]
Las personas emancipadas no pueden: aprobar las cuentas de sus tutores ni darles finiquito; hacer donación de bienes que hubieran recibido a título gratuito; ser fiadores.
\end{regla}

Fundamentos: el menor que tuvo tutor no puede aprobar sus cuentas hasta ser mayor (regla del art.~130); la prohibición de donar lo heredado proviene del Código de Vélez y evita que disponga de bienes cuyo valor no comprende cabalmente; la de ser fiador busca custodiar su patrimonio.

\begin{regla}[Art.~29 CCyC: disposición de bienes recibidos a título gratuito]
El emancipado puede disponer de ellos a título oneroso (venderlos, no donarlos), pero requiere autorización judicial, que solo procede en caso de absoluta necesidad o ventaja evidente.
\end{regla}

\begin{ejemplo}[El inmueble heredado por el emancipado]
No puede venderlo libremente: solo con autorización judicial y si existe absoluta necesidad (como una operación médica urgente) o ventaja evidente (un gran negocio). El CCyC unificó el sistema eliminando la excepción del Código de Vélez que admitía prescindir de la dispensa cuando uno de los cónyuges era mayor de edad.
\end{ejemplo}

\begin{table}[H]
\centering
\begin{tabularx}{\textwidth}{Y B{0.14\textwidth} B{0.16\textwidth} Y}
\toprule
\textbf{Acto} & \textbf{Artículo} & \textbf{¿Puede?} & \textbf{Condición} \\
\midrule
Aprobar cuentas del tutor & 28 & No & Nunca, ni con autorización judicial \\
Donar bienes recibidos gratuitamente & 28 & No & Nunca, ni con autorización judicial \\
Ser fiador & 28 & No & Nunca, ni con autorización judicial \\
Vender bienes recibidos gratuitamente & 29 & Sí, con límites & Autorización judicial: necesidad absoluta o ventaja evidente \\
\bottomrule
\end{tabularx}
\caption{Restricciones del emancipado}
\end{table}

\subsection{Mayoría de edad anticipada para egresados de institutos (Ley 27.364)}

Estableció una mayoría de edad anticipada para adolescentes que egresan de institutos y no cuentan con cuidados parentales. La doctrina la criticó: por tratarse de adolescentes sin cuidados parentales y provenientes de la institucionalización, podría funcionar como forma de abandono, aunque la ley contempla un programa estatal de protección.

\cornell{La capacidad jurídica y el régimen de los menores}{
\cq{¿Qué distingue la capacidad de derecho de la de ejercicio?}{La de derecho es la aptitud para ser titular de derechos y deberes (se reconoce a toda persona; sus limitaciones son siempre relativas, protegen el interés general y su violación causa nulidad absoluta). La de ejercicio es la aptitud para ejercerlos por sí; sus limitaciones protegen a la propia persona y se suplen con representante.}
\cq{¿Quiénes son incapaces de ejercicio (art.~24)?}{La persona por nacer, la que no tiene edad y madurez suficiente y la declarada incapaz por sentencia (art.~32, última parte, solo en casos extremos con curador).}
\cq{¿Qué diferencia hay entre representación y asistencia?}{En la representación el representante actúa en nombre del incapaz; en la asistencia, la persona con capacidad restringida actúa junto con su apoyo y su voluntad es necesaria.}
\cq{¿Cómo actúa el Ministerio Público (art.~103)?}{De forma judicial complementaria (vista), judicial principal (cuando falta o no actúa el representante) y extrajudicial (extraordinaria).}
\cq{¿Qué principios rigen la capacidad de los menores?}{Interés superior del niño, autonomía progresiva y derecho a ser oído (arts.~26 y 639).}
\cq{¿Qué puede hacer un adolescente según su edad?}{Desde los 13: tratamientos no invasivos, acción de orígenes, reconocer hijos; desde los 16: cuidado del propio cuerpo como adulto, trabajar, título habilitante (art.~30); matrimonio con autorización o dispensa.}
\cq{¿Cómo termina la incapacidad y qué límites tiene el emancipado?}{Por mayoría a los 18 años o emancipación por matrimonio (art.~27). El emancipado no puede aprobar cuentas del tutor, donar lo recibido gratuitamente ni ser fiador (art.~28) y necesita autorización judicial para vender lo recibido gratuitamente (art.~29).}
}{La capacidad se compone de capacidad de derecho (reconocida a todos) y de ejercicio (limitada para proteger a la persona), con representación, asistencia y Ministerio Público como mecanismos de protección. En los menores rige la autonomía progresiva: la incapacidad es una medida de protección que cede con la edad y la madurez, y termina con la mayoría de edad o la emancipación.}

% ------------------------------------------------------------
\chapter{Las restricciones a la capacidad de las personas con discapacidad}
\label{cap:restricciones}

Este capítulo reconstruye el sistema de restricciones a la capacidad del CCyC: el contexto convencional que ofrece la Convención sobre los Derechos de las Personas con Discapacidad (CDPD), los tres tipos de restricción (capacidad restringida, incapacidad e inhabilitación por prodigalidad), el procedimiento judicial que las rodea y el sistema de apoyos que reemplaza, como criterio general, a la antigua representación absoluta. Quien ejerza la abogacía en cualquier área que involucre contratos, sucesiones, salud o familia se encontrará con una persona cuya capacidad está restringida, o con la necesidad de evaluar si conviene promover esa restricción. Comprender el sistema permite evitar dos errores opuestos: privar de autonomía a quien todavía puede decidir por sí mismo, o desproteger a quien necesita un apoyo real.

Dentro de la capacidad de ejercicio la regla fue siempre la capacidad, con algunos casos de incapaces: las personas por nacer, las menores de edad ---con capacidad progresiva--- y el universo de situaciones de discapacidad por razones de salud mental. El tema se ubica en la Sección Tercera, Capítulo Segundo, Título Primero, Libro Primero del CCyC (véanse los capítulos~\ref{cap:capacidad} y~\ref{cap:menores}).

\begin{nota}
El sistema puede organizarse en cuatro bloques: (1) el contexto convencional que exige acompañar a la persona mediante apoyos en lugar de apartarla de la toma de decisiones; (2) los tipos de restricción, según se trate de quien conserva gran parte de su autonomía (capacidad restringida), de quien no puede en absoluto interaccionar con su entorno (incapacidad) o de quien administra sus bienes en perjuicio de su familia (inhabilitación por prodigalidad); (3) el procedimiento judicial y sus garantías; (4) los efectos de actuar contra la sentencia, la recuperación de la capacidad y el funcionamiento del sistema de apoyos.
\end{nota}

\section{El marco convencional y los antecedentes}

\subsection{La Convención sobre los Derechos de las Personas con Discapacidad}

La CDPD es del año 2006; la Argentina la ratificó en 2008 mediante la Ley 26.378. En 2014, por el procedimiento del artículo 75 inciso 22 de la Constitución Nacional, obtuvo jerarquía constitucional, equiparándose a la primera parte de la Constitución (declaraciones, derechos y garantías). % TODO: verificar consistencia entre fuentes originales: la fuente otorga la jerarquía constitucional en 2014 (véase también el capítulo sobre constitucionalización, que menciona la ratificación en 2008).

La Convención parte de que las personas con discapacidad deben gozar de igualdad de condiciones y ser incluidas en la comunidad, lo que supone superar barreras. La discapacidad surge de la interacción entre una deficiencia ---física, psíquica, sensorial o motriz--- y una barrera que impide la plena incorporación de la persona; las rampas obligatorias en las construcciones o en los colectivos son un modo de superar la barrera que supone, por ejemplo, una escalera. Introduce dos nociones fundamentales:

\begin{definicion}[Apoyo (noción convencional)]
La ayuda necesaria para que la persona con discapacidad pueda integrarse por sí misma. Es fundamental en el ejercicio de la capacidad jurídica.
\end{definicion}

\begin{definicion}[Ajuste razonable]
Adaptación del procedimiento que permite la plena integración y participación de la persona. Se encuentra en los artículos 35 y 36 del CCyC, referidos al proceso judicial. Por ejemplo, si la persona es sordomuda, corresponde disponer de un traductor de lenguaje de señas para que entienda lo que sucede y participe en la audiencia.
\end{definicion}

\subsection{El artículo 12 de la CDPD: igual reconocimiento como persona ante la ley}

\begin{table}[H]
\centering
\begin{tabularx}{\textwidth}{B{0.14\textwidth} Y}
\toprule
\textbf{Art.~12} & \textbf{Contenido} \\
\midrule
Inciso 1 & Las personas con discapacidad tienen derecho al reconocimiento de su personalidad jurídica. \\
Inciso 2 & Tienen capacidad jurídica en igualdad de condiciones con los demás en todos los aspectos de la vida. \\
Inciso 3 & Los Estados deben garantizar los apoyos necesarios para el ejercicio de su capacidad jurídica. \\
Inciso 4 & Salvaguardias: las medidas deben respetar los derechos, la voluntad y las preferencias de la persona; estar libres de conflicto de intereses e influencia indebida; ser proporcionales y adaptadas a las circunstancias; aplicarse en el plazo más corto posible; y estar sujetas a exámenes periódicos por una autoridad u órgano judicial competente, independiente e imparcial. \\
\bottomrule
\end{tabularx}
\caption{Contenido del artículo 12 de la CDPD}
\end{table}

\paragraph{Alcance de «capacidad jurídica» (inc.~2).} Hay dos posturas: quienes entienden que se refiere solo a la capacidad de derecho y quienes sostienen que comprende ambas dimensiones, de derecho y de ejercicio. Agustina Palacios y Francisco Bariffi son referentes argentinos en el estudio de la Convención; Palacios sostiene que el concepto comprende ambas, es decir, que corresponde partir de un principio de capacidad. El Comité de los Derechos de las Personas con Discapacidad, creado por la Convención, sostiene la misma interpretación.

\paragraph{El apoyo (inc.~3).} Es un concepto amplio que involucra un conjunto de medidas ---una persona con discapacidad visual puede usar apoyos informáticos o perros guía---; aquí se trata de los apoyos para el ejercicio de la capacidad jurídica, esto es, para la toma de decisiones.

\begin{ejemplo}[Apoyo para la venta de un inmueble]
Una persona con discapacidad mental dueña de un inmueble que quiere venderlo puede recibir del juez un apoyo que la asesore y preste su conformidad, para evitar que alguien se aproveche y le compre el bien muy por debajo de su valor. La decisión de vender sigue siendo de la persona; el apoyo la asiste.
\end{ejemplo}

\paragraph{Salvaguardias (inc.~4).} Deben ser adecuadas y efectivas, y la medida debe ajustarse a la persona concreta. A diferencia de la Convención sobre los Derechos del Niño, no aparece la noción de «interés superior» sino la de respeto de la voluntad y las preferencias.

\begin{ejemplo}[Respeto de la voluntad y las preferencias]
Una persona dueña de un campo prefiere sembrar soja porque le gusta esa tarea. Si el apoyo considera más conveniente destinarlo al ganado por rendir más, de todos modos debe respetar los derechos, la voluntad y las preferencias de la titular.
\end{ejemplo}

\begin{definicion}[Conflicto de intereses e influencia indebida]
El \textbf{conflicto de intereses} es más objetivo: existe, por ejemplo, cuando quien es designado apoyo es a la vez quien alquila un inmueble que administra o sobre el que debe decidir. La \textbf{influencia indebida} es más difícil de determinar, porque supone cierta manipulación de la voluntad, como puede ocurrir cuando quien ejerce una influencia decisiva convive con la persona.
\end{definicion}

\subsection{El modelo de apoyo a la decisión}

Una observación general del Comité de la CDPD (2014) señala que los Estados deben reemplazar los regímenes basados en la sustitución de decisiones ---la representación, en la que otro sucede a la titular--- por regímenes basados en el apoyo, de modo que la persona decida por sí misma. Es lo que instrumenta el nuevo código. Existe un delicado equilibrio entre autonomía y protección: la Convención enfatiza la autonomía, pero prevé protecciones, y corresponde articular ambos objetivos.

\subsection{El régimen derogado: incapaces e inhabilitados}

El código anterior preveía dos figuras:

\begin{itemize}
\item \textbf{Incapaces de hecho} (art.~141), también llamados insanos o interdictos: personas declaradas incapaces por la concurrencia de un factor biológico (enfermedad mental) y uno jurídico (ineptitud para administrar su persona o sus bienes). Se les designaba un curador y eran básicamente incapaces.
\item \textbf{Inhabilitados} (art.~152 bis): partían de un principio de capacidad y se les nombraba, por sentencia, un curador con funciones delimitadas, en tres supuestos: \textit{embriaguez habitual o uso de estupefacientes} que expusieran a actos perjudiciales (por ejemplo, un alcohólico que vende bienes a bajo precio o los regala); \textit{disminución de facultades} que no llegara a la incapacidad del artículo 141 pero en la que el juez estimara posible un daño; y \textit{prodigalidad}, que solo procedía si la persona tenía familia y ya había dilapidado parte de sus bienes (caso típico: la ludopatía).
\end{itemize}

\paragraph{La Ley de Salud Mental 26.657 y el artículo 152 ter.} Dos años después de aprobada la Convención, en 2010, la Ley de Salud Mental incorporó el artículo 152 ter: la declaración de incapacidad o inhabilitación requiere un examen de facultativos con evaluaciones interdisciplinarias; la sentencia no puede extenderse por más de tres años; y debe especificar las funciones y actos que se limitan. Fue un primer intento parcial, no del todo coherente con el resto del sistema, ya que la incapacidad y su regla general seguían existiendo, pero anticipaba la idea central del nuevo código: la regla es la capacidad y corresponde precisar qué no puede hacerse. Aun así, muchos jueces continuaron declarando a la persona «incapaz para todos los actos de la vida civil».

\section{Los tipos de restricción en el CCyC}

\subsection{Reglas generales (art.~31)}

\begin{regla}[Art.~31 CCyC]
La capacidad de ejercicio se presume, incluso cuando la persona se encuentra internada. Las limitaciones a la capacidad son de carácter excepcional y se imponen siempre en beneficio de la persona. La intervención estatal tiene carácter interdisciplinario, tanto en el tratamiento como en el proceso judicial. La persona tiene derecho a recibir información a través de medios y tecnologías adecuadas, y a participar en el proceso judicial con asistencia letrada, que debe proporcionar el Estado si carece de ella. Deben priorizarse las alternativas terapéuticas menos restrictivas de los derechos y libertades.
\end{regla}

La internación en un establecimiento psiquiátrico no implica necesariamente una restricción: puede haber internación sin restricción, restricción sin internación, o ambas. La intervención interdisciplinaria suele reunir psiquiatría (diagnóstico y tratamiento), psicología (dinamismos psíquicos y bienestar integral) y trabajo social (entorno de la persona, pues la intensidad de la restricción puede variar según cuente o no con una red de contención).

\subsection{Capacidad restringida (art.~32, primer párrafo)}

\begin{regla}[Art.~32, primer párrafo, CCyC]
El juez puede restringir la capacidad para determinados actos de una persona mayor de trece años que padece una adicción o una alteración mental permanente o prolongada, de suficiente gravedad, siempre que estime que del ejercicio de su plena capacidad puede resultar un daño a su persona o a sus bienes.
\end{regla}

Es el supuesto base del sistema y requiere sentencia judicial, pues la Convención exige que toda medida que afecte la capacidad de ejercicio sea dictada por una autoridad judicial independiente. Sus requisitos:

\begin{enumerate}
\item \textbf{Ser mayor de trece años.} Entre los 13 y los 18 la regla es la incapacidad de ejercicio con capacidad progresiva, y la sentencia restringirá lo vinculado con las atribuciones propias del adolescente. Antes de los 13, el artículo 261 presume la falta de discernimiento para los actos lícitos, por lo que no hace falta restricción judicial.
\item \textbf{Requisito intrínseco o biológico:} una adicción o una alteración mental permanente o prolongada, de suficiente gravedad. El código no precisa «adicción»; la doctrina la interpreta a la luz del artículo 4 de la Ley de Salud Mental (consumo problemático de drogas, legales o ilegales); una afición intrascendente carece de la gravedad necesaria. La alteración mental refiere a una afectación del funcionamiento o de las capacidades mentales a los fines de un proceso civil y no debe confundirse con el certificado de discapacidad.
\item \textbf{Requisito extrínseco:} que del ejercicio de la plena capacidad pueda resultar un daño a la persona o a sus bienes. No contempla el daño a terceros, que corresponde a la responsabilidad civil y, si el riesgo fuera cierto e inminente, a la internación.
\end{enumerate}

\paragraph{Régimen.} Las personas con capacidad restringida son, como regla, capaces: la sentencia limita solo los actos y funciones indicados, de modo proporcional y ajustado a su situación. Se designa un \textbf{apoyo} ---no un curador--- cuya modalidad fija el juez: concurrencia con la persona, o apoyo de consejo. El artículo 101 inciso b admite excepcionalmente el apoyo con funciones de representación, que debe estar especialmente fundado para no convertirse en una curatela con otro nombre.

\begin{ejemplo}[Apoyo de asistencia en la venta de un inmueble]
Si a una persona con una adicción se le designa un apoyo de asistencia para vender su inmueble, en la escritura concurren el escribano, la persona con capacidad restringida y el apoyo, y ambos firman en el ámbito de su respectiva intervención.
\end{ejemplo}

\subsection{Incapacidad (art.~32, último párrafo)}

Es un supuesto excepcional (véase el capítulo~\ref{cap:capacidad}). Según las fuentes, un análisis de las sentencias de la Cámara Nacional de Apelaciones en lo Civil en los dos primeros años de vigencia del código mostró que los casos de capacidad restringida eran ampliamente mayoritarios y los de incapacidad, poco frecuentes. Se trata de casos muy severos: la persona no puede entablar relación con su entorno ni expresar su voluntad por ningún medio, y el sistema de apoyos resulta ineficaz. Sin representante quedaría como un «muerto civil», sin que se protejan sus derechos ---por ejemplo, para participar en la asamblea de una sociedad anónima o administrar bienes si estuviera en coma---. La sentencia también debe establecer qué actos y funciones se limitan, con una limitación mayor que en la capacidad restringida, y se designa un \textbf{curador} con funciones de representación; en el nuevo código la curatela queda restringida a este supuesto.

\subsection{Inhabilitación por prodigalidad (art.~48)}

\begin{regla}[Art.~48 CCyC]
Pueden ser inhabilitados quienes por la prodigalidad en la gestión de sus bienes expongan a su cónyuge, conviviente o hijos menores de edad o con discapacidad a la pérdida del patrimonio. A estos inhabilitados se les designarán apoyos. Solo podrán otorgar los actos de disposición entre vivos previstos en la sentencia con la conformidad del apoyo. De ser necesario, el juez puede ampliar la nómina de actos que requieren la conformidad del apoyo.
\end{regla}

La denominación proviene del código anterior, pero ahora queda restringida a la prodigalidad; los otros dos supuestos del artículo 152 bis (embriaguez o toxicomanía y disminución de facultades) pasan a la capacidad restringida. Pródigo es quien derrocha sus bienes ---no debe confundirse con «prodigio»---; la figura protege también a su familia. A diferencia del régimen anterior, no exige que ya se haya dilapidado parte del patrimonio: basta con que se ponga en riesgo.

Las personas inhabilitadas son, como regla, capaces, pero tienen limitada la disposición de bienes ---actos que impliquen enajenarlos--- y los demás actos que fije la sentencia, con apoyos designados.

\begin{ejemplo}[El pródigo y el patrimonio familiar]
Ante quien dilapida bienes jugando en el casino, la familia puede iniciar la acción; la sentencia puede disponer, por ejemplo, que para vender su automóvil o sus bienes registrables, o para hipotecar, necesite el consentimiento de su cónyuge. Si no tuviera familia, no podría ser inhabilitada en estos términos.
\end{ejemplo}

\section{El proceso de determinación de la capacidad}

El CCyC incorpora normas procesales aunque los códigos de procedimiento corresponden a las provincias y a la Nación, para dar una regulación común a una materia que afecta directamente la capacidad.

\begin{longtable}{@{}B{0.20\textwidth}L{0.70\textwidth}@{}}
\toprule
\textbf{Etapa / artículo} & \textbf{Contenido} \\
\midrule
\endfirsthead
\toprule
\textbf{Etapa / artículo} & \textbf{Contenido} \\
\midrule
\endhead
\bottomrule
\endfoot
Juez competente (art.~36 según las fuentes) & El del domicilio o el del lugar de internación; por inmediatez, suele resultar competente el más cercano a la internación. \\
Legitimados (art.~33) & El propio interesado; el cónyuge no separado de hecho y el conviviente mientras no haya cesado la convivencia; los parientes dentro del cuarto grado y, por afinidad, dentro del segundo; el Ministerio Público. \\
Medidas cautelares (art.~34) & Durante el proceso el juez puede ordenar restricciones al ejercicio de la capacidad para garantizar los derechos personales y patrimoniales, determinando con precisión los actos y designando apoyos o, excepcionalmente, curador. \\
Entrevista personal (art.~35) & El juez debe garantizar la inmediatez y entrevistar personalmente al interesado antes de resolver, asegurando accesibilidad y ajustes razonables; deben estar presentes en las audiencias el Ministerio Público y al menos un letrado que asista al interesado. \\
Persona como parte (art.~36) & La persona es parte y puede aportar todas las pruebas; si no tiene abogado propio, el juez debe nombrarle uno; en caso de incapaces es necesaria la presencia del Ministerio Público. \\
Dictamen interdisciplinario (art.~37) & Necesario para declarar incapacidad o capacidad restringida; debe contener diagnóstico y pronóstico, época en que se manifestó la situación, recursos personales, familiares y sociales, y régimen aconsejable para la protección y asistencia con la mayor autonomía posible. \\
Sentencia (art.~38) & Determina la extensión y alcance de la restricción, especifica funciones y actos limitados con la menor afectación posible de la autonomía, designa apoyos o curadores y señala las condiciones de validez de los actos sujetos a restricción, indicando si se actúa con asistencia o representación. \\
Registro (art.~39) & Se inscribe en el Registro de Estado Civil, al margen del acta de nacimiento; produce efectos frente a terceros y se cancela al levantarse la restricción. \\
Revisión (art.~40) & En plazo no superior a tres años, con nuevos dictámenes y audiencia personal; en cualquier momento a instancia del interesado; el Ministerio Público fiscaliza e insta la revisión. \\
Internación (art.~41) & La internación sin consentimiento procede solo con los recaudos de la legislación especial y cuando, a criterio del equipo de salud, medie riesgo cierto e inminente para la persona o terceros. \\
\end{longtable}
% TODO: verificar consistencia entre fuentes originales: la fuente atribuye la competencia judicial al artículo 36, y el mismo artículo 36 a la condición de parte; también menciona «traslados (artículo 42)» sin desarrollarlos.

\paragraph{Sobre los legitimados.} Que el propio interesado pueda iniciar el proceso es una novedad: tiene sentido, por ejemplo, cuando una persona diagnosticada con deterioro cognitivo desea anticiparse y proponer a alguien de su confianza como apoyo. El término «conviviente» es ambiguo y la doctrina discute si abarca a toda persona que convive o solo a quienes tienen unión convivencial registrada; la interpretación parecería amplia. Los parientes legitimados son los que tendrían vocación hereditaria. El código anterior admitía como legitimado a cualquier vecino cuando la persona era considerada «furiosa»; esa posibilidad fue eliminada, y quien no encuadre debe acudir al Ministerio Público.

\paragraph{Sobre las cautelares.} Durante el proceso la persona es, en principio, capaz; las restricciones cautelares son excepcionales y justificadas, por ejemplo, la designación de un apoyo para decisiones sobre un inmueble o un fondo de inversión mientras la persona atraviesa un problema de salud mental.

\paragraph{Sobre la entrevista y la persona como parte.} Aunque la entrevista ya era habitual en algunos fueros, el código la exige con carácter general. El artículo 36 evita que la persona quede ausente de su propio proceso; el abogado que se le provee ejerce su representación judicial, distinta de la del apoyo o del curador.

\paragraph{Sobre el dictamen.} La Convención no exige literalmente un análisis interdisciplinario, pero de su espíritu se desprende que no debe limitarse a lo médico. El diagnóstico y pronóstico distinguen, por ejemplo, una crisis aguda de una situación permanente o un deterioro severo; la época de manifestación es relevante para anular actos anteriores (art.~45); los recursos permiten conocer la red de apoyo natural.

\paragraph{Sobre el registro.} Plantea una tensión entre la protección de la intimidad ---dato altamente sensible--- y la publicidad necesaria para que quienes contratan sepan si su interlocutor tiene una restricción.

\paragraph{Sobre la internación.} No es posible internar a una persona contra su voluntad si no hay riesgo cierto e inminente de un daño de entidad para ella o para terceros; el estándar proviene de la Ley de Salud Mental y busca evitar abusos del pasado, como internar a personas por carecer de vivienda o resultar molestas. La internación involuntaria debe ser breve, acotada, orientada a la recuperación y reintegración, con supervisión periódica, debido proceso, control judicial y defensa ---entre otros mecanismos, la Unidad de Letrados del artículo 22 de la Ley de Salud Mental---.

\section{Validez de los actos y cese de las restricciones}

\subsection{Actos posteriores a la inscripción (art.~44)}

\begin{regla}[Art.~44 CCyC]
Son nulos los actos de la persona incapaz o con capacidad restringida que contrarían lo dispuesto en la sentencia realizada con posterioridad a su inscripción en el Registro de Estado Civil y Capacidad de las Personas.
\end{regla}

La nulidad es relativa, deja sin efecto el acto y retrotrae la situación al estado anterior (si se entregó un inmueble o dinero, debe restituirse). Requiere que el acto sea posterior a la inscripción y contrario a la sentencia: si esta solo restringe la venta de inmuebles, el consentimiento prestado para una intervención quirúrgica es válido.

\begin{ejemplo}[Venta de un inmueble contraria a la sentencia]
Una persona con discapacidad mental que vive de la renta de varias cocheras tiene una sentencia que exige el apoyo de un hermano o hermana para vender inmuebles. Si vende uno sin ese apoyo y el escribano no advierte la restricción, el acto es nulo y las partes deben restituirse lo entregado.
\end{ejemplo}

\subsection{Actos anteriores a la sentencia (art.~45)}

\begin{regla}[Art.~45 CCyC]
Los actos anteriores a la inscripción de la sentencia pueden ser declarados nulos si perjudican a la persona incapaz o con capacidad restringida y se cumple alguno de estos requisitos: a) la enfermedad mental era ostensible al momento de la celebración del acto; b) quien contrató con ella era de mala fe; c) el acto fue a título gratuito.
\end{regla}

Son en principio válidos, pero anulables si concurre alguno de los requisitos; el dictamen del artículo 37 permite determinar si existían causales al celebrarse el acto.

\begin{ejemplo}[Contratante de mala fe]
Un profesional de la salud que atiende a una persona con capacidad restringida y, conociendo su situación, le compra un inmueble a bajo precio, actúa de mala fe.
\end{ejemplo}

\subsection{Cese de las restricciones (art.~47)}

\begin{regla}[Art.~47 CCyC]
El cese de la incapacidad o de la restricción a la capacidad debe decretarlo el juez que la declaró, previo examen de un equipo interdisciplinario integrado conforme al artículo 37, que dictamine sobre el restablecimiento de la persona. Si el restablecimiento es total, el juez deja sin efecto la sentencia. Si es parcial, o el interesado solicita la revisión, el juez puede ampliar la nómina de actos que la persona puede realizar por sí o con la asistencia o representación de los apoyos designados.
\end{regla}

Además de la revisión obligatoria del artículo 40, el interesado puede pedir el cese en cualquier momento por el mismo procedimiento.

\section{El sistema de apoyos (art.~43)}

\begin{regla}[Art.~43 CCyC: concepto de apoyo]
Se entiende por apoyo cualquier medida de carácter judicial o extrajudicial que facilite a la persona la toma de decisiones para dirigir su persona, administrar sus bienes y celebrar actos jurídicos en general. Las medidas de apoyo tienen como función promover la autonomía y facilitar la comunicación, la comprensión y la manifestación de voluntad de la persona para el ejercicio de sus derechos. El interesado puede proponer al juez la designación de una o más personas de su confianza para que le presten apoyo. El juez debe evaluar los alcances de la designación y procurar la protección de la persona respecto de eventuales conflictos de intereses o influencia indebida. La resolución debe establecer la condición y la calidad de las medidas de apoyo y, de ser necesario, ser inscripta en el Registro de Estado Civil y Capacidad de las Personas.
\end{regla}

El apoyo judicial es el designado en la sentencia; el extrajudicial es más difícil de precisar, porque toda persona recibe apoyos informales sin control judicial. La cuestión se presenta cuando el apoyo adquiere relevancia jurídica.

\begin{ejemplo}[Apoyo de hecho sin relevancia jurídica]
Durante un aislamiento, un vecino que hace las compras de una persona mayor que no puede salir presta un apoyo de hecho que no requiere control judicial, siempre que no involucre actos como vender un inmueble, comprar un automóvil o hipotecar una propiedad.
\end{ejemplo}

Para designar apoyos el juez debe valorar el conflicto de intereses y la influencia indebida (véase el apartado de salvaguardias): por ejemplo, la influencia indebida puede plantearse cuando una persona mayor con discapacidad mental convive con uno de sus hijos y los demás consideran que existe manipulación. El criterio de actuación de los apoyos es el respeto de la voluntad y las preferencias de la persona (art.~12 CDPD), que puede proponer a su apoyo. Que una persona jurídica actúe como apoyo se discute en la doctrina, y las fuentes no registran jurisprudencia que lo haya admitido respecto de la institución donde la persona está internada.

\subsection{Conclusiones}

El código se adapta a la Convención con su terminología ---apoyos, ajustes razonables, accesibilidad, autonomía, recursos personales y familiares---; ratifica el criterio de la Ley de Salud Mental según el cual la regla es la capacidad y las limitaciones, excepcionales y por actos y funciones concretos; el sistema general es el de apoyos; se reconoce el derecho de la persona a participar y ser escuchada; la sentencia se revisa periódicamente; y se consolida la intervención interdisciplinaria. Subsiste la discusión sobre si la incapacidad y la sustitución de la voluntad respetan plenamente la Convención; según las fuentes, la doctrina mayoritaria entiende que no la vulneran, por tratarse de situaciones absolutamente excepcionales en las que, sin representante, la persona quedaría sin protección jurídica, y porque conserva el derecho a participar y puede dejar directivas anticipadas, recibidas de modo limitado en el artículo 139, que permite señalar anticipadamente quién podría ser su curador ante la eventual pérdida de conciencia.

\begin{longtable}{@{}B{0.18\textwidth}L{0.35\textwidth}L{0.35\textwidth}@{}}
\toprule
\textbf{Situación} & \textbf{Código Civil derogado} & \textbf{CCyC vigente} \\
\midrule
\endfirsthead
\toprule
\textbf{Situación} & \textbf{Código Civil derogado} & \textbf{CCyC vigente} \\
\midrule
\endhead
\bottomrule
\endfoot
Regla general & Incapacidad como eje; el inhabilitado era la excepción acotada a tres casos. & La capacidad es la regla; toda restricción es excepcional y en beneficio de la persona (art.~31). \\
Figura de protección & Curador con función de representación, para incapaces e inhabilitados. & Apoyos para capacidad restringida e inhabilitados; curador solo para incapaces (art.~32, último párrafo). \\
Intervención profesional & Evaluación médica (luego interdisciplinaria, desde el art.~152 ter de 2010). & Dictamen interdisciplinario obligatorio: psiquiatría, psicología y trabajo social (art.~37). \\
Duración de la sentencia & Sin plazo, hasta el art.~152 ter (2010), que fijó un máximo de tres años. & Revisión obligatoria cada tres años como máximo, o antes a pedido del interesado (art.~40). \\
Voz de la persona & Sin entrevista personal obligatoria ni condición de parte. & Entrevista personal obligatoria (art.~35); la persona es parte (art.~36). \\
Sistema de apoyos & No existía; solo representación a través del curador. & Sistema de apoyos como regla; respeto de la voluntad y preferencias (art.~43). \\
\caption{Régimen anterior y régimen vigente de las restricciones a la capacidad}\\
\end{longtable}

\cornell{Las restricciones a la capacidad}{
\cq{¿Qué exige el art.~12 de la CDPD?}{Reconocimiento de la personalidad jurídica, capacidad jurídica en igualdad de condiciones, apoyos para ejercerla y salvaguardias que respeten la voluntad y preferencias de la persona, sin conflicto de intereses ni influencia indebida, proporcionales, de corta duración y revisables.}
\cq{¿Qué cambió respecto del código derogado?}{Se pasa del eje de la incapacidad y la curatela (arts.~141 y 152 bis) a la capacidad como regla y el sistema de apoyos; el art.~152 ter (Ley 26.657, 2010) fue un primer intento parcial.}
\cq{¿Cuáles son las reglas generales del art.~31?}{La capacidad se presume, incluso en la internación; las limitaciones son excepcionales y en beneficio de la persona; intervención interdisciplinaria; información, asistencia letrada y alternativas terapéuticas menos restrictivas.}
\cq{¿Cuáles son los tres tipos de restricción y su figura de protección?}{Capacidad restringida (art.~32, 1.er párrafo: mayor de 13 años, adicción o alteración mental grave, riesgo de daño a sí o a sus bienes; apoyo); incapacidad (art.~32, último párrafo: excepcional; curador); inhabilitación por prodigalidad (art.~48; apoyo).}
\cq{¿Cómo es el proceso?}{Legitimados (art.~33), cautelares (34), entrevista personal (35), persona como parte con abogado (36), dictamen interdisciplinario (37), sentencia específica (38), registro (39) y revisión cada tres años (40).}
\cq{¿Qué valor tienen los actos celebrados contra la sentencia o antes de ella?}{Los posteriores a la inscripción que contrarían la sentencia son nulos (art.~44); los anteriores son anulables si la enfermedad era ostensible, el contratante obró de mala fe o el acto fue gratuito (art.~45).}
\cq{¿Qué es un apoyo y cómo actúa?}{Cualquier medida judicial o extrajudicial que facilite la toma de decisiones (art.~43); asiste sin sustituir y respeta la voluntad y preferencias de la persona, con salvaguardas frente al conflicto de intereses y la influencia indebida.}
}{El CCyC adapta el régimen a la CDPD: la capacidad es la regla, las restricciones son excepcionales, se fijan por sentencia con dictamen interdisciplinario y se revisan periódicamente, y el sistema general es el de apoyos que asisten sin sustituir la voluntad. El curador y la representación quedan para los casos extremos de incapacidad.}

% ============================================================
\appendix
\chapter{Normativa citada en las fuentes}
\label{app:normativa}

Las fuentes de estos apuntes no incluyen bibliografía doctrinal en forma de lista; se conservan únicamente las referencias normativas y de autores efectivamente mencionadas en el texto. Se agrupan aquí las principales normas citadas.

\begin{itemize}
\item Constitución Nacional: arts.~19, 75 inc.~22 y 75 inc.~23.
\item Código Civil y Comercial de la Nación (Ley 26.994): Título Preliminar (arts.~1 a 18) y, entre otros, los arts.~19 a 21, 22 a 32, 33 a 48, 51, 57, 62 a 72, 73 a 78, 98, 101, 103, 332, 344, 386, 423, 574, 584, 591, 592, 597, 639, 641, 644, 684, 958, 960, 1002, 1004, 1014, 1742, 2279 y 2600.
\item Código Civil de Vélez Sarsfield (Ley 340) y sus reformas: arts.~3, 4, 5, 30, 51, 63, 70, 141, 152 bis, 152 ter, 954, 4044 y 4045.
\item Leyes: 17.711, 18.248, 23.264, 23.515, 23.849, 26.061, 26.378, 26.390, 26.413, 26.579, 26.657, 26.743, 27.077 y 27.364.
\item Instrumentos internacionales: Declaración Universal de Derechos Humanos (1948); Convención Americana sobre Derechos Humanos (Pacto de San José de Costa Rica); Pacto Internacional de Derechos Civiles y Políticos; Convención sobre los Derechos del Niño; Convención sobre los Derechos de las Personas con Discapacidad.
\item Jurisprudencia citada: \textit{Artavia Murillo y otros c.\ Costa Rica} (Corte Interamericana de Derechos Humanos, 2012).
\end{itemize}

\end{document}
